\input{preamble}

% Début du document.
%
\begin{document}

\title{Théories cohomologiques de Weil}


\maketitle

\phantomsection
\label{section-phantom}

\tableofcontents



\section{Introduction}
\label{section-introduction}

\noindent
Nous étudions dans ce chapitre les théories cohomologiques de Weil pour les schémas
projectifs lisses sur un corps de base. En bref, nous entendons par là une théorie
cohomologique $H^*$ munie d'une formule de K\"unneth, d'une dualité de Poincar\'e
et de classes de cycles (satisfaisant aux compatibilités convenables). Nous
avertissons le lecteur que la notion de « théorie cohomologique de Weil »
ne fait pas l'objet d'un accord universel dans la littérature.

\medskip\noindent
Avant d'aborder ce chapitre, le lecteur consultera
Catégories, section \ref{categories-section-monoidal}, et
Homologie, section \ref{homology-section-monoidal}, où
nous définissons les catégories monoïdales (symétriques) et développons
les rudiments de leur langage nécessaires au présent chapitre.
À l'aide de ce langage, nous construisons, à la
section \ref{section-correspondences}, la catégorie graduée monoïdale
symétrique dont les objets sont les schémas projectifs lisses et
les morphismes les correspondances. À la section \ref{section-chow-motives},
nous adjoignons les images des projecteurs et inversons le motif de Lefschetz
pour obtenir la catégorie karoubienne monoïdale symétrique $M_k$
des motifs de Chow. Cette catégorie est munie d'un foncteur contravariant
$$
h : \{\text{schémas projectifs lisses sur }k\} \longrightarrow M_k
$$
Comme nous le verrons ci-dessous, une propriété essentielle d'une théorie
cohomologique de Weil est de se factoriser par $h$.

\medskip\noindent
Tout d'abord, lorsque le corps de base est algébriquement clos, nous définissons
ce que nous appelons une « théorie cohomologique de Weil classique » ;
voir la section \ref{section-axioms-classical}. Cette notion est
celle introduite dans \cite[Section 1.2]{Kleiman-cycles} et
coïncide avec celle de \cite[page 65]{Kleiman-motives}.
En revanche, elle ne coïncide pas a priori avec celle introduite dans
\cite[page 10]{Kleiman-standard}, car l'auteur y ajoute deux axiomes
de type Lefschetz, et l'on ignore si toute théorie cohomologique de Weil
classique au sens du présent chapitre satisfait à ces axiomes.
À la fin de la section \ref{section-axioms-classical}, nous montrons
qu'une théorie cohomologique de Weil classique est de la forme $H^* = G \circ h$,
où $G$ est un foncteur monoïdal symétrique de $M_k$ dans la catégorie
des espaces vectoriels gradués sur le corps de coefficients de $H^*$.

\medskip\noindent
À la section \ref{section-cycles-nonclosed}, nous démontrons quelques lemmes
sur les groupes de cycles sur des corps non clos, qui serviront à étudier
les théories cohomologiques de Weil sur les schémas projectifs lisses
sur un corps quelconque.

\medskip\noindent
Les propriétés suivantes motivent notre choix d'axiomes pour une théorie
cohomologique de Weil $H^*$ sur un corps de base quelconque $k$ :
\begin{enumerate}
\item $H^* = G \circ h$ pour un foncteur monoïdal symétrique $G$ de $M_k$
dans la catégorie des espaces vectoriels gradués sur le corps de coefficients $F$ de $H^*$ ;
\item $G$ doit envoyer le motif de Tate (inverse du motif de Lefschetz)
sur un espace vectoriel de dimension $1$, noté $F(1)$ et placé en degré $-2$ ;
\item lorsque $k$ est algébriquement clos, on doit retrouver la notion
étudiée à la section \ref{section-axioms-classical},
au choix près d'un élément de base de $F(1)$.
\end{enumerate}
Nous analysons d'abord les deux premières conditions à la section \ref{section-axioms}.
Après avoir établi quelques résultats supplémentaires à la section \ref{section-further},
nous ajoutons, à la section \ref{section-old}, les axiomes nécessaires
pour obtenir la propriété (3).

\medskip\noindent
Dans la dernière section \ref{section-c1}, nous exposons en détail
une autre approche des théories cohomologiques de Weil, fondée sur
une application première classe de Chern au lieu des classes de cycles.
C'est cette approche qui se prêtera le mieux à la démonstration, dans
les chapitres ultérieurs, du fait que certaines théories cohomologiques
sont des théories cohomologiques de Weil ; voir
Cohomologie de de Rham, section \ref{derham-section-weil}.





\section{Conventions et notations}
\label{section-conventions}

\noindent
Soit $F$ un corps. Dans ce chapitre, nous considérons la catégorie des espaces
vectoriels gradués sur $F$ comme une catégorie monoïdale symétrique $F$-linéaire,
munie de la contrainte d'associativité usuelle et d'une contrainte de commutativité
faisant intervenir des signes.
Voir Homologie, exemple \ref{homology-example-graded-vector-spaces}.

\medskip\noindent
Soit $R$ un anneau. Dans ce chapitre, une
{\it $R$-algèbre graduée commutative} $A$ est une
$R$-algèbre différentielle graduée commutative
(Algèbre différentielle graduée, définitions \ref{dga-definition-dga} et
\ref{dga-definition-cdga}) dont la différentielle est nulle. Ainsi, $A$
est un $R$-module muni d'une graduation
$A = \bigoplus_{n \in \mathbf{Z}} A^n$ par des
sous-$R$-modules. La multiplication $R$-bilinéaire
$$
A^n \times A^m \longrightarrow A^{n + m},\quad
\alpha \times \beta \longmapsto \alpha \cup \beta
$$
sera appelée le {\it cup-produit} dans ce chapitre.
La contrainte de commutativité s'écrit
$\alpha \cup \beta = (-1)^{nm} \beta \cup \alpha$ si
$\alpha \in A^n$ et $\beta \in A^m$. Enfin, il existe
un élément unité $1 \in A^0$ pour la multiplication, ou, ce qui revient au même,
une application additive et multiplicative $R \to A^0$, compatible avec
la structure de $R$-module de $A$.

\medskip\noindent
Soit $k$ un corps et soit $X$ un schéma de type fini sur $k$.
Les groupes de Chow $\CH_k(X)$ de $X$, formés des cycles de dimension $k$
modulo l'équivalence rationnelle, ont été définis dans
Homologie de Chow, définition \ref{chow-definition-rational-equivalence}.
Si $X$ est normal ou de Cohen-Macaulay, on peut aussi considérer
les groupes de Chow $\CH^p(X)$ de cycles de codimension $p$
(Homologie de Chow, section \ref{chow-section-cycles-codimension}) ;
alors $[X] \in \CH^0(X)$ désigne la « classe fondamentale » de $X$ ; voir
Homologie de Chow, remarque \ref{chow-remark-fundamental-class}.
Si $X$ est lisse et si $\alpha$ et $\beta$ sont des cycles sur $X$,
$\alpha \cdot \beta$ désigne le produit d'intersection de
$\alpha$ et $\beta$ ; voir
Homologie de Chow, section \ref{chow-section-intersection-product}.













\section{Correspondances}
\label{section-correspondences}

\noindent
Soit $k$ un corps. Pour des schémas $X$ et $Y$ sur $k$, nous notons
$X \times Y$ le produit de $X$ et $Y$ dans la catégorie des schémas
sur $k$. Nous construisons dans cette section la catégorie graduée sur
$\mathbf{Q}$ dont les objets sont les schémas projectifs lisses sur $k$ et dont
les morphismes sont les correspondances.

\medskip\noindent
Soient $X$ et $Y$ des schémas projectifs lisses sur $k$.
Soit $X = \coprod X_d$ la décomposition de $X$ en
sous-schémas ouverts et fermés équidimensionnels
tels que $\dim(X_d) = d$. Nous définissons l'espace vectoriel sur $\mathbf{Q}$
{\it des correspondances de degré $r$ de $X$ dans $Y$}
par la formule :
$$
\text{Corr}^r(X, Y) =
\bigoplus\nolimits_d \CH^{d + r}(X_d \times Y) \otimes \mathbf{Q}
\subset
\CH^*(X \times Y) \otimes \mathbf{Q}
$$
Étant donnés $c \in \text{Corr}^r(X, Y)$ et $\beta \in \CH_j(Y) \otimes \mathbf{Q}$,
nous définissons l'{\it image réciproque} de $\beta$ par $c$ par la formule
$$
c^*(\beta) = \text{pr}_{1, *}(c \cdot \text{pr}_2^*\beta)
\quad\text{dans}\quad
\CH_{j - r}(X) \otimes \mathbf{Q}
$$
Cette expression a un sens, car $\text{pr}_2$ est plat de dimension relative
$d$ sur $X_d \times Y$ ; ainsi $\text{pr}_2^*\beta$ est un cycle de
dimension $d + j$ sur $X_d \times Y$, donc $c \cdot \text{pr}_2^*\beta$
est un cycle de dimension $j - r$ sur $X_d \times Y$, dont l'image directe
par le morphisme propre $\text{pr}_1$ est un cycle de même dimension.
De même, en passant à la graduation par la codimension,
pour $\alpha \in \CH^i(X) \otimes \mathbf{Q}$, nous définissons
l'{\it image directe} de $\alpha$ par $c$ par la formule
$$
c_*(\alpha) = \text{pr}_{2, *}(c \cdot \text{pr}_1^*\alpha)
\quad\text{dans}\quad
\CH^{i + r}(Y) \otimes \mathbf{Q}
$$
Cette expression a un sens, car $\text{pr}_1^*\alpha$ est un cycle de codimension
$i$ sur $X \times Y$ ; ainsi $c \cdot \text{pr}_1^*\alpha$ est un cycle
de codimension $i + d + r$ sur $X_d \times Y$, dont l'image directe
est un cycle de codimension $i + r$ sur $Y$.

\medskip\noindent
Étant donnés trois schémas projectifs lisses $X, Y, Z$ sur $k$, nous définissons une
{\it composition des correspondances}
$$
\text{Corr}^s(Y, Z)
\times
\text{Corr}^r(X, Y)
\longrightarrow
\text{Corr}^{r + s}(X, Z)
$$
par la règle
$$
(c', c)
\longmapsto
c' \circ c =
\text{pr}_{13, *}(\text{pr}_{12}^*c \cdot \text{pr}_{23}^*c')
$$
où $\text{pr}_{12} : X \times Y \times Z \to X \times Y$
est la projection, et de même pour $\text{pr}_{13}$ et $\text{pr}_{23}$.

\begin{lemma}
\label{lemma-composition-correspondences}
Les correspondances possèdent les propriétés suivantes :
\begin{enumerate}
\item la composition des correspondances est $\mathbf{Q}$-bilinéaire
et associative ;
\item il existe un isomorphisme canonique
$$
\CH_{-r}(X) \otimes \mathbf{Q} = \text{Corr}^r(X, \Spec(k))
$$
par lequel l'image réciproque par les correspondances correspond à la composition ;
\item il existe un isomorphisme canonique
$$
\CH^r(X) \otimes \mathbf{Q} = \text{Corr}^r(\Spec(k), X)
$$
par lequel l'image directe par les correspondances correspond à la composition ;
\item la composition des correspondances est compatible avec l'image directe
et l'image réciproque des cycles.
\end{enumerate}
\end{lemma}

\begin{proof}
La bilinéarité résulte immédiatement de la linéarité de l'image directe
et de l'image réciproque, et de la bilinéarité du produit d'intersection.
Pour démontrer l'associativité, donnons-nous
$X, Y, Z, W$ ainsi que $c \in \text{Corr}(X, Y)$, $c' \in \text{Corr}(Y, Z)$ et
$c'' \in \text{Corr}(Z, W)$. On a alors
\begin{align*}
c'' \circ (c' \circ c)
& =
\text{pr}^{134}_{14, *}(
\text{pr}^{134, *}_{13}
\text{pr}^{123}_{13, *}(\text{pr}^{123, *}_{12}c \cdot
\text{pr}^{123, *}_{23}c')
\cdot \text{pr}^{134, *}_{34}c'') \\
& =
\text{pr}^{134}_{14, *}(
\text{pr}^{1234}_{134, *}
\text{pr}^{1234, *}_{123}(\text{pr}^{123, *}_{12}c \cdot
\text{pr}^{123, *}_{23}c')
\cdot \text{pr}^{134, *}_{34}c'') \\
& =
\text{pr}^{134}_{14, *}(
\text{pr}^{1234}_{134, *}
(\text{pr}^{1234, *}_{12}c \cdot
\text{pr}^{1234, *}_{23}c')
\cdot \text{pr}^{134, *}_{34}c'') \\
& =
\text{pr}^{134}_{14, *}
\text{pr}^{1234}_{134, *}
((\text{pr}^{1234, *}_{12}c \cdot
\text{pr}^{1234, *}_{23}c')
\cdot \text{pr}^{1234, *}_{34}c'') \\
& =
\text{pr}^{1234}_{14, *}(
(\text{pr}^{1234, *}_{12}c \cdot
\text{pr}^{1234, *}_{23}c') \cdot
\text{pr}^{1234, *}_{34}c'')
\end{align*}
Nous notons ici
$$
p^{1234}_{134} : X \times Y \times Z \times W \to X \times Z \times W
\quad\text{et}\quad
p^{134}_{14} : X \times Z \times W \to X \times W
$$
{\emergencystretch=1em les projections, et de même pour les autres indices.
La première égalité est la définition de la composition.
La deuxième résulte de l'égalité
$\text{pr}^{134, *}_{13} \text{pr}^{123}_{13, *} =
\text{pr}^{1234}_{134, *} \text{pr}^{1234, *}_{123}$,
d'après Homologie de Chow, lemme \ref{chow-lemma-flat-pullback-proper-pushforward}.
La troisième résulte de la compatibilité du produit d'intersection
avec l'application de Gysin de $p^{1234}_{123}$ (qui est donnée par l'image réciproque plate) ; voir
Homologie de Chow, lemme \ref{chow-lemma-lci-gysin-product}.
La quatrième égalité résulte de la formule de projection pour
$p^{1234}_{134}$ ; voir Homologie de Chow, lemme \ref{chow-lemma-projection-formula}.
La quatrième égalité exprime la compatibilité de l'image directe propre
avec la composition ; voir Homologie de Chow, lemme \ref{chow-lemma-compose-pushforward}.
Le produit d'intersection étant associatif d'après
Homologie de Chow, lemme \ref{chow-lemma-associative},
cela achève la démonstration de l'associativité de la composition des correspondances.\par}

\medskip\noindent
Nous omettons les démonstrations de (2) et (3), qui consistent essentiellement
à vérifier avec soin les espaces où vivent les différents cycles et leurs (co)dimensions.

\medskip\noindent
L'assertion relative à l'image directe et à l'image réciproque des cycles
signifie que $(c' \circ c)^*(\alpha) = c^*((c')^*(\alpha))$ et
$(c' \circ c)_*(\alpha) = (c')_*(c_*(\alpha))$.
Elle résulte de la combinaison de (1), (2) et (3).
\end{proof}

\begin{example}
\label{example-graph-correspondence}
Soit $f : Y \to X$ un morphisme de schémas projectifs lisses sur $k$.
Notons $\Gamma_f \subset X \times Y$ le graphe de $f$. Plus précisément,
$\Gamma_f$ est l'image de l'immersion fermée
$$
(f, \text{id}_Y) : Y \longrightarrow X \times Y
$$
Soit $X = \coprod X_d$ la décomposition de $X$ en ses
parties ouvertes et fermées $X_d$, équidimensionnelles de dimension $d$.
Alors $\Gamma_f \cap (X_d \times Y)$ est de codimension pure $d$. Par conséquent,
$[\Gamma_f] \in \CH^*(X \times Y) \otimes \mathbf{Q}$
appartient à $\text{Corr}^0(X \times Y)$ ; autrement dit, $[\Gamma_f]$
est une correspondance de degré $0$ de $X$ dans $Y$.
\end{example}

\begin{lemma}
\label{lemma-category-correspondences}
Les schémas projectifs lisses sur $k$, munis des correspondances et
de leur composition définies ci-dessus, forment une catégorie graduée sur
$\mathbf{Q}$
(Algèbre différentielle graduée, définition \ref{dga-definition-graded-category}).
\end{lemma}

\begin{proof}
Tout résulte de la construction et du
lemme \ref{lemma-composition-correspondences},
sauf l'existence des morphismes identiques.
Pour un schéma projectif lisse $X$, considérons
la classe $[\Delta]$ de la diagonale $\Delta \subset X \times X$
dans $\text{Corr}^0(X, X)$. Observons que $\Delta$
est le graphe de l'identité $\text{id}_X : X \to X$ ;
nous utiliserons ce fait ci-dessous.

\medskip\noindent
Pour démontrer que $[\Delta]$ joue le rôle d'identité, il faut établir
$[\Delta] \circ c = c$ et $c' \circ [\Delta] = c'$ pour toutes correspondances
$c \in \text{Corr}^r(Y, X)$ et $c' \in \text{Corr}^s(X, Y)$.
Dans le deuxième cas, il faut montrer que
$$
c' = \text{pr}_{13, *}(\text{pr}_{12}^*[\Delta] \cdot \text{pr}_{23}^*c')
$$,
où $\text{pr}_{12} : X \times X \times Y \to X \times X$ est la
projection, et de même pour $\text{pr}_{13}$ et $\text{pr}_{23}$.
On peut écrire $c' = \sum a_i [Z_i]$, pour certains sous-schémas fermés intègres
$Z_i \subset X \times Y$ et nombres rationnels $a_i$. Il suffit donc
manifestement de montrer que
$$
[Z] = \text{pr}_{13, *}(\text{pr}_{12}^*[\Delta] \cdot \text{pr}_{23}^*[Z])
$$
dans le groupe de Chow de $X \times Y$, pour tout sous-schéma fermé intègre $Z$
de $X \times Y$. En remplaçant $X$ et $Y$ par
la composante irréductible contenant l'image de $Z$ par chacune des deux projections,
on peut supposer $X$ et $Y$ également intègres. Il faut alors montrer
$$
[Z] = \text{pr}_{13, *}([\Delta \times Y] \cdot [X \times Z])
$$
Notons $Z' \subset X \times X \times Y$ l'image de $Z$ par le morphisme
$(\Delta, 1) : X \times Y \to X \times X \times Y$. Alors $Z'$
est un sous-schéma fermé de $X \times X \times Y$ isomorphe à $Z$, et
$Z' = \Delta \times Y \cap X \times Z$ au sens schématique.
D'après Homologie de Chow, lemme \ref{chow-lemma-intersect-properly}\footnote{Le
lecteur vérifiera que $\dim(Z') = \dim(\Delta \times Y) + \dim(X \times Z) -
\dim(X \times X \times Y)$ et que $Z'$ possède un unique point générique,
dont l'image est le point générique de $Z$ (où l'anneau local est de Cohen-Macaulay)
et un point de $X$ (où l'anneau local est de Cohen-Macaulay).
Toutes les hypothèses du lemme sont donc bien vérifiées.},
on en déduit
$$
[Z'] = [\Delta \times Y] \cdot [X \times Z]
$$
Puisque $Z'$ s'envoie aussi isomorphiquement sur $Z$ par $\text{pr}_{13}$,
on obtient la conclusion. La vérification de
$[\Delta] \circ c = c$ est analogue et nous l'omettons.
\end{proof}

\begin{lemma}
\label{lemma-contravariant-functor}
Il existe un foncteur contravariant de la catégorie des schémas
projectifs lisses sur $k$ dans la catégorie des correspondances,
qui est l'identité sur les objets et envoie $f : Y \to X$ sur
l'élément $[\Gamma_f] \in \text{Corr}^0(X, Y)$.
\end{lemma}

\begin{proof}
Dans la démonstration du lemme \ref{lemma-category-correspondences},
nous avons vu que cette construction envoie les identités sur
les identités. Pour achever la démonstration, il faut montrer que, si $g : Z \to Y$
est un autre morphisme de schémas projectifs lisses sur $k$, alors
$[\Gamma_g] \circ [\Gamma_f] = [\Gamma_{f \circ g}]$ dans
$\text{Corr}^0(X, Z)$. En raisonnant comme dans la démonstration du
lemme \ref{lemma-category-correspondences}, on voit qu'il
suffit de montrer
$$
[\Gamma_{f \circ g}] =
\text{pr}_{13, *}([\Gamma_f \times Z] \cdot [X \times \Gamma_g])
$$
dans $\CH^*(X \times Z)$ lorsque $X$, $Y$ et $Z$ sont intègres.
Notons $Z' \subset X \times Y \times Z$ l'image de l'immersion fermée
$(f \circ g, g, 1) : Z \to X \times Y \times Z$.
Alors $Z' = \Gamma_f \times Z \cap X \times \Gamma_g$
au sens schématique, et l'on déduit de
Homologie de Chow, lemme \ref{chow-lemma-intersect-properly},
que
$$
[Z'] = [\Gamma_f \times Z] \cdot [X \times \Gamma_g]
$$
Comme il est clair que $\text{pr}_{13, *}([Z']) = [\Gamma_{f \circ g}]$,
la démonstration est terminée.
\end{proof}

\begin{remark}
\label{remark-transpose}
Soient $X$ et $Y$ des schémas projectifs lisses sur $k$.
Supposons $X$ équidimensionnel de dimension $d$ et
$Y$ équidimensionnel de dimension $e$. L'isomorphisme
$X \times Y \to Y \times X$ qui échange les facteurs détermine alors
un isomorphisme
$$
\text{Corr}^r(X, Y) \longrightarrow \text{Corr}^{d - e + r}(Y, X),\quad
c \longmapsto c^t
$$,
appelé la {\it transposition}. Il agit aussi bien sur les cycles que sur leurs classes.
Un exemple parfois utile est la transposée
$[\Gamma_f]^t = [\Gamma_f^t]$ du graphe d'un morphisme $f : Y \to X$.
\end{remark}

\begin{lemma}
\label{lemma-functor-and-cycles}
Soit $f : Y \to X$ un morphisme de schémas projectifs lisses sur $k$.
Soit $[\Gamma_f] \in \text{Corr}^0(X, Y)$ comme dans
l'exemple \ref{example-graph-correspondence}. Alors :
\begin{enumerate}
\item l'image directe des cycles par la correspondance $[\Gamma_f]$
coïncide avec l'application de Gysin $f^! : \CH^*(X) \to \CH^*(Y)$ ;
\item l'image réciproque des cycles par la correspondance $[\Gamma_f]$
coïncide avec l'application image directe $f_* : \CH_*(Y) \to \CH_*(X)$ ;
\item si $X$ et $Y$ sont équidimensionnels de dimensions $d$ et $e$,
alors :
\begin{enumerate}
\item l'image directe des cycles par la correspondance
$[\Gamma_f^t]$ de la remarque \ref{remark-transpose}
correspond à l'image directe des cycles par $f$ ;
\item l'image réciproque des cycles par la correspondance
$[\Gamma_f^t]$ de la remarque \ref{remark-transpose}
correspond à l'application de Gysin $f^!$.
\end{enumerate}
\end{enumerate}
\end{lemma}

\begin{proof}
Démontrons (1). Rappelons que
$[\Gamma_f]_*(\alpha) =
\text{pr}_{2, *}([\Gamma_f] \cdot \text{pr}_1^*\alpha)$.
On a
$$
[\Gamma_f] \cdot \text{pr}_1^*\alpha =
(f, 1)_*((f, 1)^! \text{pr}_1^*\alpha) =
(f, 1)_*((f, 1)^! \text{pr}_1^!\alpha) =
(f, 1)_*(f^!\alpha)
$$
La première égalité résulte de Homologie de Chow, lemme
\ref{chow-lemma-intersect-regularly-embedded}.
La deuxième résulte de
Homologie de Chow, lemme \ref{chow-lemma-lci-gysin-flat}.
La troisième résulte de $\text{pr}_1 \circ (f, 1) = f$ et de
Homologie de Chow, lemme \ref{chow-lemma-lci-gysin-composition}.
On conclut alors puisque
$\text{pr}_{2, *} \circ (f, 1)_* = 1_*$, d'après
Homologie de Chow, lemme \ref{chow-lemma-compose-pushforward}.

\medskip\noindent
Démontrons (2). Rappelons que $[\Gamma_f]_*(\beta) =
\text{pr}_{1, *}([\Gamma_f] \cdot \text{pr}_2^*\beta)$.
En raisonnant exactement comme ci-dessus, on a
$$
[\Gamma_f] \cdot \text{pr}_2^*\beta = (f, 1)_*\beta
$$
Le résultat s'en déduit comme précédemment.

\medskip\noindent
Démontrons (3). La démonstration est exactement analogue à la précédente.
\end{proof}

\begin{example}
\label{example-decompose-P1}
Soit $X = \mathbf{P}^1_k$. On a alors
$$
\text{Corr}^0(X, X) = \CH^1(X \times X) \otimes \mathbf{Q} =
\CH_1(X \times X) \otimes \mathbf{Q}
$$
Choisissons un point $k$-rationnel $x \in X$ et
considérons les cycles $c_0 = [x \times X]$ et $c_2 = [X \times x]$.
Un calcul montre que $1 = [\Delta] = c_0 + c_2$ dans $\text{Corr}^0(X, X)$
et que l'on a les règles de composition suivantes :
$c_0 \circ c_0 = c_0$,
$c_0 \circ c_2 = 0$,
$c_2 \circ c_0 = 0$ et
$c_2 \circ c_2 = c_2$.
Autrement dit, $c_0$ et $c_2$ sont des idempotents orthogonaux dans
l'algèbre $\text{Corr}^0(X, X)$ ; en fait, on obtient
$$
\text{Corr}^0(X, X) = \mathbf{Q} \times \mathbf{Q}
$$
comme $\mathbf{Q}$-algèbre.
\end{example}

\noindent
La catégorie des correspondances est une catégorie monoïdale symétrique.
Pour des schémas projectifs lisses $X$ et $Y$ sur $k$, on pose
$X \otimes Y = X \times Y$. Pour quatre schémas projectifs lisses
$X, X', Y, Y'$ sur $k$, on définit un produit tensoriel
$$
\otimes :
\text{Corr}^r(X, Y) \times \text{Corr}^{r'}(X', Y')
\longrightarrow
\text{Corr}^{r + r'}(X \times X', Y \times Y')
$$
par la règle
$$
(c, c') \longmapsto
c \otimes c' = \text{pr}_{13}^*c \cdot \text{pr}_{24}^*c'
$$,
où $\text{pr}_{13} : X \times X' \times Y \times Y' \to X \times Y$
et $\text{pr}_{24} : X \times X' \times Y \times Y' \to X' \times Y'$
sont les projections. Comme contrainte d'associativité
$$
X \otimes (Y \otimes Z) = (X \otimes Y) \otimes Z
$$,
{\emergencystretch=1em on prend la contrainte d'associativité usuelle des produits de schémas.
La contrainte de commutativité est donnée par l'isomorphisme
$X \times Y \to Y \times X$ qui échange les facteurs.\par}

\begin{lemma}
\label{lemma-tensor-product}
Le produit tensoriel des correspondances défini ci-dessus fait de la catégorie
des correspondances une catégorie monoïdale symétrique d'unité $\Spec(k)$.
\end{lemma}

\begin{proof}
Démonstration omise.
\end{proof}

\begin{lemma}
\label{lemma-prep-dual}
Soit $f : Y \to X$ un morphisme de schémas projectifs lisses sur $k$.
Supposons $X$ et $Y$ équidimensionnels de dimensions $d$ et $e$.
Notons $a = [\Gamma_f] \in \text{Corr}^0(X, Y)$ et
$a^t = [\Gamma_f^t] \in \text{Corr}^{d - e}(Y, X)$.
Posons
$\eta_X = [\Gamma_{X \to X \times X}] \in \text{Corr}^0(X \times X, X)$,
$\eta_Y = [\Gamma_{Y \to Y \times Y}] \in \text{Corr}^0(Y \times Y, Y)$,
$[X] \in \text{Corr}^{-d}(X, \Spec(k))$ et
$[Y] \in \text{Corr}^{-e}(Y, \Spec(k))$. Le diagramme
$$
\xymatrix{
X \otimes Y \ar[r]_{a \otimes \text{id}} \ar[d]_{\text{id} \otimes a^t} &
Y \otimes Y \ar[r]_{\eta_Y} &
Y \ar[d]^{[Y]} \\
X \otimes X \ar[r]^{\eta_X} &
X \ar[r]^{[X]} &
\Spec(k)
}
$$
est commutatif dans la catégorie des correspondances.
\end{lemma}

\begin{proof}
{\emergencystretch=1em Rappelons que $\text{Corr}^r(W, \Spec(k)) = \CH_{-r}(W)$ pour tout
schéma projectif lisse $W$ sur $k$
et que, pour $c \in \text{Corr}^s(W', W)$, la composition
avec $c$ coïncide avec l'image réciproque par $c$, comme application
$\CH_{-r}(W) \to \CH_{-r - s}(W')$
(lemme \ref{lemma-composition-correspondences}).
Enfin, le lemme \ref{lemma-functor-and-cycles}
indique comment exprimer cela au moyen des images directes
et réciproques usuelles des cycles.
On a
$$
(a \otimes \text{id})^* \eta_Y^* [Y] =
(a \otimes \text{id})^* [\Delta_Y] =
(f \times \text{id})_*\Delta_Y = [\Gamma_f]
$$,
et, par l'autre chemin, on obtient
$$
(\text{id} \otimes a^t)^* \eta_X^* [X] =
(\text{id} \otimes a^t)^* [\Delta_X] =
(\text{id} \times f)^![\Delta_X] = [\Gamma_f]
$$
La dernière égalité résulte de
Homologie de Chow, lemme \ref{chow-lemma-lci-gysin-easy}.
Autrement dit, les deux chemins dans le diagramme donnent
l'élément de $\text{Corr}^d(X \times Y, \Spec(k))$
correspondant au cycle $\Gamma_f \subset X \times Y$.\par}
\end{proof}







\section{Motifs de Chow}
\label{section-chow-motives}

\noindent
Fixons un corps de base $k$. Nous construisons dans cette section une catégorie
additive karoubienne $\mathbf{Q}$-linéaire $M_k$, munie
d'une structure monoïdale symétrique et d'un foncteur contravariant
$$
h : \{\text{schémas projectifs lisses sur }k\} \longrightarrow M_k
$$
qui transforme les produits en produits tensoriels et les réunions disjointes en sommes directes.
Notre construction sera caractérisée par le fait que $h$ se factorise par
la catégorie monoïdale symétrique dont les objets sont les variétés projectives lisses
et les morphismes les correspondances de degré $0$, de telle sorte que
l'image du projecteur $c_2$ sur $h(\mathbf{P}^1_k)$ de
l'exemple \ref{example-decompose-P1} soit inversible dans $M_k$ ; voir
le lemme \ref{lemma-characterize-motives}.
À la fin de cette section, nous montrerons que tout motif, c'est-à-dire
tout objet de $M_k$, admet un dual (à gauche) ; voir le lemme \ref{lemma-dual-general}.

\medskip\noindent
Un {\it motif}, ou {\it motif de Chow}, sur $k$ sera un triplet
$(X, p, m)$, où
\begin{enumerate}
\item $X$ est un schéma projectif lisse sur $k$ ;
\item $p \in \text{Corr}^0(X, X)$ satisfait à $p \circ p = p$ ;
\item $m \in \mathbf{Z}$.
\end{enumerate}
Étant donné un second motif $(Y, q, n)$, nous appelons
{\it morphisme de motifs}, ou {\it morphisme de motifs de Chow},
un élément de
$$
\Hom((X, p, m), (Y, q, n)) =
q \circ \text{Corr}^{n - m}(X, Y) \circ p \subset \text{Corr}^{n - m}(X, Y)
$$
La composition des morphismes de motifs est définie au moyen
de la composition des correspondances définie ci-dessus.

\begin{lemma}
\label{lemma-motives}
La catégorie $M_k$ dont les objets sont les motifs sur $k$ et dont les morphismes
sont les morphismes de motifs sur $k$ est une catégorie $\mathbf{Q}$-linéaire.
Il existe un foncteur contravariant
$$
h : \{\text{schémas projectifs lisses sur }k\} \longrightarrow M_k
$$
défini par $h(X) = (X, 1, 0)$ et $h(f) = [\Gamma_f]$.
\end{lemma}

\begin{proof}
Cela résulte immédiatement du lemme \ref{lemma-contravariant-functor}.
\end{proof}

\begin{lemma}
\label{lemma-Karoubian}
La catégorie $M_k$ est karoubienne.
\end{lemma}

\begin{proof}
Soit $M = (X, p, m)$ un motif et soit $a \in \Mor(M, M)$
un projecteur. Alors $a = a \circ a$ aussi bien dans $\Mor(M, M)$
que dans $\text{Corr}^0(X, X)$. Posons $N = (X, a, m)$.
Puisque $a = p \circ a \circ a$ dans $\text{Corr}^0(X, X)$,
on voit que $a : N \to M$ est un morphisme de $M_k$.
Supposons ensuite que $b : (Y, q, n) \to M$ soit un morphisme
tel que $(1 - a) \circ b = 0$. Alors $b = a \circ b$, et l'on a aussi
$b = b \circ q$. Par conséquent, $b$ est un morphisme $b : (Y, q, n) \to N$.
Le projecteur $1 - a$ admet donc un noyau, à savoir $N$,
et $M_k$ est karoubienne ; voir
Homologie, définition \ref{homology-definition-karoubian}.
\end{proof}

\noindent
Nous définissons un foncteur
$$
\otimes : M_k \times M_k \longrightarrow M_k
$$
Sur les objets, nous employons la formule
$$
(X, p, m) \otimes (Y, q, n) = (X \times Y, p \otimes q, m + n)
$$
Sur les morphismes, nous employons l'application
$$
\xymatrix{
\Mor((X, p, m), (Y, q, n)) \times
\Mor((X', p', m'), (Y', q', n')) \ar[d] \\
\Mor(
(X \times X', p \otimes p', m + m'),
(Y \times Y', q \otimes q', n + n'))
}
$$,
donnée par la règle $(a, a') \longmapsto a \otimes a'$, où
$\otimes$ est le produit des correspondances de la section \ref{section-correspondences}.
Cela a un sens : par définition des morphismes de motifs,
on peut écrire $a = q \circ c \circ p$ et $a' = q' \circ c' \circ p'$,
avec $c \in \text{Corr}^{n - m}(X, Y)$ et
$c' \in \text{Corr}^{n' - m'}(X', Y')$ ;
on obtient alors
$$
a \otimes a' =
(q \circ c \circ p) \otimes (q' \circ c' \circ p') =
(q \otimes q') \circ (c \otimes c') \circ (p \otimes p')
$$,
qui est bien un morphisme de motifs de
$(X \times X', p \otimes p', m + m')$ dans
$(Y \times Y', q \otimes q', n + n')$.

\begin{lemma}
\label{lemma-motives-monoidal}
La catégorie $M_k$, munie du produit tensoriel défini ci-dessus,
est monoïdale symétrique pour les contraintes évidentes d'associativité
et de commutativité, et pour l'unité $\mathbf{1} = (\Spec(k), 1, 0)$.
\end{lemma}

\begin{proof}
Cela résulte aisément du lemme \ref{lemma-tensor-product}. Nous omettons les détails.
\end{proof}

\noindent
Les motifs $\mathbf{1}(n) = (\Spec(k), 1, n)$ seront utiles. Observons que
$$
\mathbf{1} = \mathbf{1}(0)
\quad\text{et}\quad
\mathbf{1}(n + m) = \mathbf{1}(n) \otimes \mathbf{1}(m)
$$
Le produit tensoriel par $\mathbf{1}(1)$ est donc une autoéquivalence de
la catégorie des motifs. Pour un motif $M$, nous écrivons parfois
$M(n) = M \otimes \mathbf{1}(n)$. Observons que, si $M = (X, p, m)$,
alors $M(n) = (X, p, m + n)$.

\begin{lemma}
\label{lemma-inverse-h2}
Avec les notations de l'exemple \ref{example-decompose-P1} :
\begin{enumerate}
\item
le motif $(X, c_0, 0)$ est isomorphe au motif
$\mathbf{1} = (\Spec(k), 1, 0)$ ;
\item
le motif $(X, c_2, 0)$ est isomorphe au motif
$\mathbf{1}(-1) = (\Spec(k), 1, -1)$.
\end{enumerate}
\end{lemma}

\begin{proof}
{\emergencystretch=2em Nous utiliserons le lemme \ref{lemma-contravariant-functor} sans le mentionner de nouveau.
Le morphisme structural $X \to \Spec(k)$ donne une correspondance
$a \in \text{Corr}^0(\Spec(k), X)$. D'autre part, le point rationnel
$x$ est un morphisme $\Spec(k) \to X$, qui donne une correspondance
$b \in \text{Corr}^0(X, \Spec(k))$. On a $b \circ a = 1$ comme
correspondance sur $\Spec(k)$. La composée $a \circ b$ correspond
au graphe de la composée $X \to x \to X$, qui est
$c_0 = [x \times X]$. Ainsi, $a = a \circ b \circ a = c_0 \circ a$
et $b = b \circ a \circ b = b \circ c_0$.
En revenant aux définitions, on voit donc que
$a$ et $b$ sont des morphismes inverses l'un de l'autre
$a : (\Spec(k), 1, 0) \to (X, c_0, 0)$ et
$b : (X, c_0, 0) \to (\Spec(k), 1, 0)$.\par}

\medskip\noindent
Nous procédons exactement comme ci-dessus pour démontrer la seconde assertion.
Notons
$$
a' \in \text{Corr}^1(\Spec(k), X) = \CH^1(X)
$$
la classe du point $x$. Notons
$$
b' \in \text{Corr}^{-1}(X, \Spec(k)) = \CH_1(X)
$$
la classe de $[X]$. On a $b' \circ a' = 1$ comme correspondance
sur $\Spec(k)$, puisque $[x] \cdot [X] = [x]$ sur
$X = \Spec(k) \times X \times \Spec(k)$. Le calcul du
produit d'intersection $\text{pr}_{12}^*b' \cdot \text{pr}_{23}^*a'$
sur $X \times \Spec(k) \times X$ donne le cycle
$X \times \Spec(k) \times x$. Par conséquent,
la composée $a' \circ b'$ est égale à $c_2$ comme
correspondance sur $X$. Ainsi, $a' = a' \circ b \circ a' = c_2 \circ a'$
et $b' = b' \circ a' \circ b' = b' \circ c_2$. Rappelons que
$$
\Mor((\Spec(k), 1, -1), (X, c_2, 0)) =
c_2 \circ \text{Corr}^1(\Spec(k), X)
\subset
\text{Corr}^1(\Spec(k), X)
$$
et
$$
\Mor((X, c_2, 0), (\Spec(k), 1, -1)) =
\text{Corr}^{-1}(X, \Spec(k)) \circ c_2
\subset
\text{Corr}^{-1}(X, \Spec(k))
$$
On voit donc que $a'$ et $b'$ sont des morphismes inverses l'un de l'autre
$a' : (\Spec(k), 1, -1) \to (X, c_2, 0)$ et
$b' : (X, c_2, 0) \to (\Spec(k), 1, -1)$.
\end{proof}

\begin{remark}[Motifs de Lefschetz et de Tate]
\label{remark-lefschetz-tate}
Soient $X = \mathbf{P}^1_k$ et $c_2$ comme dans l'exemple \ref{example-decompose-P1}.
Dans la littérature, le motif $(X, c_2, 0)$ est parfois appelé
{\it motif de Lefschetz} et, selon les références, les notations
$L$, $\mathbf{L}$, $\mathbf{Q}(-1)$ ou $h^2(\mathbf{P}^1_k)$
peuvent le désigner. D'après le lemme \ref{lemma-inverse-h2}, le motif de Lefschetz
est isomorphe à $\mathbf{1}(-1)$. Il est donc
inversible (Catégories, définition \ref{categories-definition-invertible}),
d'inverse
$\mathbf{1}(1)$. Le motif $\mathbf{1}(1)$ est parfois appelé
{\it motif de Tate} et, selon les références, les notations
$L^{-1}$, $\mathbf{L}^{-1}$, $\mathbf{T}$ ou $\mathbf{Q}(1)$
peuvent le désigner.
\end{remark}

\begin{lemma}
\label{lemma-additive}
La catégorie $M_k$ est additive.
\end{lemma}

\begin{proof}
Soient $(Y, p, m)$ et $(Z, q, n)$ des motifs. Si $n = m$, une
somme directe est donnée par $(Y \amalg Z, p + q, m)$, avec des notations évidentes.
Nous omettons les détails.

\medskip\noindent
Supposons $n < m$. Soient $X$ et $c_2$ comme dans
l'exemple \ref{example-decompose-P1}. Considérons alors
\begin{align*}
(Z, q, n)
& =
(Z, q, m) \otimes (\Spec(k), 1, -1) \otimes \ldots \otimes
(\Spec(k), 1, -1) \\
& \cong
(Z, q, m) \otimes (X, c_2, 0) \otimes \ldots \otimes (X, c_2, 0) \\
& \cong
(Z \times X^{m - n}, q \otimes c_2 \otimes \ldots \otimes c_2, m)
\end{align*},
où nous avons utilisé le lemme \ref{lemma-inverse-h2}.
Nous sommes ainsi ramenés au cas traité dans le premier paragraphe.
\end{proof}

\begin{lemma}
\label{lemma-decompose-P1}
Dans $M_k$, on a $h(\mathbf{P}^1_k) \cong \mathbf{1} \oplus \mathbf{1}(-1)$.
\end{lemma}

\begin{proof}
Cela résulte de l'exemple \ref{example-decompose-P1} et du
lemme \ref{lemma-inverse-h2}.
\end{proof}

\begin{lemma}
\label{lemma-characterize-motives}
Soient $X$ et $c_2$ comme dans l'exemple \ref{example-decompose-P1}.
Soit $\mathcal{C}$ une catégorie monoïdale symétrique karoubienne
$\mathbf{Q}$-linéaire. Tout foncteur $\mathbf{Q}$-linéaire
$$
F :
\left\{
\begin{matrix}
\text{schémas projectifs lisses sur }k\\
\text{les morphismes sont les correspondances de degré }0
\end{matrix}
\right\}
\longrightarrow
\mathcal{C}
$$
de catégories monoïdales symétriques tel que l'image de $F(c_2)$ dans
$F(X)$ soit un objet inversible se factorise de manière unique par un foncteur
$F : M_k \to \mathcal{C}$ de catégories monoïdales symétriques.
\end{lemma}

\begin{proof}
Notons $U$ l'objet inversible de $\mathcal{C}$ dont l'existence
est supposée dans l'énoncé du lemme. Nous prolongeons $F$ aux motifs en posant
$$
F(X, p, m) = \left(\text{l'image du
projecteur }F(p)\text{ dans }F(X)\right) \otimes U^{\otimes -m}
$$,
ce qui a un sens puisque $U$ est inversible et que $\mathcal{C}$
est karoubienne. Une propriété importante de ce choix est que
$F(X, c_2, 0) = U$. Observons que
\begin{align*}
F((X, p, m) \otimes (Y, q, n))
& =
F(X \times Y, p \otimes q, m + n) \\
& =
\left(\text{l'image de }F(p \otimes q)\text{ dans }F(X \times Y)\right)
\otimes U^{\otimes -m - n} \\
& =
F(X, p, m) \otimes F(Y, q, n)
\end{align*}
Notre règle est donc compatible avec les produits tensoriels
au niveau des objets (nous omettons les détails).

\medskip\noindent
Prolongeons ensuite $F$ aux morphismes de motifs. Supposons que
$$
a \in
\Hom((Y, p, m), (Z, q, n)) =
q \circ \text{Corr}^{n - m}(Y, Z) \circ p \subset \text{Corr}^{n - m}(Y, Z)
$$
soit un morphisme. Si $n = m$, alors $a$ est une correspondance de degré $0$
et l'on peut utiliser $F(a) : F(Y) \to F(Z)$ pour obtenir l'application voulue
$F(Y, p, m) \to F(Z, q, n)$. Si $n < m$, on a des identifications canoniques
\begin{align*}
s : F((Z, q, n))
& \to
F(Z, q, m) \otimes U^{m - n} \\
& \to
F(Z, q, m) \otimes F(X, c_2, 0) \otimes \ldots \otimes F(X, c_2, 0) \\
& \to
F((Z, q, m) \otimes (X, c_2, 0) \otimes \ldots \otimes (X, c_2, 0)) \\
& \to
F((Z \times X^{m - n}, q \otimes c_2 \otimes \ldots \otimes c_2, m))
\end{align*}
En effet, pour le premier isomorphisme, on utilise la définition de $F$ sur les motifs
donnée ci-dessus. Pour le deuxième, on utilise le choix de $U$. Pour le troisième,
on utilise la compatibilité de $F$ avec les produits tensoriels de motifs.
Le quatrième est la définition du produit tensoriel des motifs. D'autre part,
on a de même un isomorphisme
$$
\sigma : (Z, q, n) \to
(Z \times X^{m - n}, q \otimes c_2 \otimes \ldots \otimes c_2, m)
$$
(voir la démonstration du lemme \ref{lemma-additive}). En composant $a$ avec cet isomorphisme,
on obtient
$$
\sigma \circ a \in
\Hom((Y, p, m),
(Z \times X^{m - n}, q \otimes c_2 \otimes \ldots \otimes c_2, m))
$$
En réunissant ces constructions, on obtient
$$
s^{-1} \circ F(\sigma \circ a) :
F(Y, p, m) \to
F(Z, q, n)
$$
Si $n > m$, on définit de même des isomorphismes
$$
t : F((Y, p, m)) \to
F((Y \times X^{n - m}, p \otimes c_2 \otimes \ldots \otimes c_2, n))
$$
et
$$
\tau : (Y, p, m)) \to
(Y \times X^{n - m}, p \otimes c_2 \otimes \ldots \otimes c_2, n)
$$,
et l'on pose $F(a) = F(a \circ \tau^{-1}) \circ t$.
Nous omettons la vérification que cette construction définit un foncteur
de catégories monoïdales symétriques.
\end{proof}

\begin{lemma}
\label{lemma-dual}
Soit $X$ un schéma projectif lisse sur $k$, équidimensionnel
de dimension $d$. Alors $h(X)(d)$ est un dual à gauche de $h(X)$ dans $M_k$.
\end{lemma}

\begin{proof}
Nous utiliserons le lemme \ref{lemma-composition-correspondences}
sans le mentionner de nouveau. On calcule
$$
\Hom(\mathbf{1}, h(X) \otimes h(X)(d)) =
\text{Corr}^d(\Spec(k), X \times X) = \CH^d(X \times X)
$$
On dispose ici de $\eta = [\Delta]$. D'autre part, on a
$$
\Hom(h(X)(d) \otimes h(X), \mathbf{1}) =
\text{Corr}^{-d}(X \times X, \Spec(k)) = \CH_d(X \times X)
$$,
où l'on dispose également de la classe $\epsilon = [\Delta]$
de la diagonale. La composition de la correspondance
$[\Delta] \otimes 1$ avec $1 \otimes [\Delta]$, dans un ordre ou dans l'autre,
est la correspondance $[\Delta] = 1$ dans $\text{Corr}^0(X, X)$ ; les diagrammes de
Catégories, définition \ref{categories-definition-dual},
sont donc commutatifs.
En effet, observons que
$$
[\Delta] \otimes 1 \in \text{Corr}^d(X, X \times X \times X) =
\CH^{2d}(X \times X \times X \times X)
$$
est donnée par la classe du cycle
$\text{pr}^{1234, -1}_{23}(\Delta) \cap \text{pr}^{1234, -1}_{14}(\Delta)$, avec
des notations évidentes. De même, la classe
$$
1 \otimes [\Delta] \in \text{Corr}^{-d}(X \times X \times X, X) =
\CH^{2d}(X \times X \times X \times X)
$$
est donnée par la classe du cycle
$\text{pr}^{1234, -1}_{23}(\Delta) \cap \text{pr}^{1234, -1}_{14}(\Delta)$.
La composée $(1 \otimes [\Delta]) \circ ([\Delta] \otimes 1)$
est, par définition, l'image directe par $\text{pr}^{12345}_{15, *}$
du produit d'intersection
$$
[\text{pr}^{12345, -1}_{23}(\Delta) \cap \text{pr}^{12345, -1}_{14}(\Delta)]
\cdot
[\text{pr}^{12345, -1}_{34}(\Delta) \cap \text{pr}^{12345, -1}_{15}(\Delta)]
=
[\text{petite diagonale dans } X^5]
$$,
qui est égale à $\Delta$, comme voulu. Nous omettons la démonstration
de la formule pour la composition dans l'autre ordre.
\end{proof}

\begin{lemma}
\label{lemma-dual-general}
Tout objet de $M_k$ admet un dual à gauche.
\end{lemma}

\begin{proof}
Soit $M = (X, p, m)$ un objet de $M_k$. Alors $M$ est un facteur direct de
$(X, 1, m) = h(X)(m)$.
D'après Homologie, lemme \ref{homology-lemma-Karoubian-dual},
il suffit de montrer que
$h(X)(m) = h(X) \otimes \mathbf{1}(m)$ admet un dual.
Par construction, $\mathbf{1}(-m)$ est un dual à gauche de $\mathbf{1}(m)$.
Il suffit donc de montrer que $h(X)$ admet un dual à gauche ; voir
Catégories, lemme \ref{categories-lemma-tensor-dual}.
Soit $X = \coprod X_i$ la décomposition de $X$ en
composantes irréductibles. Alors $h(X) = \bigoplus h(X_i)$,
et il suffit de montrer que $h(X_i)$ admet un dual à gauche ; voir
Homologie, lemme \ref{homology-lemma-additive-dual}.
Cela résulte du lemme \ref{lemma-dual}.
\end{proof}






\section{Groupes de Chow des motifs}
\label{section-chow-groups-motives}

\noindent
Nous définissons les groupes de Chow d'un motif comme suit.

\begin{definition}
\label{definition-chow-group-motives}
Soit $k$ un corps de base et soit $M = (X, p, m)$ un motif de Chow sur $k$.
Pour $i \in \mathbf{Z}$, nous définissons le {\it $i$-ième groupe de Chow de $M$}
par la formule
$$
\CH^i(M) = p\left(\CH^{i + m}(X) \otimes \mathbf{Q}\right)
$$
\end{definition}

\noindent
On a $\CH^i(h(X)) = \CH^i(X) \otimes \mathbf{Q}$
si $X$ est un schéma projectif lisse sur $k$.

\medskip\noindent
Observons que $\CH^i(-)$ est un foncteur de $M_k$ dans les espaces vectoriels sur $\mathbf{Q}$.
En effet, si $c : M \to N$ est un morphisme de motifs
$M = (X, p, m)$ et $N = (Y, q, n)$, alors $c$ est une correspondance de
degré $n - m$ de $X$ dans $Y$ ; l'image directe par $c$
(section \ref{section-correspondences}) est donc une famille d'applications
$$
c_* :
\CH^{i + m}(X) \otimes \mathbf{Q}
\longrightarrow
\CH^{i + n}(Y) \otimes \mathbf{Q}
$$
Puisque $c = q \circ c \circ p$ par définition des morphismes de motifs,
on obtient bien
$$
c_* : \CH^i(M) \to \CH^i(N)
$$
pour tout $i \in \mathbf{Z}$. Cette construction est compatible avec la composition des
morphismes de motifs, d'après le lemme \ref{lemma-composition-correspondences}.
Cette fonctorialité des groupes de Chow peut également se déduire
du lemme suivant.

\begin{lemma}
\label{lemma-chow-groups-representable}
Soit $k$ un corps de base. Le foncteur $\CH^i(-)$ sur la catégorie
des motifs $M_k$ est représentable par $\mathbf{1}(-i)$ ; autrement dit,
on a
$$
\CH^i(M) = \Hom_{M_k}(\mathbf{1}(-i), M)
$$,
fonctoriellement en $M$ dans $M_k$.
\end{lemma}

\begin{proof}
Cela résulte immédiatement des définitions et du
lemme \ref{lemma-composition-correspondences}.
\end{proof}

\noindent
Le lecteur peut prévoir que le
lemme \ref{lemma-chow-groups-representable}, le lemme de Yoneda et
la dualité du lemme \ref{lemma-dual} permettent d'obtenir le résultat suivant.

\begin{lemma}[Manin]
\label{lemma-manin}
Soit $k$ un corps de base et soit $c : M \to N$ un morphisme de motifs.
Si, pour tout schéma projectif lisse $X$ sur $k$, l'application
$c \otimes 1 : M \otimes h(X) \to N \otimes h(X)$ induit un isomorphisme sur
les groupes de Chow, alors $c$ est un isomorphisme.
\end{lemma}

\begin{proof}
Tout objet $L$ de $M_k$ est un facteur direct de $h(X)(m)$ pour un schéma projectif lisse
$X$ sur $k$ et un entier $m \in \mathbf{Z}$ convenables. Observons que les groupes de Chow
de $M \otimes h(X)(m)$ sont les mêmes que ceux de $M \otimes h(X)$,
à un décalage des degrés près. Notre hypothèse entraîne donc
que $c \otimes 1 : M \otimes L \to N \otimes L$ induit un isomorphisme sur
les groupes de Chow pour tout objet $L$ de $M_k$. D'après le
lemme \ref{lemma-chow-groups-representable},
l'application
$$
\Hom_{M_k}(\mathbf{1}, M \otimes L) \to
\Hom_{M_k}(\mathbf{1}, N \otimes L)
$$
est un isomorphisme pour tout $L$. Tout objet de $M_k$ admettant un dual à gauche
(lemme \ref{lemma-dual-general}), il en résulte que
$$
\Hom_{M_k}(K, M) \to \Hom_{M_k}(K, N)
$$
est un isomorphisme pour tout objet $K$ de $M_k$ ; voir
Catégories, lemme \ref{categories-lemma-left-dual}.
On conclut par le lemme de Yoneda
(Catégories, lemme \ref{categories-lemma-yoneda}).
\end{proof}




\section{Formule du fibré projectif}
\label{section-projective-space-bundle}

\noindent
Soit $k$ un corps de base et soit $X$ un schéma projectif lisse sur $k$.
Soit $\mathcal{E}$ un $\mathcal{O}_X$-module localement libre de rang $r$.
Nous adoptons la convention selon laquelle le {\it fibré projectif associé à
$\mathcal{E}$} est le morphisme
$$
\xymatrix{
P = \mathbf{P}(\mathcal{E}) =
\underline{\text{Proj}}_X(\text{Sym}^*(\mathcal{E}))
\ar[r]^-p
& X
}
$$
sur $X$, où $\mathcal{O}_P(1)$ est normalisé de telle sorte que
$p_*(\mathcal{O}_P(1)) = \mathcal{E}$. Rappelons que
$$
[\Gamma_p] \in \text{Corr}^0(X, P) \subset \CH^*(X \times P) \otimes \mathbf{Q}
$$
Voir l'exemple \ref{example-graph-correspondence}.
Pour $i = 0, \ldots, r - 1$, considérons les correspondances
$$
c_i = c_1(\text{pr}_2^*\mathcal{O}_P(1))^i \cap [\Gamma_p]
\in \text{Corr}^i(X, P)
$$
Nous considérons désormais $c_i$ comme un morphisme $h(X)(-i) \to h(P)$.

\begin{lemma}[Formule du fibré projectif]
\label{lemma-projective-space-bundle-formula}
Dans la situation ci-dessus, l'application
$$
\sum\nolimits_{i = 0, \ldots, r - 1} c_i :
\bigoplus\nolimits_{i = 0, \ldots, r - 1} h(X)(-i)
\longrightarrow
h(P)
$$
est un isomorphisme dans la catégorie des motifs.
\end{lemma}

\begin{proof}
D'après le lemme \ref{lemma-manin}, il suffit de montrer que
notre application définit un isomorphisme sur les groupes de Chow des motifs
après produit par un schéma projectif lisse quelconque $Z$.
Observons que $P \times Z \to X \times Z$ est le fibré
projectif associé à l'image réciproque de $\mathcal{E}$ sur $X \times Z$.
L'assertion sur les groupes de Chow résulte donc de la formule
du fibré projectif donnée dans
Homologie de Chow, lemme \ref{chow-lemma-chow-ring-projective-bundle}.
En effet, l'image directe des cycles par $[\Gamma_p]$ est donnée par l'image réciproque
des cycles par $p$, d'après le lemme \ref{lemma-functor-and-cycles} et
Homologie de Chow, lemme \ref{chow-lemma-lci-gysin-flat}. L'image directe
par $c_i$ envoie donc $\alpha$ sur $c_1(\mathcal{O}_P(1))^i \cap p^*\alpha$.
Nous omettons certains détails.
\end{proof}

\noindent
Dans la situation ci-dessus, pour $j = 0, \ldots, r - 1$, considérons
les correspondances
$$
c'_j = c_1(\text{pr}_1^*\mathcal{O}_P(1))^{r - 1 - j} \cap [\Gamma_p^t] \in
\text{Corr}^{-j}(P, X)
$$
Pour $i, j \in \{0, \ldots, r - 1\}$, on a
$$
c'_j \circ c_i =
\text{pr}_{13, *}\left(
c_1(\text{pr}_2^*\mathcal{O}_P(1))^{i + r - 1 - j} \cap
(\text{pr}_{12}^*[\Gamma_p] \cdot \text{pr}_{23}^*[\Gamma_p^t])
\right)
$$
Les cycles $\text{pr}_{12}^{-1}\Gamma_p$ et
$\text{pr}_{23}^{-1}\Gamma_p^t$ se coupent transversalement,
et leur intersection est l'image de
$(p, 1, p) : P \to X \times P \times X$.
Observons que les fibres de
$(p, p) = \text{pr}_{13} \circ (p, 1, p) : P \to X \times X$
sont de dimension $r - 1$. On en déduit immédiatement
$c'_j \circ c_i = 0$ pour $i + r - 1 - j < r - 1$, autrement dit lorsque $i < j$.
D'autre part, d'après la formule du fibré projectif
(Homologie de Chow, lemme \ref{chow-lemma-chow-ring-projective-bundle}),
le cycle $c_1(\mathcal{O}_P(1))^{r - 1} \cap [P]$ a pour image
$[X]$ dans $X$. Ainsi, pour $i = j$, l'image directe ci-dessus
donne la classe de la diagonale ; on a donc
$$
c'_i \circ c_i = 1 \in \text{Corr}^0(X, X)
$$
pour tout $i \in \{0, \ldots, r - 1\}$. La matrice
de la composée
$$
\bigoplus h(X)(-i)
\xrightarrow{\bigoplus c_i}
h(P)
\xrightarrow{\bigoplus c'_j}
\bigoplus h(X)(-j)
$$
est donc inversible (elle est triangulaire supérieure, avec des $1$ sur la diagonale).
La formule du fibré projectif
(lemme \ref{lemma-projective-space-bundle-formula})
entraîne que la composée dans l'autre ordre est elle aussi
inversible, mais une démonstration directe semble un peu plus difficile.

\begin{lemma}
\label{lemma-diagonal-projective-bundle}
Soit $p : P \to X$ comme dans le lemme \ref{lemma-projective-space-bundle-formula}.
La classe $[\Delta_P]$ de la diagonale de $P$ dans $\CH^*(P \times P)$
s'écrit
$$
[\Delta_P] =
\left(\sum\nolimits_{i = 0, \ldots, r - 1}
{r - 1 \choose i} c_{r - 1 - i}(\text{pr}_1^*\mathcal{S}^\vee) \cap
c_1(\text{pr}_2^*\mathcal{O}_P(1))^i\right)
\cap
(p \times p)^*[\Delta_X]
$$,
où $\mathcal{S}$ est le noyau de la surjection canonique
$p^*\mathcal{E} \to \mathcal{O}_P(1)$.
\end{lemma}

\begin{proof}
Observons que $(p \times p)^*[\Delta_X] = [P \times_X P]$.
Puisque $\Delta_P \subset P \times_X P \subset P \times P$
et que le cap-produit par les classes de Chern commute à l'image directe propre
(Homologie de Chow, lemme \ref{chow-lemma-pushforward-cap-cj}),
il suffit de montrer que la classe de
$\Delta_P \subset P \times_X P$ dans $\CH^*(P \times_X P)$
est égale à
$$
\left(\sum\nolimits_{i = 0, \ldots, r - 1}
{r - 1 \choose i} c_{r - 1 - i}(q_1^*\mathcal{S}^\vee) \cap
c_1(q_2^*\mathcal{O}_P(1))^i\right)
\cap
[P \times_X P]
$$,
où $q_i : P \times_X P \to P$, pour $i = 1, 2$, sont les projections.
Posons $q = p \circ q_1 = p \circ q_2 : P \times_X P \to X$.
Considérons les applications
$$
q_1^*\mathcal{S} \otimes q_2^*\mathcal{O}_P(-1) \to
q^*\mathcal{E} \otimes q^*\mathcal{E}^\vee \to
\mathcal{O}_{P \times_X P}
$$,
où la dernière flèche est l'image réciproque par $q$ de l'application d'évaluation
$\mathcal{E} \otimes_{\mathcal{O}_X} \mathcal{E}^\vee \to \mathcal{O}_X$.
La source de la composée est un module localement libre de rang $r - 1$,
et un calcul local montre que cette application s'annule exactement le long de
$\Delta_P$. D'après Homologie de Chow, lemme \ref{chow-lemma-top-chern-class},
la classe $[\Delta_P]$ est la classe de Chern de degré maximal du dual
$$
q_1^*\mathcal{S}^\vee \otimes q_2^*\mathcal{O}_P(1)
$$
Le résultat voulu découle de Homologie de Chow, lemme
\ref{chow-lemma-chern-classes-E-tensor-L}.
\end{proof}








\section{Théories cohomologiques de Weil classiques}
\label{section-axioms-classical}

\noindent
Nous définissons dans cette section ce que nous appellerons une théorie
cohomologique de Weil classique. Il s'agit exactement de ce qui est appelé
une théorie cohomologique de Weil dans \cite[Section 1.2]{Kleiman-cycles}.

\medskip\noindent
Fixons un corps algébriquement clos $k$ (le corps de base).
Dans cette section, {\it variété} signifie variété sur $k$ ; voir
Variétés, section \ref{varieties-section-varieties}.
Fixons un corps $F$ de caractéristique $0$ (le corps de coefficients).
Une théorie cohomologique de Weil est définie par des données (D1), (D2) et (D3)
soumises aux axiomes (A), (B) et (C).

\medskip\noindent
Les données sont les suivantes :
\begin{enumerate}
\item[(D1)] Un foncteur contravariant $H^*$ de la catégorie
des variétés projectives lisses dans la catégorie des
$F$-algèbres graduées commutatives.
\item[(D2)] Pour toute variété projective lisse $X$,
un homomorphisme de groupes $\gamma : \CH^i(X) \to H^{2i}(X)$.
\item[(D3)] Pour toute variété projective lisse $X$ de dimension $d$,
une application $\int_X : H^{2d}(X) \to F$.
\end{enumerate}
Précisons le sens de ces données par quelques remarques
et introduisons la terminologie correspondante.

\medskip\noindent
Remarques sur (D1). Pour une variété projective lisse $X$,
on appelle $H^*(X)$ la {\it cohomologie} de $X$. Pour un morphisme
$f : X \to Y$ de variétés projectives lisses, on note
$f^* : H^*(Y) \to H^*(X)$ l'application $H^*(f)$ et on l'appelle
l'{\it application image réciproque}.

\medskip\noindent
Remarques sur (D2). L'application $\gamma$ est appelée l'{\it application classe de cycle}.
On appelle $\gamma(\alpha)$ la {\it classe de cohomologie} de $\alpha$.
Si $Z \subset Y \subset X$ sont des sous-schémas fermés, avec $Y$ et $X$
des variétés projectives lisses et $Z$ intègre, alors $[Z]$ peut
désigner la classe du cycle $[Z]$ dans $\CH^*(Y)$ ou dans $\CH^*(X)$.
Dans ce cas, la notation $\gamma([Z])$ est ambiguë, et son sens
doit être déduit du contexte.

\medskip\noindent
Remarques sur (D3). L'application $\int_X$ est parfois appelée
l'{\it application trace} et est parfois notée $\text{Tr}_X$.

\medskip\noindent
Le premier axiome est souvent appelé {\it dualité de Poincar\'e} :
\begin{enumerate}
\item[(A)] Soit $X$ une variété projective lisse de dimension $d$. Alors :
\begin{enumerate}
\item $\dim_F H^i(X) < \infty$ pour tout $i$ ;
\item $H^i(X) \times H^{2d - i}(X) \rightarrow H^{2d}(X) \rightarrow F$
est un accouplement parfait pour tout $i$, la dernière
application étant l'application trace $\int_X$ ;
\item $H^i(X) = 0$ sauf si $i \in [0, 2d]$ ;
\item $\int_X : H^{2d}(X) \to F$ est un isomorphisme.
\end{enumerate}
\end{enumerate}
Soit $f : X \to Y$ un morphisme de variétés projectives lisses,
avec $\dim(X) = d$ et $\dim(Y) = e$. À l'aide de la dualité de Poincar\'e,
on peut définir une {\it image directe}
$$
f_* : H^{2d - i}(X) \longrightarrow H^{2e - i}(Y)
$$
comme l'application contragrédiente de l'application linéaire $f^* : H^i(Y) \to H^i(X)$.
Explicitement, pour $a \in H^{2d - i}(X)$, l'élément $f_*a \in H^{2e - i}(Y)$
est caractérisé par
$$
\int_X f^*b \cup a = \int_Y b \cup f_*a
$$
pour tout $b \in H^i(Y)$.

\begin{lemma}
\label{lemma-pushforward-classical}
Supposons que les données (D1) et (D3) satisfassent à (A). Pour $f : X \to Y$
un morphisme de variétés projectives lisses, on a
$f_*(f^*b \cup a) = b \cup f_*a$. Si $g : Y \to Z$ est un second morphisme
de variétés projectives lisses, alors $g_* \circ f_* = (g \circ f)_*$.
\end{lemma}

\begin{proof}
La première égalité résulte de
$$
\int_Y c \cup b \cup f_*a =
\int_X f^*c \cup f^*b \cup a =
\int_Y c \cup f_*(f^*b \cup a).
$$
La seconde résulte de
$$
\int_Z c \cup (g \circ f)_*a = \int_X (g \circ f)^*c \cup a =
\int_X f^* g^* c \cup a = \int_Y g^*c \cup f_*a = \int_Z c \cup g_*f_*a
$$
Cela achève la démonstration.
\end{proof}

\noindent
Le deuxième axiome exprime que $H^*$ respecte la structure monoïdale
donnée par les produits, au moyen de la {\it formule de K\"unneth} :
\begin{enumerate}
\item[(B)] Soient $X$ et $Y$ des variétés projectives lisses. L'application
$$
H^*(X) \otimes_F H^*(Y) \to H^*(X \times Y),\quad
a \otimes b \mapsto \text{pr}_1^*a \cup \text{pr}_2^*b
$$
est un isomorphisme.
\end{enumerate}
Le troisième axiome porte sur les applications classes de cycles :
\begin{enumerate}
\item[(C)] Les applications classes de cycles satisfont aux règles suivantes :
\begin{enumerate}
\item pour un morphisme $f : X \to Y$ de variétés projectives lisses,
on a $\gamma(f^!\beta) = f^*\gamma(\beta)$ pour $\beta \in \CH^*(Y)$ ;
\item pour un morphisme $f : X \to Y$ de variétés projectives lisses, on a
$\gamma(f_*\alpha) = f_*\gamma(\alpha)$ pour $\alpha \in \CH^*(X)$ ;
\item pour toute variété projective lisse $X$, on a
$\gamma(\alpha \cdot \beta) = \gamma(\alpha) \cup \gamma(\beta)$
pour $\alpha, \beta \in \CH^*(X)$ ;
\item $\int_{\Spec(k)} \gamma([\Spec(k)]) = 1$.
\end{enumerate}
\end{enumerate}

\begin{remark}
\label{remark-replace-cup-product-classical}
Soit $X$ une variété projective lisse. On obtient des applications
$$
H^*(X) \otimes_F H^*(X) \longrightarrow H^*(X \times X)
\xrightarrow{\Delta^*} H^*(X)
$$,
où la première flèche est celle de l'axiome (B) et $\Delta^*$
est l'image réciproque par le morphisme diagonal $\Delta : X \to X \times X$.
La composée est le cup-produit, puisque l'image réciproque est un homomorphisme d'algèbres
et que $\text{pr}_i \circ \Delta = \text{id}$.
D'autre part, pour des cycles $\alpha, \beta$ sur $X$,
le produit d'intersection est défini par la formule
$$
\alpha \cdot \beta =
\Delta^!(\alpha \times \beta)
$$
Autrement dit, $\alpha \cdot \beta$ est l'image réciproque du
produit extérieur $\alpha \times \beta$ sur $X \times X$ par
la diagonale. Notons aussi que
$\alpha \times \beta = \text{pr}_1^*\alpha \cdot \text{pr}_2^*\beta$
dans $\CH^*(X \times X)$ (nous omettons la démonstration). Ainsi, sous l'axiome (C)(a),
l'axiome (C)(c) équivaut à la compatibilité de $\gamma$
avec le produit extérieur, au sens où
$\gamma(\alpha \times \beta)$ est égal à
$\text{pr}_1^*\gamma(\alpha) \cup \text{pr}_2^*\gamma(\beta)$.
C'est ainsi que l'axiome (C)(c) est formulé dans \cite{Kleiman-cycles}.
\end{remark}

\begin{definition}
\label{definition-weil-cohomology-theory-classical}
Soit $k$ un corps algébriquement clos.
Soit $F$ un corps de caractéristique $0$.
Une {\it théorie cohomologique de Weil classique} sur $k$ à coefficients dans $F$
est définie par des données (D1), (D2) et (D3) satisfaisant
à la dualité de Poincar\'e, à la formule de K\"unneth et à la compatibilité
avec les classes de cycles, plus précisément aux axiomes (A), (B) et (C).
\end{definition}

\noindent
Établissons quelques conséquences élémentaires.

\begin{lemma}
\label{lemma-degrees-cycles-classical}
Soit $H^*$ une théorie cohomologique de Weil classique
(définition \ref{definition-weil-cohomology-theory-classical}).
Soit $X$ une variété projective lisse de dimension $d$. Le diagramme
$$
\xymatrix{
\CH^d(X) \ar[r]_-\gamma \ar@{=}[d] &
H^{2d}(X) \ar[d]^{\int_X} \\
\CH_0(X) \ar[r]^\deg & F
}
$$
est commutatif, où $\deg : \CH_0(X) \to \mathbf{Z}$ est le degré des
zéro-cycles étudié dans Homologie de Chow, section
\ref{chow-section-degree-zero-cycles}.
\end{lemma}

\begin{proof}
Le résultat vaut pour $\Spec(k)$ d'après l'axiome (C)(d). Soit $x : \Spec(k) \to X$
un point fermé de $X$. Alors $\gamma([x]) = x_*\gamma([\Spec(k)])$
dans $H^{2d}(X)$, d'après l'axiome (C)(b). Par conséquent, $\int_X \gamma([x]) = 1$
par définition de $x_*$.
\end{proof}

\begin{lemma}
\label{lemma-trace-product-classical}
Soit $H^*$ une théorie cohomologique de Weil classique
(définition \ref{definition-weil-cohomology-theory-classical}).
Soient $X$ et $Y$ des variétés projectives lisses.
Alors $\int_{X \times Y} = \int_X \otimes \int_Y$.
\end{lemma}

\begin{proof}
Posons $\dim(X) = d$ et $\dim(Y) = e$. D'après l'axiome (B), on a
$H^{2d + 2e}(X \times Y) = H^{2d}(X) \otimes H^{2e}(Y)$,
et cet espace est de dimension $1$ d'après l'axiome (A)(d).
D'après le lemme \ref{lemma-degrees-cycles-classical},
cet espace vectoriel de dimension $1$ est engendré par
la classe $\gamma([x \times y])$ d'un point fermé $(x, y)$, et
$\int_{X \times Y} \gamma([x \times y]) = 1$.
Puisque $\gamma([x \times y]) = \gamma([x]) \otimes \gamma([y])$
d'après les axiomes (C)(a) et (C)(c), et puisque $\int_X \gamma([x]) = 1$ et
$\int_Y \gamma([y]) = 1$, on obtient la conclusion.
\end{proof}

\begin{lemma}
\label{lemma-pr2star-classical}
Soit $H^*$ une théorie cohomologique de Weil classique
(définition \ref{definition-weil-cohomology-theory-classical}).
Soient $X$ et $Y$ des variétés projectives lisses.
Alors $\text{pr}_{2, *} : H^*(X \times Y) \to H^*(Y)$
envoie $a \otimes b$ sur $(\int_X a) b$.
\end{lemma}

\begin{proof}
Cela équivaut au résultat du lemme \ref{lemma-trace-product-classical}.
\end{proof}

\begin{lemma}
\label{lemma-class-diagonal-classical}
Soit $H^*$ une théorie cohomologique de Weil classique
(définition \ref{definition-weil-cohomology-theory-classical}).
Soit $X$ une variété projective lisse de dimension $d$.
Choisissons une base $e_{i, j}, j = 1, \ldots, \beta_i$ de $H^i(X)$ sur $F$.
À l'aide de la formule de K\"unneth, écrivons
$$
\gamma([\Delta]) =
\sum\nolimits_{i = 0, \ldots, 2d}
\sum\nolimits_j e_{i, j} \otimes e'_{2d - i , j}
\quad\text{dans}\quad
\bigoplus\nolimits_i H^i(X) \otimes_F H^{2d - i}(X)
$$
avec $e'_{2d - i, j} \in H^{2d - i}(X)$.
Alors $\int_X e_{i, j} \cup e'_{2d - i, j'} = (-1)^i\delta_{jj'}$.
\end{lemma}

\begin{proof}
Rappelons que $\Delta^* : H^*(X \times X) \to H^*(X)$ est l'application
cup-produit $H^*(X) \otimes_F H^*(X) \to H^*(X)$ ; voir
la remarque \ref{remark-replace-cup-product-classical}. D'autre part, on a
$\gamma([\Delta]) = \Delta_*\gamma([X]) = \Delta_*1$ d'après
l'axiome (C)(b) et le fait que $\gamma([X]) = 1$. En effet,
$[X] \cdot [X] = [X]$ ; ainsi, par l'axiome (C)(c), la classe de cohomologie
$\gamma([X])$ vaut $0$ ou $1$ dans l'algèbre de dimension $1$ sur $F$, $H^0(X)$ ;
nous avons aussi utilisé ici les axiomes (A)(d) et (A)(b).
Mais $\gamma([X])$ ne peut être nulle, car $[X] \cdot [x] = [x]$
pour un point fermé $x$ de $X$, et $\gamma([x])$
est non nulle d'après le lemme \ref{lemma-degrees-cycles-classical}.
Par conséquent,
$$
\int_{X \times X} \gamma([\Delta]) \cup a \otimes b =
\int_{X \times X} \Delta_*1 \cup a \otimes b =
\int_X a \cup b
$$,
par définition de $\Delta_*$. D'autre part, on a
$$
\int_{X \times X} (\sum e_{i, j} \otimes e'_{2d -i , j}) \cup a \otimes b =
\sum (\int_X a \cup e_{i, j})(\int_X e'_{2d - i, j} \cup b)
$$,
d'après le lemme \ref{lemma-trace-product-classical} ; notons que nous avons
effectué deux échanges, de sorte que le signe est $1$.
Ainsi, si l'on choisit $a$ de sorte que $\int_X a \cup e_{i, j} = 1$
et que tous les autres accouplements soient nuls, on en déduit
$\int_X e'_{2d - i, j} \cup b = \int_X a \cup b$ pour tout $b$, c'est-à-dire
$e'_{2d - i, j} = a$. Le lemme est démontré.
\end{proof}

\begin{lemma}
\label{lemma-square-diagonal-classical}
Soit $H^*$ une théorie cohomologique de Weil classique
(définition \ref{definition-weil-cohomology-theory-classical}).
Soit $X$ une variété projective lisse. On a
$$
\sum\nolimits_{i = 0, \ldots, 2\dim(X)} (-1)^i\dim_F H^i(X) =
\deg([\Delta] \cdot [\Delta]) = \deg(c_d(\mathcal{T}_X) \cap [X])
$$
\end{lemma}

\begin{proof}
Démontrons l'égalité de droite. On a
$[\Delta] \cdot [\Delta] = \Delta_*(\Delta^![\Delta])$
(Homologie de Chow, lemme \ref{chow-lemma-intersect-regularly-embedded}).
Puisque $\Delta_*$ préserve les degrés des $0$-cycles, il suffit de calculer
le degré de $\Delta^![\Delta]$. La classe $\Delta^![\Delta]$ est donnée
par le cap-produit de $[\Delta]$ avec
la classe de Chern de degré maximal du faisceau normal de $\Delta \subset X \times X$
(Homologie de Chow, lemme \ref{chow-lemma-gysin-fundamental}).
Puisque le faisceau conormal de $\Delta$ est $\Omega_{X/k}$
(Morphismes, lemme \ref{morphisms-lemma-differentials-diagonal}),
le faisceau normal est égal au faisceau tangent
$\mathcal{T}_X = \SheafHom_{\mathcal{O}_X}(\Omega_{X/k}, \mathcal{O}_X)$,
comme voulu.

\medskip\noindent
Démontrons l'égalité de gauche. D'après le lemme \ref{lemma-degrees-cycles-classical}, on a
\begin{align*}
\deg([\Delta] \cdot [\Delta])
& =
\int_{X \times X} \gamma([\Delta]) \cup \gamma([\Delta]) \\
& =
\int_{X \times X} \Delta_*1 \cup \gamma([\Delta]) \\
& =
\int_{X \times X} \Delta_*(\Delta^*\gamma([\Delta])) \\
& =
\int_X \Delta^*\gamma([\Delta])
\end{align*}
Écrivons $\gamma([\Delta]) = \sum  e_{i, j} \otimes e'_{2d - i , j}$
comme dans le lemme \ref{lemma-class-diagonal-classical}.
En rappelant que $\Delta^*$ est donné par le cup-produit, on obtient
$$
\int_X \sum\nolimits_{i, j} e_{i, j} \cup e'_{2d - i, j} =
\sum\nolimits_{i, j} \int_X e_{i, j} \cup e'_{2d - i, j} =
\sum\nolimits_{i, j} (-1)^i = \sum (-1)^i\beta_i
$$,
comme voulu.
\end{proof}




\noindent
Relions maintenant les théories cohomologiques de Weil classiques aux motifs.

\begin{lemma}
\label{lemma-from-functor-to-weil-classical}
Soit $k$ un corps algébriquement clos et soit $F$ un corps de
caractéristique $0$. Considérons un foncteur $\mathbf{Q}$-linéaire
$$
G : M_k \longrightarrow \text{espaces vectoriels gradués sur }F\text{}
$$
de catégories monoïdales symétriques tel que $G(\mathbf{1}(1))$
ne soit non nul qu'en degré $-2$. On obtient alors des données (D1), (D2), (D3)
satisfaisant à tous les axiomes (A), (B), (C), sauf peut-être (A)(c) et (A)(d).
\end{lemma}

\begin{proof}
On obtient un foncteur contravariant de la catégorie des variétés
projectives lisses dans celle des espaces vectoriels gradués sur $F$
en posant $H^*(X) = G(h(X))$. Par hypothèse, on dispose d'un isomorphisme
canonique
$$
H^*(X \times Y) = G(h(X \times Y)) = G(h(X) \otimes h(Y)) =
G(h(X)) \otimes G(h(Y)) = H^*(X) \otimes H^*(Y)
$$,
compatible avec les images réciproques. L'image réciproque par la diagonale
$\Delta : X \to X \times X$ donne une application canonique
$$
H^*(X) \otimes H^*(X) = H^*(X \times X) \to H^*(X)
$$
d'espaces vectoriels gradués, compatible avec les images réciproques.
Cela définit une structure fonctorielle de $F$-algèbre graduée sur
$H^*(X)$. Puisque $\Delta$ commute à la contrainte de commutativité
$h(X) \otimes h(X) \to h(X) \otimes h(X)$ (qui échange les facteurs),
que $G$ est un foncteur de catégories monoïdales symétriques (donc compatible avec
les contraintes de commutativité), et compte tenu de notre convention dans
Homologie, exemple \ref{homology-example-graded-vector-spaces},
on en déduit que $H^*(X)$ est une algèbre graduée
commutative ! Nous obtenons ainsi la donnée (D1).

\medskip\noindent
Puisque $\mathbf{1}(1)$ est inversible dans la catégorie des motifs,
$G(\mathbf{1}(1))$ est inversible dans la catégorie des
espaces vectoriels gradués sur $F$. Ainsi, $\sum_i \dim_F G^i(\mathbf{1}(1)) = 1$.
Par hypothèse, le seul degré où cet objet est non nul est $-2$, et l'on peut
choisir un isomorphisme $F[2] \to G(\mathbf{1}(1))$ d'espaces vectoriels gradués sur $F$.
Ici et dans la suite, $F[n]$ désigne l'espace vectoriel gradué sur $F$ ayant
$F$ en degré $-n$ et zéro ailleurs. La compatibilité avec
les produits tensoriels donne, pour tout $n \in \mathbf{Z}$, un isomorphisme
$F[2n] \to G(\mathbf{1}(n))$ compatible avec les produits tensoriels.

\medskip\noindent
Soit $X$ une variété projective lisse. D'après
le lemme \ref{lemma-composition-correspondences}, on a
$$
\CH^r(X) \otimes \mathbf{Q} = \text{Corr}^r(\Spec(k), X) =
\Hom(\mathbf{1}(-r), h(X))
$$
En appliquant le foncteur $G$, on obtient
$$
\gamma :
\CH^r(X) \otimes \mathbf{Q} \longrightarrow
\Hom(G(\mathbf{1}(-r)), H^*(X)) = H^{2r}(X)
$$
C'est la donnée (D2).

\medskip\noindent
Soit $X$ une variété projective lisse de dimension $d$. D'après
le lemme \ref{lemma-composition-correspondences}, on a
$$
\Mor(h(X)(d), \mathbf{1}) = \Mor((X, 1, d), (\Spec(k), 1, 0)) =
\text{Corr}^{-d}(X, \Spec(k)) = \CH_d(X)
$$
La classe du cycle $[X]$ dans $\CH_d(X)$ définit donc un morphisme
$h(X)(d) \to \mathbf{1}$. En appliquant $G$, on obtient
$$
H^*(X) \otimes F[-2d] = G(h(X)(d)) \longrightarrow G(\mathbf{1}) = F
$$
Cette application est nulle sauf en degré $0$, où l'on obtient
$\int_X : H^{2d}(X) \to F$. C'est la donnée (D3).

\medskip\noindent
Soit $X$ une variété projective lisse de dimension $d$.
D'après le lemme \ref{lemma-dual},
$h(X)(d)$ est un dual à gauche de $h(X)$. Par conséquent,
$G(h(X)(d)) = H^*(X) \otimes F[-2d]$ est un dual à gauche de
$H^*(X)$ dans la catégorie des espaces vectoriels gradués sur $F$.
D'après Homologie, lemme \ref{homology-lemma-left-dual-graded-vector-spaces},
on a $\sum_i \dim_F H^i(X) < \infty$, et
$\epsilon : h(X)(d) \otimes h(X) \to \mathbf{1}$ donne
des accouplements non dégénérés $H^{2d - i}(X) \otimes_F H^i(X) \to F$.
Dans la démonstration du lemme \ref{lemma-dual}, nous avons vu que
$\epsilon$ est donné par $[\Delta]$ au moyen des identifications
$$
\Hom(h(X)(d) \otimes h(X), \mathbf{1}) =
\text{Corr}^{-d}(X \times X, \Spec(k)) =
\CH_d(X \times X)
$$
Ainsi, $\epsilon$ est la composée de $[X] : h(X)(d) \to \mathbf{1}$
et de $h(\Delta)(d) : h(X)(d) \otimes h(X) \to h(X)(d)$. Il s'ensuit
que les accouplements ci-dessus sont donnés par le cup-produit suivi de
$\int_X$. Cela démontre les parties (a) et (b) de l'axiome (A).

\medskip\noindent
L'axiome (B) résulte de la compatibilité de $G$ avec les structures
tensorielles, supposée par hypothèse, et de notre construction du cup-produit ci-dessus.

\medskip\noindent
Axiome (C). Notre construction de $\gamma$ prend un cycle $\alpha$ sur $X$,
l'interprète comme une correspondance $a$ de $\Spec(k)$ dans $X$ d'un certain degré,
puis lui applique $G$. Si $f : Y \to X$ est un morphisme de variétés projectives
lisses, alors $f^!\alpha$ est l'image directe (!) de $\alpha$
par la correspondance $[\Gamma_f]$ de $X$ dans $Y$ ; voir
le lemme \ref{lemma-functor-and-cycles}. Ainsi,
$f^!\alpha$, considéré comme une correspondance de $\Spec(k)$ dans $Y$,
est égal à $a \circ [\Gamma_f]$ ; voir
le lemme \ref{lemma-composition-correspondences}.
Puisque $G$ est un foncteur, on en déduit que
$\gamma$ est compatible avec les images réciproques, ce qui est l'axiome (C)(a).

\medskip\noindent
Soit $f : Y \to X$ un morphisme de variétés projectives lisses et
soit $\beta \in \CH^r(Y)$ un cycle sur $Y$. Il faut montrer que
$$
\int_Y \gamma(\beta) \cup f^*c = \int_X \gamma(f_*\beta) \cup c
$$
pour tout $c \in H^*(X)$. Soient $a, a^t, \eta_X, \eta_Y, [X], [Y]$
comme dans le lemme \ref{lemma-prep-dual}.
Soit $b$ le cycle $\beta$ considéré comme une correspondance de $\Spec(k)$ dans $Y$
de degré $r$. Alors $f_*\beta$, considéré comme une correspondance de
$\Spec(k)$ dans $X$, est égal à $a^t \circ b$ ; voir
les lemmes \ref{lemma-functor-and-cycles} et
\ref{lemma-composition-correspondences}.
L'égalité affichée ci-dessus sera vérifiée si l'on montre que
$$
h(X) = \mathbf{1} \otimes h(X)
\xrightarrow{b \otimes 1}
h(Y)(r) \otimes h(X)
\xrightarrow{1 \otimes a}
h(Y)(r) \otimes h(Y)
\xrightarrow{\eta_Y}
h(Y)(r)
\xrightarrow{[Y]}
\mathbf{1}(r - e)
$$
est égal à
$$
h(X) = \mathbf{1} \otimes h(X)
\xrightarrow{a^t \circ b \otimes 1}
h(X)(r + d - e) \otimes h(X)
\xrightarrow{\eta_X}
h(X)(r + d - e)
\xrightarrow{[X]}
\mathbf{1}(r - e)
$$
Cela résulte immédiatement du lemme \ref{lemma-prep-dual}.
L'axiome (C)(b) est donc vérifié.

\medskip\noindent
Pour démontrer l'axiome (C)(c), nous utilisons la discussion de
la remarque \ref{remark-replace-cup-product-classical}.
Il suffit donc de démontrer que $\gamma$ est compatible avec
les produits extérieurs. Soient $X$ et $Y$ des variétés projectives lisses
et soient $\alpha$ et $\beta$ des cycles sur ces variétés. Notons
$a$ et $b$ les correspondances associées de $\Spec(k)$ dans
$X$ et $Y$. Alors $\alpha \times \beta$ correspond à
la correspondance $a \otimes b$ de $\Spec(k)$ dans $X \otimes Y = X \times Y$.
La propriété voulue résulte donc de la compatibilité de $G$
avec les structures tensorielles de part et d'autre.

\medskip\noindent
L'axiome (C)(d) résulte de ce que le cycle $[\Spec(k)]$
correspond au morphisme identique de $h(\Spec(k))$.
Cela achève la démonstration du lemme.
\end{proof}

\begin{lemma}
\label{lemma-from-weil-to-functor-classical}
Soit $k$ un corps algébriquement clos et soit $F$ un corps de
caractéristique $0$. Soit $H^*$ une théorie cohomologique de Weil classique.
On peut alors construire un foncteur $\mathbf{Q}$-linéaire
$$
G : M_k \longrightarrow \text{espaces vectoriels gradués sur }F\text{}
$$
de catégories monoïdales symétriques tel que $H^*(X) = G(h(X))$.
\end{lemma}

\begin{proof}
D'après le lemme \ref{lemma-characterize-motives}, il suffit de construire un foncteur
$G$ sur la catégorie des schémas projectifs lisses sur $k$,
dont les morphismes sont les correspondances de degré $0$, tel que
l'image de $G(c_2)$ dans $G(\mathbf{P}^1)$ soit un espace vectoriel gradué
inversible sur $F$.
Tout schéma projectif lisse étant canoniquement une réunion
disjointe de variétés projectives lisses, il suffit de construire
$G$ sur la catégorie dont les objets sont les variétés projectives lisses
et les morphismes les correspondances de degré $0$.
(Nous omettons certains détails.)

\medskip\noindent
Pour une variété projective lisse $X$, nous posons $G(X) = H^*(X)$.

\medskip\noindent
Pour une correspondance $c \in \text{Corr}^0(X, Y)$ entre variétés projectives
lisses, considérons l'application
$G(c) : G(X) = H^*(X) \to G(Y) = H^*(Y)$ donnée par la règle
$$
a \longmapsto
G(c)(a) = \text{pr}_{2, *}(\gamma(c) \cup \text{pr}_1^*a)
$$
Il est clair que $G(c)$ est additive en $c$, donc $\mathbf{Q}$-linéaire.
La compatibilité de $\gamma$ avec les images réciproques, les images directes
et les produits d'intersection, donnée par les axiomes (C)(a), (C)(b) et (C)(c),
montre que
$G(c' \circ c) = G(c') \circ G(c)$ si $c' \in \text{Corr}^0(Y, Z)$.
En effet, pour $a \in H^*(X)$, on a
\begin{align*}
(G(c') \circ G(c))(a)
& =
\text{pr}^{23}_{3, *}(\gamma(c') \cup
\text{pr}^{23, *}_2(\text{pr}^{12}_{2, *}(\gamma(c) \cup
\text{pr}^{12, *}_1a))) \\
& =
\text{pr}^{23}_{3, *}(\gamma(c') \cup
\text{pr}^{123}_{23, *}(\text{pr}^{123, *}_{12}(\gamma(c) \cup
\text{pr}^{12, *}_1 a))) \\
& =
\text{pr}^{23}_{3, *}
\text{pr}^{123}_{23, *}(
\text{pr}^{123, *}_{23}\gamma(c') \cup
\text{pr}^{123, *}_{12}\gamma(c) \cup
\text{pr}^{123, *}_1 a) \\
& =
\text{pr}^{23}_{3, *}
\text{pr}^{123}_{23, *}(
\gamma(\text{pr}^{123, *}_{23}c') \cup
\gamma(\text{pr}^{123, *}_{12}c) \cup
\text{pr}^{123, *}_1 a) \\
& =
\text{pr}^{13}_{3, *}
\text{pr}^{123}_{13, *}(
\gamma(\text{pr}^{123, *}_{23}c' \cdot \text{pr}^{123, *}_{12}c) \cup
\text{pr}^{123, *}_1 a) \\
& =
\text{pr}^{13}_{3, *}(
\gamma(\text{pr}^{123}_{13, *}(
\text{pr}^{123, *}_{23}c' \cdot \text{pr}^{123, *}_{12}c)) \cup
\text{pr}^{13, *}_1 a) \\
& =
G(c' \circ c)(a)
\end{align*},
avec des notations évidentes. La première égalité résulte des définitions.
La deuxième résulte de l'égalité
$\text{pr}^{23, *}_2 \circ \text{pr}^{12}_{2, *} =
\text{pr}^{123}_{23, *} \circ \text{pr}^{123, *}_{12}$,
qui découle immédiatement de la description de l'image directe
par les projections donnée dans le lemme \ref{lemma-pr2star-classical}.
La troisième égalité résulte du lemme \ref{lemma-pushforward-classical}
et du fait que $H^*$ est un foncteur.
La quatrième égalité résulte de l'axiome (C)(a) et du fait que
l'application de Gysin coïncide avec l'image réciproque plate pour les morphismes plats
(Homologie de Chow, lemme \ref{chow-lemma-lci-gysin-flat}).
La cinquième utilise l'axiome (C)(c), ainsi que
le lemme \ref{lemma-pushforward-classical}, pour obtenir
$\text{pr}^{23}_{3, *} \circ \text{pr}^{123}_{23, *} =
\text{pr}^{13}_{3, *} \circ \text{pr}^{123}_{13, *}$.
La sixième utilise la formule de projection du
lemme \ref{lemma-pushforward-classical}, ainsi que
l'axiome (C)(b), pour obtenir $
\text{pr}^{123}_{13, *}
\gamma(\text{pr}^{123, *}_{23}c' \cdot \text{pr}^{123, *}_{12}c) =
\gamma(\text{pr}^{123}_{13, *}(
\text{pr}^{123, *}_{23}c' \cdot \text{pr}^{123, *}_{12}c))$.
Enfin, la dernière égalité est la définition.

\medskip\noindent
Pour achever la démonstration que $G$ est un foncteur,
il faut montrer qu'il préserve les identités. Autrement dit, si
$1 = [\Delta] \in \text{Corr}^0(X, X)$ est l'identité
dans la catégorie des correspondances (voir
le lemme \ref{lemma-category-correspondences} et sa démonstration),
il faut montrer que $G([\Delta]) = \text{id}$.
Cela résulte de la détermination de
$\gamma([\Delta])$ dans le lemme \ref{lemma-class-diagonal-classical}
et du lemme \ref{lemma-pr2star-classical}.
La construction de $G$ comme foncteur sur les
variétés projectives lisses et les correspondances de degré $0$ est ainsi achevée.

\medskip\noindent
Il résulte des axiomes (A)(c) et (A)(d) que
$G(\Spec(k)) = H^*(\Spec(k))$ est canoniquement isomorphe à $F$
comme $F$-algèbre.
L'axiome de K\"unneth (B) montre que notre foncteur est compatible avec les produits tensoriels.
Il est donc un foncteur de catégories monoïdales symétriques.

\medskip\noindent
Il reste à vérifier que l'image de $G(c_2)$ dans $G(\mathbf{P}^1)$
est un espace vectoriel gradué inversible sur $F$ (en particulier, nous ne savons pas encore
que $G$ se prolonge à $M_k$).
D'après l'axiome (A)(d), l'application $\int_{\mathbf{P}^1} : H^2(\mathbf{P}^1) \to F$
est un isomorphisme. D'après l'axiome (A)(b), on a $\dim_F H^0(\mathbf{P}^1) = 1$.
Le lemme \ref{lemma-square-diagonal-classical} et l'axiome (A)(c)
donnent $2 - \dim_F H^1(\mathbf{P}^1) = c_1(T_{\mathbf{P}^1}) = 2$.
Ainsi, $H^1(\mathbf{P}^1) = 0$. Par conséquent,
$$
G(\mathbf{P}^1) = H^0(\mathbf{P}^1) \oplus H^2(\mathbf{P}^1)
$$
Rappelons que $1 = c_0 + c_2$ est une décomposition de l'identité
en somme d'idempotents orthogonaux dans
$\text{Corr}^0(\mathbf{P}^1, \mathbf{P}^1)$ ; voir
l'exemple \ref{example-decompose-P1}. On a $c_0 = a \circ b$, où
$a \in \text{Corr}^0(\Spec(k), \mathbf{P}^1)$ et
$b \in \text{Corr}^0(\mathbf{P}^1, \Spec(k))$, et où
$b \circ a = 1$ dans $\text{Corr}^0(\Spec(k), \Spec(k))$ ; voir la démonstration du
lemme \ref{lemma-inverse-h2}. Puisque $F = G(\Spec(k))$, la
fonctorialité entraîne que $G(c_0)$ est le projecteur sur le facteur direct
$H^0(\mathbf{P}^1) \subset G(\mathbf{P}^1)$. Par conséquent,
$G(c_2)$ est nécessairement la projection sur $H^2(\mathbf{P}^1)$,
ce qui achève la démonstration.
\end{proof}

\begin{proposition}
\label{proposition-weil-cohomology-theory-classical}
Soit $k$ un corps algébriquement clos et soit $F$ un corps de
caractéristique $0$. Une théorie cohomologique de Weil classique revient à se donner
un foncteur $\mathbf{Q}$-linéaire
$$
G : M_k \longrightarrow \text{espaces vectoriels gradués sur }F\text{}
$$
de catégories monoïdales symétriques, muni d'un isomorphisme
$F[2] \to G(\mathbf{1}(1))$ d'espaces vectoriels gradués sur $F$, tel que,
de plus,
\begin{enumerate}
\item $G(h(X))$ soit concentré en degrés positifs ou nuls ;
\item $\dim_F G^0(h(X)) = 1$
\end{enumerate}
pour toute variété projective lisse $X$.
\end{proposition}

\begin{proof}
Étant donnés $G$ et $F[2] \to G(\mathbf{1}(1))$, on pose $H^*(X) = G(h(X))$
et l'on obtient des données (D1), (D2) et (D3) satisfaisant à tous les axiomes (A), (B) et (C),
sauf peut-être (A)(c) et (A)(d) ; voir
le lemme \ref{lemma-from-functor-to-weil-classical} et sa démonstration.
Observons que les hypothèses (1) et (2) entraînent les axiomes (A)(c) et (A)(d),
compte tenu des axiomes (A)(a) et (A)(b) déjà établis.

\medskip\noindent
Réciproquement, à partir de $H^*$, on obtient un foncteur $G$ par la construction du
lemme \ref{lemma-from-weil-to-functor-classical}.
Soient $X = \mathbf{P}^1, c_0, c_2$ comme dans l'exemple \ref{example-decompose-P1}.
Nous avons construit un isomorphisme $1(-1) \to (X, c_2, 0)$ de motifs au
lemme \ref{lemma-inverse-h2}. Dans la démonstration du
lemme \ref{lemma-from-weil-to-functor-classical}, nous avons vu que
$G(1(-1)) = G(X, c_2, 0) = H^2(\mathbf{P}^1)[-2]$.
Ainsi, l'isomorphisme $\int_{\mathbf{P}^1} : H^2(\mathbf{P}^1) \to F$
de l'axiome (A)(d) donne un isomorphisme $G(1(-1)) \to F[-2]$,
qui détermine un isomorphisme $F[2] \to G(\mathbf{1}(1))$.
Enfin, puisque $G(h(X)) = H^*(X)$, les hypothèses (1) et (2)
résultent de l'axiome (A).
\end{proof}











\section{Cycles sur les corps non clos}
\label{section-cycles-nonclosed}

\noindent
Voici quelques lemmes qui nous seront utiles pour étudier les motifs
sur des corps de base non algébriquement clos.

\begin{lemma}
\label{lemma-generated-by-separable}
Soit $k$ un corps et soit $X$ un schéma projectif lisse sur $k$.
Alors $\CH_0(X)$ est engendré par les classes de points fermés dont les corps
résiduels sont séparables sur $k$.
\end{lemma}

\begin{proof}
Le lemme est immédiat si $k$ est de caractéristique $0$ ou s'il est parfait.
On peut donc supposer que $k$ est un corps infini de caractéristique $p > 0$.

\medskip\noindent
On peut supposer $X$ irréductible de dimension $d$.
Alors $k' = H^0(X, \mathcal{O}_X)$ est une extension finie séparable
de $k$ et $X$ est géométriquement intègre sur $k'$.
Voir Variétés, lemmes \ref{varieties-lemma-smooth-geometrically-normal},
\ref{varieties-lemma-proper-geometrically-reduced-global-sections} et
\ref{varieties-lemma-baby-stein}. Remplaçons donc $k$ par $k'$
et supposons désormais $X$ géométriquement intègre.

\medskip\noindent
Soit $x \in X$ un point fermé. Pour démontrer le lemme, nous allons montrer que
$[x] \in \CH_0(X)$ est rationnellement équivalent à une combinaison linéaire
à coefficients entiers de classes de points fermés dont les corps résiduels
sont séparables sur $k$. Choisissons un
$\mathcal{O}_X$-module inversible ample $\mathcal{L}$. Posons
$$
V = \{s \in H^0(X, \mathcal{L}) \mid s(x) = 0 \}
$$
En remplaçant $\mathcal{L}$ par une puissance, on peut supposer
que (a) $\mathcal{L}$ est très ample, (b) $V$ engendre
$\mathcal{L}$ sur $X \setminus x$, (c) le morphisme
$X \setminus x \to \mathbf{P}(V)$ est une immersion et (d)
l'application $V \to \mathfrak m_x\mathcal{L}_x/\mathfrak m_x^2\mathcal{L}_x$
est surjective ; voir
Morphismes, lemme \ref{morphisms-lemma-finite-type-ample-very-ample},
Variétés, lemme \ref{varieties-lemma-generate-over-complement}, et
Propriétés, proposition \ref{properties-proposition-characterize-ample}.
Considérons l'ensemble
$$
V^d \supset U =
\{
(s_1, \ldots, s_d) \in V^d \mid s_1, \ldots, s_d
\text{ engendrent }
\mathfrak m_x\mathcal{L}_x/\mathfrak m_x^2\mathcal{L}_x
\text{ sur }\kappa(x)
\}
$$
Puisque $\mathcal{O}_{X, x}$ est un anneau local régulier de dimension $d$,
on a $\dim_{\kappa(x)}(\mathfrak m_x/\mathfrak m_x^2) = d$ ;
ainsi, $U$ est un ouvert de Zariski non vide de $V^d$.
Pour $(s_1, \ldots, s_d) \in U$, posons $H_i = Z(s_i)$. Puisque
$s_1, \ldots, s_d$ engendrent $\mathfrak m_x\mathcal{L}_x$,
on a
$$
H_1 \cap \ldots \cap H_d = x \amalg Z
$$
au sens schématique, pour un certain sous-schéma fermé $Z \subset X$.
D'après Bertini (sous la forme de Variétés, lemme \ref{varieties-lemma-bertini}),
pour un élément général $s_1 \in V$, le schéma $H_1 \cap (X \setminus x)$
est lisse sur $k$, de dimension $d - 1$.
Une fois $s_1$ choisi, pour un élément général
$s_2 \in V$, le schéma $H_1 \cap H_2 \cap (X \setminus x)$
est lisse sur $k$, de dimension $d - 2$. Et ainsi de suite.
Il en résulte que, pour un choix suffisamment général de
$(s_1, \ldots, s_d) \in U$, le schéma $Z$ est \'etale sur $\Spec(k)$.
En particulier, $H_1 \cap \ldots \cap H_d$ est de dimension $0$,
d'où
$$
[H_1] \cdot \ldots \cdot [H_d] = [x] + [Z]
$$
dans $\CH_0(X)$, par applications successives de
Homologie de Chow, lemme \ref{chow-lemma-intersect-properly} (nous omettons les détails).
Cela achève la démonstration, car on obtient $[x] \sim_{rat} - [Z] + [Z']$,
où $Z' = H'_1 \cap \ldots \cap H'_d$ est une intersection complète générale
de lieux d'annulation de sections suffisamment générales
de $\mathcal{L}$, qui est \'etale sur $k$ par le même
raisonnement que précédemment.
\end{proof}

\begin{lemma}
\label{lemma-chow-limit}
Soit $K/k$ une extension algébrique de corps et soit $X$ un schéma
de type fini sur $k$. Alors $\CH_i(X_K) = \colim \CH_i(X_{k'})$, où la
limite inductive porte sur les sous-extensions $K/k'/k$ telles que $k'/k$ soit finie.
\end{lemma}

\begin{proof}
C'est un cas particulier de
Homologie de Chow, lemme \ref{chow-lemma-chow-limit}.
\end{proof}

\begin{lemma}
\label{lemma-divide-difference-points}
Soit $k$ un corps et soit $X$ un schéma projectif lisse
géométriquement irréductible sur $k$. Soient $x, x' \in X$ des points $k$-rationnels.
Soit $n$ un entier inversible dans $k$.
Il existe alors une extension finie séparable $k'/k$ telle que
l'image réciproque de $[x] - [x']$ sur $X_{k'}$
soit divisible par $n$ dans $\CH_0(X_{k'})$.
\end{lemma}

\begin{proof}
Soit $k'$ une clôture séparable de $k$. Supposons que l'on sache montrer
que l'image réciproque de $[x] - [x']$ sur $X_{k'}$ est divisible par $n$ dans
$\CH_0(X_{k'})$. On conclut alors par le lemme \ref{lemma-chow-limit}.
On peut donc supposer désormais $k$ séparablement clos.

\medskip\noindent
Supposons $\dim(X) > 1$. Soit $\mathcal{L}$ un faisceau inversible ample sur $X$.
Posons
$$
V = \{s \in H^0(X, \mathcal{L}) \mid s(x) = 0\text{ et }s(x') = 0 \}
$$
En remplaçant $\mathcal{L}$ par une puissance, on voit que, pour
un élément général $v \in V$, le diviseur correspondant $H_v \subset X$ est lisse
en dehors de $x$ et de $x'$ ; voir
Variétés, lemmes \ref{varieties-lemma-generate-over-complement} et
\ref{varieties-lemma-bertini}. Pour trouver $v$, on utilise le fait que $k$ est infini,
car séparablement clos.
Si l'on choisit $s$ général, l'image de $s$ dans
$\mathfrak m_x\mathcal{L}_x/\mathfrak m_x^2\mathcal{L}_x$
est non nulle, ce qui entraîne que $H_v$ est lisse en $x$
(nous omettons les détails). Il en est de même pour $x'$. Ainsi, $H_v$ est lisse.
D'après Variétés, lemme \ref{varieties-lemma-connectedness-ample-divisor}
(appliqué après changement de base de tous les objets
à la clôture algébrique de $k$),
$H_v$ est géométriquement connexe.
Il suffit de démontrer le résultat pour
$[x] - [x']$ considéré comme un élément de $\CH_0(H_v)$.
On se ramène ainsi au cas d'une courbe.

\medskip\noindent
Supposons que $X$ soit une courbe. Alors $\mathcal{O}_X(x - x')$
définit un point $k$-rationnel $g$ de $J = \underline{\Pic}^0_{X/k}$ ; voir
Schémas de Picard des courbes, lemme \ref{pic-lemma-picard-pieces}.
Rappelons que $J$ est une variété propre et lisse sur $k$,
qui est aussi un schéma en groupes sur $k$ (même référence).
Par conséquent, $J$ est géométriquement intègre
(voir Variétés, lemmes \ref{varieties-lemma-geometrically-connected-criterion}
et \ref{varieties-lemma-smooth-geometrically-normal}).
Autrement dit, $J$ est une variété abélienne ; voir
Groupoïdes, définition \ref{groupoids-definition-abelian-variety}.
Ainsi, $[n] : J \to J$ est fini \'etale d'après
Groupoïdes, proposition \ref{groupoids-proposition-review-abelian-varieties}
(c'est ici que l'on utilise l'inversibilité de $n$ dans $k$).
Puisque $k$ est séparablement clos, on en déduit $g = [n](g')$
pour un certain $g' \in J(k)$. Si $\mathcal{L}$ est le module inversible de degré $0$
sur $X$ correspondant à $g'$, on en déduit
$\mathcal{O}_X(x - x') \cong \mathcal{L}^{\otimes n}$, comme voulu.
\end{proof}

\begin{lemma}
\label{lemma-kernel-to-closure}
Soit $K/k$ une extension algébrique de corps.
Soit $X$ un schéma de type fini sur $k$.
Le noyau de l'application $\CH_i(X) \to \CH_i(X_K)$
construite au lemme \ref{lemma-chow-limit}
est de torsion.
\end{lemma}

\begin{proof}
Il suffit manifestement de montrer que le noyau
de l'image réciproque plate $\CH_i(X) \to \CH_i(X_{k'})$
par $\pi : X_{k'} \to X$ est de torsion,
pour toute extension finie $k'/k$. Cela résulte de
$\pi_* \pi^* \alpha = [k' : k] \alpha$, d'après
Homologie de Chow, lemme \ref{chow-lemma-finite-flat}.
\end{proof}

\begin{lemma}[Voevodsky]
\label{lemma-smash-nilpotence}
\begin{reference}
\cite{nilpotence}
\end{reference}
Soit $k$ un corps et soit $X$ un schéma projectif lisse
géométriquement irréductible sur $k$. Soient $x, x' \in X$ des points $k$-rationnels.
Pour $n$ assez grand, la classe du zéro-cycle
$$
([x] - [x']) \times \ldots \times ([x] - [x']) \in
\CH_0(X^n)
$$
est de torsion.
\end{lemma}

\begin{proof}
Si l'on sait le montrer après changement de base à la clôture algébrique de $k$,
le résultat s'ensuit sur $k$, car le noyau de l'image réciproque
est de torsion d'après le lemme \ref{lemma-kernel-to-closure}.
On peut donc supposer désormais $k$ algébriquement clos.

\medskip\noindent
Le théorème de Bertini permet de choisir une courbe lisse $C \subset X$ passant par
$x$ et $x'$. Voir la démonstration du lemme \ref{lemma-divide-difference-points}.
On peut donc supposer que $X$ est une courbe.

\medskip\noindent
Supposons que $X$ soit une courbe et que $k$ soit algébriquement clos.
Écrivons $S^n(X) = \underline{\Hilbfunctor}^n_{X/k}$, avec les notations de
Schémas de Picard des courbes, sections \ref{pic-section-hilbert-scheme-points}
et \ref{pic-section-divisors}. Il existe un morphisme canonique
$$
\pi : X^n \longrightarrow S^n(X)
$$
qui envoie le point $k$-rationnel $(x_1, \ldots, x_n)$ sur le point $k$-rationnel
correspondant au diviseur $[x_1] + \ldots + [x_n]$ sur $X$.
Le groupe symétrique $S_n$ agit fidèlement sur $X^n$.
Le morphisme $\pi$ est $S_n$-invariant, et les fibres de $\pi$ sont
des orbites sous $S_n$ (ensemblistement). Enfin, $\pi$ est fini plat de
degré $n!$ ; voir Schémas de Picard des courbes, lemme
\ref{pic-lemma-universal-object}.

\medskip\noindent
Soit $\alpha_n$ le zéro-cycle sur $X^n$ donné par la formule de
l'énoncé du lemme. Posons $\mathcal{L} = \mathcal{O}_X(x - x')$. Alors
$c_1(\mathcal{L}) \cap [X] = [x] - [x']$. Ainsi,
$$
\alpha_n = c_1(\mathcal{L}_1) \cap \ldots \cap c_1(\mathcal{L}_n) \cap [X^n]
$$,
où $\mathcal{L}_i = \text{pr}_i^*\mathcal{L}$ et où $\text{pr}_i : X^n \to X$
est la $i$-ième projection. D'après
Diviseurs, lemme \ref{divisors-lemma-finite-locally-free-has-norm}, ou
Diviseurs, lemme \ref{divisors-lemma-norm-in-normal-case},
il existe une norme pour $\pi$. Posons
$\mathcal{N} = \text{Norm}_\pi(\mathcal{L}_1)$ ;
voir Diviseurs, lemme \ref{divisors-lemma-norm-invertible}. Un calcul donne
$$
\pi^*\mathcal{N} =
(\mathcal{L}_1 \otimes \ldots \otimes \mathcal{L}_n)^{\otimes (n - 1)!}
$$
dans $\Pic(X^n)$. Nous omettons les détails ; indiquons que cela résulte
du fait que
$\text{Norm}_\pi : \pi_*\mathcal{O}_{X^n} \to \mathcal{O}_{S^n(X)}$,
composé avec l'application naturelle $\pi_*\mathcal{O}_{S^n(X)} \to \mathcal{O}_{X^n}$,
est égal au produit, pour tous les $\sigma \in S_n$, des actions de
$\sigma$ sur $\pi_*\mathcal{O}_{X^n}$. Considérons
$$
\beta_n = c_1(\mathcal{N})^n \cap [S^n(X)]
$$
dans $\CH_0(S^n(X))$. Observons que
$c_1(\mathcal{L}_i) \cap c_1(\mathcal{L}_i) = 0$,
car $\mathcal{L}_i$ provient d'une courbe par image réciproque ; voir
Homologie de Chow, lemme \ref{chow-lemma-vanish-above-dimension}. On obtient donc
\begin{align*}
\pi^*\beta_n
& =
((n - 1)!)^n
(\sum\nolimits_{i = 1, \ldots, n} c_1(\mathcal{L}_i))^n \cap [X^n] \\
& =
((n - 1)!)^n n^n 
c_1(\mathcal{L}_1) \cap \ldots \cap c_1(\mathcal{L}_n) \cap [X^n] \\
& =
(n!)^n \alpha_n
\end{align*}
Il suffit ainsi de montrer que $\beta_n$ est de torsion.

\medskip\noindent
Il existe un morphisme canonique
$$
f : S^n(X) \longrightarrow \underline{\Picardfunctor}^n_{X/k}
$$
Voir Schémas de Picard des courbes, lemme \ref{pic-lemma-picard-pieces}.
Pour $n \geq 2g - 1$, ce morphisme est un fibré projectif
(nous omettons les détails ; comparer avec la
démonstration de Schémas de Picard des courbes, lemme \ref{pic-lemma-picard-pieces}).
Le faisceau inversible $\mathcal{N}$ est trivial sur les
fibres de $f$, comme on le verra ci-dessous. Par la formule du fibré projectif
(Homologie de Chow, lemme \ref{chow-lemma-chow-ring-projective-bundle}),
on a donc $\mathcal{N} = f^*\mathcal{M}$ pour un certain module
inversible $\mathcal{M}$ sur $\underline{\Picardfunctor}^n_{X/k}$.
Il s'ensuit évidemment que
$$
c_1(\mathcal{N})^n = f^*(c_1(\mathcal{M})^n)
$$
est nul, puisque $n > g = \dim(\underline{\Picardfunctor}^n_{X/k})$
et que l'on peut appliquer Homologie de Chow, lemme \ref{chow-lemma-vanish-above-dimension},
comme précédemment.

\medskip\noindent
Il reste à montrer que $\mathcal{N}$ est trivial sur une fibre $F$
de $f$. Puisque les fibres de $f$ sont des espaces projectifs et que
$\Pic(\mathbf{P}^m_k) = \mathbf{Z}$
(Diviseurs, lemme \ref{divisors-lemma-Pic-projective-space-UFD}),
on pourrait le montrer en calculant le degré de $\mathcal{N}$
sur une droite contenue dans la fibre. Nous le démontrerons plutôt
en établissant que $\mathcal{N}$ est algébriquement
équivalent à zéro. Montrons d'abord qu'il existe un schéma connexe $T$ de type fini
sur $k$, un module inversible $\mathcal{L}'$ sur $T \times X$ et des points
$k$-rationnels $p, q \in T$ tels que
$\mathcal{M}_p \cong \mathcal{O}_X$ et $\mathcal{M}_q = \mathcal{L}$.
En effet, puisque $\mathcal{L} = \mathcal{O}_X(x - x')$, on peut prendre
$T = X$, $p = x'$, $q = x$ et
$\mathcal{L}' = \mathcal{O}_{X \times X}(\Delta)
\otimes \text{pr}_2^*\mathcal{O}_X(-x')$.
On note ensuite $\mathcal{L}'_i$, sur
$T \times X^n$, pour $i = 1, \ldots, n$,
l'image réciproque de $\mathcal{L}'$ par
$\text{id}_T \times \text{pr}_i : T \times X^n \to T \times X$.
Enfin, on pose
$\mathcal{N}' = \text{Norm}_{\text{id}_T \times \pi}(\mathcal{L}'_1)$
sur $T \times S^n(X)$.
Par construction, on a $\mathcal{N}'_p = \mathcal{O}_{S^n(X)}$
et $\mathcal{N}'_q = \mathcal{N}$.
Il en résulte que
$$
\mathcal{N}'|_{T \times F}
$$
est un module inversible sur $T \times F \cong T \times \mathbf{P}^m_k$,
dont la fibre en $p$ est le module inversible trivial et dont la fibre
en $q$ est $\mathcal{N}|_F$. Puisque la caractéristique d'Euler
du fibré trivial est $1$ et que cette caractéristique d'Euler
est localement constante en famille (Catégories dérivées de schémas,
lemme \ref{perfect-lemma-chi-locally-constant-geometric}),
on obtient $\chi(F, \mathcal{N}^{\otimes s}|_F) = 1$
pour tout $s \in \mathbf{Z}$. Cela n'est possible que si
$\mathcal{N}|_F \cong \mathcal{O}_F$ (voir
Cohomologie des schémas, lemme
\ref{coherent-lemma-cohomology-projective-space-over-ring}),
ce qui achève la démonstration. Nous omettons certains détails.
\end{proof}















\section{Théories cohomologiques de Weil, I}
\label{section-axioms}

\noindent
Cette section est l'analogue de la section \ref{section-axioms-classical}
sur un corps quelconque. Autrement dit, nous déterminons les données
et les axiomes qui correspondent aux foncteurs $G$ de catégories monoïdales symétriques
de la catégorie des motifs dans celle des espaces vectoriels gradués, tels que
$G(\mathbf{1}(1))$ soit placé en degré $-2$. À la section \ref{section-old},
nous définirons une théorie cohomologique de Weil en ajoutant une seule
condition supplémentaire.

\medskip\noindent
Fixons un corps $k$ (le corps de base).
Fixons un corps $F$ de caractéristique $0$ (le corps de coefficients).
Les données sont les suivantes :
\begin{enumerate}
\item[(D0)] Un espace vectoriel de dimension $1$ sur $F$, noté $F(1)$.
\item[(D1)] Un foncteur contravariant $H^*$ de la catégorie
des schémas projectifs lisses sur $k$ dans la catégorie des
$F$-algèbres graduées commutatives.
\item[(D2)] Pour tout schéma projectif lisse $X$ sur $k$,
un homomorphisme de groupes $\gamma : \CH^i(X) \to H^{2i}(X)(i)$.
\item[(D3)] Pour tout schéma projectif lisse non vide $X$ sur $k$,
équidimensionnel de dimension $d$, une application
$\int_X : H^{2d}(X)(d) \to F$.
\end{enumerate}
Précisons le sens de ces données par quelques remarques
et introduisons la terminologie correspondante.

\medskip\noindent
Remarques sur (D0).
L'espace vectoriel $F(1)$ définit des {\it torsions de Tate} sur la catégorie
des espaces vectoriels sur $F$. En effet, pour $n \in \mathbf{Z}$, on pose
$F(n) = F(1)^{\otimes n}$ si $n \geq 0$, puis $F(-1) = \Hom_F(F(1), F)$,
et $F(n) = F(-1)^{\otimes - n}$ si $n < 0$. Comparer avec
Compléments d'algèbre, section \ref{more-algebra-section-picard}.
Pour un espace vectoriel sur $F$, noté $V$, on définit la {\it $n$-ième torsion de Tate}
$$
V(n) = V \otimes_F F(n)
$$
Nous emploierons des notations évidentes : par exemple, pour des espaces vectoriels sur $F$, notés $U$, $V$
et $W$, et une application linéaire $U \otimes_F V \to W$, on obtient une application
linéaire $U(n) \otimes_F V(m) \to W(n + m)$ pour $n, m \in \mathbf{Z}$.

\medskip\noindent
Remarques sur (D1).
Pour un schéma projectif lisse $X$ sur $k$, on appelle $H^*(X)$
la {\it cohomologie} de $X$. Pour un morphisme $f : X \to Y$
de schémas projectifs lisses sur $k$, on note $f^* : H^*(Y) \to H^*(X)$
l'application $H^*(f)$ et on l'appelle l'{\it application image réciproque}.

\medskip\noindent
Remarques sur (D2). L'application $\gamma$ est appelée l'{\it application classe de cycle}.
On appelle $\gamma(\alpha)$ la {\it classe de cohomologie} de $\alpha$.
Si $Z \subset Y \subset X$ sont des sous-schémas fermés, avec $Y$ et $X$
projectifs lisses sur $k$ et $Z$ intègre, alors $[Z]$ peut
désigner la classe du cycle $[Z]$ dans $\CH^*(Y)$ ou dans $\CH^*(X)$.
Dans ce cas, la notation $\gamma([Z])$ est ambiguë, et son sens
doit être déduit du contexte.

\medskip\noindent
Remarques sur (D3). L'application $\int_X$ est parfois appelée
l'{\it application trace} et est parfois notée $\text{Tr}_X$.

\medskip\noindent
Le premier axiome est souvent appelé {\it dualité de Poincar\'e} :
\begin{enumerate}
\item[(A)] Soit $X$ un schéma projectif lisse non vide sur $k$,
équidimensionnel de dimension $d$. Alors :
\begin{enumerate}
\item $\dim_F H^i(X) < \infty$ pour tout $i$ ;
\item $H^i(X) \times H^{2d - i}(X)(d) \rightarrow
H^{2d}(X)(d) \rightarrow F$
est un accouplement parfait pour tout $i$, la dernière
application étant l'application trace $\int_X$.
\end{enumerate}
\end{enumerate}
Soit $f : X \to Y$ un morphisme de schémas projectifs lisses non vides, avec $X$
équidimensionnel de dimension $d$ et $Y$ équidimensionnel de dimension $e$.
À l'aide de la dualité de Poincar\'e, on peut définir une {\it image directe}
$$
f_* : H^{2d - i}(X)(d) \longrightarrow H^{2e - i}(Y)(e)
$$
comme l'application contragrédiente de l'application linéaire $f^* : H^i(Y) \to H^i(X)$.
Explicitement, pour $a \in H^{2d - i}(X)(d)$, l'élément $f_*a \in H^{2e - i}(Y)(e)$
est caractérisé par
$$
\int_X f^*b \cup a = \int_Y b \cup f_*a
$$
pour tout $b \in H^i(Y)$.

\begin{lemma}
\label{lemma-pushforward}
Supposons que les données (D0), (D1) et (D3) satisfassent à (A). Pour $f : X \to Y$
un morphisme de schémas projectifs lisses non vides et équidimensionnels sur $k$,
on a $f_*(f^*b \cup a) = b \cup f_*a$. Si $g : Y \to Z$ est un second morphisme,
avec $Z$ non vide, projectif lisse et équidimensionnel, alors
$g_* \circ f_* = (g \circ f)_*$.
\end{lemma}

\begin{proof}
La première égalité résulte de
$$
\int_Y c \cup b \cup f_*a =
\int_X f^*c \cup f^*b \cup a =
\int_Y c \cup f_*(f^*b \cup a).
$$
La seconde résulte de
$$
\int_Z c \cup (g \circ f)_*a = \int_X (g \circ f)^*c \cup a =
\int_X f^* g^* c \cup a = \int_Y g^*c \cup f_*a = \int_Z c \cup g_*f_*a
$$
Cela achève la démonstration.
\end{proof}

\noindent
Le deuxième axiome exprime que $H^*$ respecte la structure monoïdale
donnée par les produits, au moyen de la {\it formule de K\"unneth} :
\begin{enumerate}
\item[(B)] Soient $X$ et $Y$ des schémas projectifs lisses sur $k$.
\begin{enumerate}
\item L'application $H^*(X) \otimes_F H^*(Y) \to H^*(X \times Y)$,
$\alpha \otimes \beta \mapsto \text{pr}_1^*\alpha \cup \text{pr}_2^*\beta$,
est un isomorphisme ;
\item si $X$ et $Y$ sont non vides et équidimensionnels, alors
$\int_{X \times Y} = \int_X \otimes \int_Y$ par l'identification (a).
\end{enumerate}
\end{enumerate}
L'axiome (B)(b) permet de calculer les images directes par les projections.

\begin{lemma}
\label{lemma-pr2star}
Supposons que les données (D0), (D1) et (D3) satisfassent à (A) et (B).
Soient $X$ et $Y$ des schémas projectifs lisses non vides sur $k$,
équidimensionnels de dimensions $d$ et $e$. Alors
$\text{pr}_{2, *} : H^*(X \times Y)(d + e) \to H^*(Y)(e)$ envoie
$a \otimes b$ sur $(\int_X a) b$.
\end{lemma}

\begin{proof}
Cela résulte des axiomes (B)(a) et (B)(b).
\end{proof}

\noindent
Le troisième axiome porte sur les applications classes de cycles :
\begin{enumerate}
\item[(C)] Les applications classes de cycles satisfont aux règles suivantes :
\begin{enumerate}
\item pour un morphisme $f : X \to Y$ de schémas projectifs lisses sur
$k$, on a $\gamma(f^!\beta) = f^*\gamma(\beta)$ pour $\beta \in \CH^*(Y)$ ;
\item pour un morphisme $f : X \to Y$ de schémas projectifs lisses non vides
et équidimensionnels sur $k$, on a
$\gamma(f_*\alpha) = f_*\gamma(\alpha)$ pour $\alpha \in \CH^*(X)$ ;
\item pour tout schéma projectif lisse $X$ sur $k$, on a
$\gamma(\alpha \cdot \beta) = \gamma(\alpha) \cup \gamma(\beta)$
pour $\alpha, \beta \in \CH^*(X)$ ;
\item $\int_{\Spec(k)} \gamma([\Spec(k)]) = 1$.
\end{enumerate}
\end{enumerate}
Explicitons l'axiome (C)(b). Soit $f : X \to Y$
comme dans (C)(b), avec $\dim(X) = d$ et $\dim(Y) = e$. L'image directe
sur les groupes de Chow donne alors
$$
f_* : \CH^{d - i}(X) = \CH_i(X) \to \CH_i(Y) = \CH^{e - i}(Y)
$$
Soit $\alpha \in \CH^{d - i}(X)$. D'une part, on a
$f_*\alpha \in \CH^{e - i}(Y)$, d'où
$\gamma(f_*\alpha) \in H^{2e - 2i}(Y)(e - i)$.
D'autre part, on a
$\gamma(\alpha) \in H^{2d - 2i}(X)(d - i)$, d'où
$f_*\gamma(\alpha) \in H^{2e - 2i}(Y)(e - i)$ également.
La condition $\gamma(f_*\alpha) = f_*\gamma(\alpha)$ a donc bien un sens.

\begin{remark}
\label{remark-replace-cup-product}
Supposons que les données (D0), (D1), (D2) et (D3) satisfassent à (A), (B) et (C)(a).
Soit $X$ un schéma projectif lisse sur $k$. On obtient des applications
$$
H^*(X) \otimes_F H^*(X) \longrightarrow H^*(X \times X)
\xrightarrow{\Delta^*} H^*(X)
$$,
où la première flèche est celle de l'axiome (B) et $\Delta^*$
est l'image réciproque par le morphisme diagonal $\Delta : X \to X \times X$.
La composée est le cup-produit, puisque l'image réciproque est un homomorphisme d'algèbres
et que $\text{pr}_i \circ \Delta = \text{id}$.
D'autre part, pour des cycles $\alpha, \beta$ sur $X$,
le produit d'intersection est défini par la formule
$$
\alpha \cdot \beta =
\Delta^!(\alpha \times \beta)
$$
Autrement dit, $\alpha \cdot \beta$ est l'image réciproque du
produit extérieur $\alpha \times \beta$ sur $X \times X$ par
la diagonale. Notons aussi que
$\alpha \times \beta = \text{pr}_1^*\alpha \cdot \text{pr}_2^*\beta$
dans $\CH^*(X \times X)$ (nous omettons la démonstration). Ainsi, sous l'axiome (C)(a),
l'axiome (C)(c) équivaut à la compatibilité de $\gamma$
avec le produit extérieur, au sens où
$\gamma(\alpha \times \beta)$ est égal à
$\text{pr}_1^*\gamma(\alpha) \cup \text{pr}_2^*\gamma(\beta)$.
\end{remark}

\begin{lemma}
\label{lemma-base}
Supposons que les données (D0), (D1), (D2) et (D3) satisfassent à (A), (B) et (C).
Alors $H^i(\Spec(k)) = 0$ pour $i \not = 0$, et il existe un
unique isomorphisme de $F$-algèbres $F = H^0(\Spec(k))$.
On a $\gamma([\Spec(k)]) = 1$ et $\int_{\Spec(k)} 1 = 1$.
\end{lemma}

\begin{proof}
D'après l'axiome (C)(d), $H^0(\Spec(k))$ est non nul et même
$\gamma([\Spec(k)])$ est non nul.
Puisque $\Spec(k) \times \Spec(k) = \Spec(k)$, on obtient
$$
H^*(\Spec(k)) \otimes_F H^*(\Spec(k)) = H^*(\Spec(k))
$$
par l'axiome (B)(a), ce qui entraîne, en considérant les dimensions, que seul
$H^0$ est non nul et qu'il est de dimension $1$. Ainsi,
$F = H^0(\Spec(k))$ par l'unique isomorphisme de $F$-algèbres
envoyant $1 \in F$ sur $1 \in H^0(\Spec(k))$.
Puisque $[\Spec(k)] \cdot [\Spec(k)] = [\Spec(k)]$ dans
l'anneau de Chow de $\Spec(k)$, on obtient
$\gamma([\Spec(k)) \cup \gamma([\Spec(k)]) = \gamma([\Spec(k)])$
par l'axiome (C)(c). Comme nous savons déjà que $\gamma([\Spec(k)])$ est non nul,
il est nécessairement égal à $1$.
Enfin, on a $\int_{\Spec(k)} 1 = 1$, puisque
$\int_{\Spec(k)} \gamma([\Spec(k)]) = 1$
d'après l'axiome (C)(d).
\end{proof}

\begin{lemma}
\label{lemma-unit}
Supposons que les données (D0), (D1), (D2) et (D3) satisfassent à (A), (B) et (C).
Soit $X$ un schéma projectif lisse sur $k$.
Si $X = \emptyset$, alors $H^*(X) = 0$.
Si $X$ est non vide, alors $\gamma([X]) = 1$ et $1 \not = 0$ dans $H^0(X)$.
\end{lemma}

\begin{proof}
Supposons d'abord $X$ non vide.
Observons que $[X]$ est l'image réciproque de $[\Spec(k)]$ par le morphisme structural
$p : X \to \Spec(k)$. On obtient donc $\gamma([X]) = 1$ par l'axiome (C)(a)
et le lemme \ref{lemma-base}. Soit $X' \subset X$ une composante irréductible.
Par fonctorialité, il suffit de montrer que $1 \not = 0$ dans $H^0(X')$.
On peut donc supposer $X$ irréductible, et en particulier
non vide et équidimensionnel, de dimension $d$.
Pour établir $1 \not = 0$, il suffit de montrer que $H^*(X)$ est non nul.

\medskip\noindent
Soit $x \in X$ un point fermé dont le corps résiduel $k'$
est séparable sur $k$ ; voir
Variétés, lemme \ref{varieties-lemma-smooth-separable-closed-points-dense}.
Soit $i : \Spec(k') \to X$ le morphisme d'inclusion.
Notons $p : X \to \Spec(k)$ le morphisme structural.
Observons que
$p_*i_*[\Spec(k')] = [k' : k][\Spec(k)]$ dans $\CH_0(\Spec(k))$.
En appliquant deux fois l'axiome (C)(b), puis le lemme \ref{lemma-base},
on en déduit que
$$
p_*i_*\gamma([\Spec(k')]) = \gamma([k' : k][\Spec(k)]) = [k' : k]
\in F = H^0(\Spec(k))
$$
est non nul. Ainsi, $i_*\gamma([\Spec(k)]) \in H^{2d}(X)(d)$ est non nul,
car son image par $p_*$ est non nulle. Cela achève la démonstration
lorsque $X$ est non vide.

\medskip\noindent
Considérons enfin le cas du schéma vide. L'axiome (B)(a) donne
$H^*(\emptyset) \otimes H^*(\emptyset) = H^*(\emptyset)$,
d'où il résulte que $H^*(\emptyset)$ est soit nul, soit de dimension $1$
en degré $0$. L'axiome (B)(a) donne à nouveau
$H^*(\emptyset) \otimes H^*(X) = H^*(\emptyset)$ pour
tout schéma projectif lisse $X$ sur $k$. L'axiome (A)(b)
et la non-nullité de $H^0(X)$ établie ci-dessus
entraînent que $H^*(X)$ est non nul en au moins deux degrés
si $\dim(X) > 0$. Cela force alors $H^*(\emptyset)$ à être nul.
\end{proof}

\begin{lemma}
\label{lemma-push-unit}
Supposons que les données (D0), (D1), (D2) et (D3) satisfassent à (A), (B) et (C).
Soit $i : X \to Y$ une immersion fermée de schémas projectifs lisses
non vides et équidimensionnels sur $k$. Alors
$\gamma([X]) = i_*1$ dans $H^{2c}(Y)(c)$, où $c = \dim(Y) - \dim(X)$.
\end{lemma}

\begin{proof}
En effet, $1 = \gamma([X])$ dans $H^0(X)$ d'après le lemme \ref{lemma-unit},
et il suffit alors d'appliquer l'axiome (C)(b).
\end{proof}

\begin{lemma}
\label{lemma-class-diagonal}
Supposons que les données (D0), (D1), (D2) et (D3) satisfassent à (A), (B) et (C).
Soit $X$ un schéma projectif lisse non vide sur $k$, équidimensionnel
de dimension $d$. Choisissons une base $e_{i, j}, j = 1, \ldots, \beta_i$ de
$H^i(X)$ sur $F$. À l'aide de la formule de K\"unneth, écrivons
$$
\gamma([\Delta]) =
\sum\nolimits_i
\sum\nolimits_j e_{i, j} \otimes e'_{2d - i , j}
\quad\text{dans}\quad
\bigoplus\nolimits_i H^i(X) \otimes_F H^{2d - i}(X)(d)
$$
avec $e'_{2d - i, j} \in H^{2d - i}(X)(d)$.
Alors $\int_X e_{i, j} \cup e'_{2d - i, j'} = (-1)^i\delta_{jj'}$.
\end{lemma}

\begin{proof}
Rappelons que $\Delta^* : H^*(X \times X) \to H^*(X)$ est l'application
cup-produit $H^*(X) \otimes_F H^*(X) \to H^*(X)$ ; voir
la remarque \ref{remark-replace-cup-product}. D'autre part,
rappelons que $\gamma([\Delta]) = \Delta_*1$ (lemme \ref{lemma-push-unit}),
d'où
$$
\int_{X \times X} \gamma([\Delta]) \cup a \otimes b =
\int_{X \times X} \Delta_*1 \cup a \otimes b =
\int_X a \cup b
$$
d'après le lemme \ref{lemma-pushforward}.
D'autre part, on a
$$
\int_{X \times X} (\sum e_{i, j} \otimes e'_{2d -i , j}) \cup a \otimes b =
\sum (\int_X a \cup e_{i, j})(\int_X e'_{2d - i, j} \cup b)
$$
par l'axiome (B)(b) ; notons que nous avons effectué deux échanges, de sorte que
chaque terme a le signe $1$. Ainsi, si l'on choisit $a$ de sorte que
$\int_X a \cup e_{i, j} = 1$ et que tous les autres accouplements soient nuls,
on obtient $\int_X e'_{2d - i, j} \cup b = \int_X a \cup b$
pour tout $b$, c'est-à-dire $e'_{2d - i, j} = a$. Le lemme est démontré.
\end{proof}

\begin{lemma}
\label{lemma-cohomology-P1}
Supposons que les données (D0), (D1), (D2) et (D3) satisfassent à (A), (B) et (C).
Alors $H^*(\mathbf{P}^1_k)$ est de dimension $1$ en degrés $0$ et $2$
et nul en tout autre degré.
\end{lemma}

\begin{proof}
Soit $x \in \mathbf{P}^1_k$ un point $k$-rationnel. Observons que
$\Delta = \text{pr}_1^*x + \text{pr}_2^*x$ comme diviseurs sur
$\mathbf{P}^1_k \times \mathbf{P}^1_k$. L'axiome (C)(a)
et l'additivité de $\gamma$ donnent
$$
\gamma([\Delta]) =
\text{pr}_1^*\gamma([x]) +
\text{pr}_2^*\gamma([x]) =
\gamma([x]) \otimes 1 + 1 \otimes \gamma([x])
$$
dans $H^*(\mathbf{P}^1_k \times \mathbf{P}^1_k) =
H^*(\mathbf{P}^1_k) \otimes_F H^*(\mathbf{P}^1_k)$.
Or le lemme \ref{lemma-class-diagonal}
montre que $\gamma([\Delta])$ ne peut s'écrire
comme somme de moins de $\sum \beta_i$ tenseurs purs,
où $\beta_i = \dim_F H^i(\mathbf{P}^1_k)$.
Ainsi, $\sum \beta_i \leq 2$.
D'après le lemme \ref{lemma-unit}, on a $H^0(\mathbf{P}^1_k) \not = 0$.
La dualité de Poincar\'e, plus précisément l'axiome (A)(b),
donne $\beta_0 = \beta_2$. Le lemme s'ensuit.
\end{proof}

\begin{lemma}
\label{lemma-weil-additive}
Supposons que les données (D0), (D1), (D2) et (D3) satisfassent à (A), (B) et (C).
Si $X$ et $Y$ sont des schémas projectifs lisses sur $k$, alors
$H^*(X \amalg Y) \to H^*(X) \times H^*(Y)$,
$a \mapsto (i^*a, j^*a)$, est un isomorphisme, où $i$ et $j$
sont les coprojections.
\end{lemma}

\begin{proof}
Si $X$ ou $Y$ est vide, cela résulte de
$H^*(\emptyset) = 0$, d'après le lemme \ref{lemma-unit}.
On peut donc supposer $X$ et $Y$ tous deux non vides.

\medskip\noindent
Montrons d'abord que l'application est injective. Observons que
l'on peut trouver des morphismes $X' \to X$ et $Y' \to Y$
de schémas projectifs lisses, tels que $X'$ et $Y'$ soient
équidimensionnels de même dimension et que
$X' \to X$ et $Y' \to Y$ admettent chacun une section. En effet,
décomposons $X = \coprod X_d$ et $Y = \coprod Y_e$
en sous-schémas ouverts et fermés équidimensionnels de
dimensions $d$ et $e$. On prend alors
$X' = \coprod X_d \times \mathbf{P}^{n - d}$
et $Y' = \coprod Y_e \times \mathbf{P}^{n - e}$ pour un entier
$n$ assez grand. L'image réciproque par
$X' \amalg Y' \to X \amalg Y$ est donc injective,
car ce morphisme admet une section, et
il suffit de démontrer l'injectivité pour $X', Y'$,
ce que nous faisons au paragraphe suivant.

\medskip\noindent
Montrons que l'application est injective lorsque $X$ et $Y$ sont équidimensionnels
de même dimension $d$.
Observons que $[X \amalg Y] = [X] + [Y]$ dans $\CH^0(X \amalg Y)$
et que $[X]$ et $[Y]$ sont des idempotents orthogonaux dans $\CH^0(X \amalg Y)$.
Ainsi,
$$
1 = \gamma([X \amalg Y] = \gamma([X]) + \gamma([Y]) = i_*1 + j_*1
$$
est une décomposition en idempotents orthogonaux. Nous avons utilisé ici
les lemmes \ref{lemma-unit} et \ref{lemma-push-unit}, ainsi que l'axiome (C)(c).
On obtient alors
$$
a = a \cup 1 = a \cup i_*1 + a \cup j_*1 =
i_*(i^*a) + j_*(j^*a)
$$
par la formule de projection (lemme \ref{lemma-pushforward}), d'où
l'injectivité.

\medskip\noindent
Montrons que l'application est surjective. Posons $e = \gamma([X])$ et $f = \gamma([Y])$,
considérés comme des éléments de $H^0(X \amalg Y)$. On a
$i^*e = 1$, $i^*f = 0$, $j^*e = 0$ et $j^*f = 1$ d'après l'axiome (C)(a).
Ainsi, si $i^* : H^*(X \amalg Y) \to  H^*(X)$
et $j^* : H^*(X \amalg Y) \to H^*(Y)$ sont surjectives,
$(i^*, j^*)$ l'est aussi. En effet, pour $a, a' \in H^*(X \amalg Y)$,
on a
$$
(i^*a, j^*a') = (i^*(a \cup e + a' \cup f), j^*(a \cup e + a' \cup f))
$$
Par symétrie, il suffit de montrer que $i^* : H^*(X \amalg Y) \to  H^*(X)$
est surjective. S'il existe un morphisme $Y \to X$, il existe un morphisme
$g : X \amalg Y \to X$ tel que $g \circ i = \text{id}_X$, et l'on conclut.
Pour achever la démonstration, observons que, pour montrer
la surjectivité de $i^*$, il suffit de le faire après produit tensoriel
par un espace vectoriel gradué non nul sur $F$. Ainsi, par l'axiome (B)(b)
et la non-nullité de la cohomologie (lemme \ref{lemma-unit}),
il suffit de démontrer la surjectivité de $i^*$ après avoir remplacé
$X$ et $Y$ par $X \times \Spec(k')$ et $Y \times \Spec(k')$
pour une extension finie séparable $k'/k$ convenable.
Si l'on choisit $k'$ de sorte qu'il existe un point fermé
$x \in X$ tel que $\kappa(x) = k'$ (ce qui est possible d'après
Variétés, lemme \ref{varieties-lemma-smooth-separable-closed-points-dense}),
il existe alors un morphisme $Y \times \Spec(k') \to X \times \Spec(k')$,
ce qui achève la démonstration.
\end{proof}

\begin{lemma}
\label{lemma-from-functor-to-weil}
Soit $k$ un corps et soit $F$ un corps de caractéristique $0$.
Supposons donné un foncteur $\mathbf{Q}$-linéaire
$$
G : M_k \longrightarrow \text{espaces vectoriels gradués sur }F\text{}
$$
de catégories monoïdales symétriques tel que $G(\mathbf{1}(1))$
ne soit non nul qu'en degré $-2$. On obtient alors des données (D0), (D1), (D2) et (D3)
satisfaisant à tous les axiomes (A), (B) et (C) ci-dessus.
\end{lemma}

\begin{proof}
Cette démonstration est la même que celle du
lemme \ref{lemma-from-functor-to-weil-classical} ;
nous conseillons vivement au lecteur de lire plutôt la démonstration de ce lemme.

\medskip\noindent
On obtient un foncteur contravariant de la catégorie des schémas projectifs
lisses sur $k$ dans celle des espaces vectoriels gradués sur $F$
en posant $H^*(X) = G(h(X))$. Par hypothèse, on dispose d'un isomorphisme
canonique
$$
H^*(X \times Y) = G(h(X \times Y)) = G(h(X) \otimes h(Y)) =
G(h(X)) \otimes G(h(Y)) = H^*(X) \otimes H^*(Y)
$$,
compatible avec les images réciproques. L'image réciproque par la diagonale
$\Delta : X \to X \times X$ donne une application canonique
$$
H^*(X) \otimes H^*(X) = H^*(X \times X) \to H^*(X)
$$
d'espaces vectoriels gradués, compatible avec les images réciproques.
Cela définit une structure fonctorielle de $F$-algèbre graduée sur
$H^*(X)$. Puisque $\Delta$ commute à la contrainte de commutativité
$h(X) \otimes h(X) \to h(X) \otimes h(X)$ (qui échange les facteurs),
que $G$ est un foncteur de catégories monoïdales symétriques (donc compatible avec
les contraintes de commutativité), et compte tenu de notre convention dans
Homologie, exemple \ref{homology-example-graded-vector-spaces},
on en déduit que $H^*(X)$ est une algèbre graduée
commutative ! Nous obtenons ainsi la donnée (D1).

\medskip\noindent
Puisque $\mathbf{1}(1)$ est inversible dans la catégorie des motifs,
$G(\mathbf{1}(1))$ est inversible dans la catégorie des
espaces vectoriels gradués sur $F$. Ainsi, $\sum_i \dim_F G^i(\mathbf{1}(1)) = 1$.
Par hypothèse, le seul degré où cet objet est non nul est $-2$.
La donnée (D0) est l'espace vectoriel $F(1) = G^{-2}(\mathbf{1}(1))$.
Puisque $G$ est un foncteur monoïdal symétrique,
on a $F(n) = G^{-2n}(\mathbf{1}(n))$ pour tout $n \in \mathbf{Z}$.
Il s'ensuit que
$$
H^{2r}(X)(r) = G^{2r}(h(X)) \otimes G^{-2r}(\mathbf{1}(r)) =
G^0(h(X)(r))
$$,
formule dont nous ferons fréquemment usage ci-dessous.

\medskip\noindent
Soit $X$ un schéma projectif lisse sur $k$. D'après
le lemme \ref{lemma-composition-correspondences}, on a
$$
\CH^r(X) \otimes \mathbf{Q} = \text{Corr}^r(\Spec(k), X) =
\Hom(\mathbf{1}(-r), h(X)) = \Hom(\mathbf{1}, h(X)(r))
$$
En appliquant le foncteur $G$, on obtient une application à valeurs dans
$\Hom(G(\mathbf{1}), G(h(X)(r)))$.
En prenant l'image de $1$ dans $G^0(\mathbf{1}) = F$
à valeurs dans $G^0(h(X)(r)) = H^{2r}(X)(r)$, on obtient
$$
\gamma :
\CH^r(X) \otimes \mathbf{Q} \longrightarrow H^{2r}(X)(r)
$$
C'est la donnée (D2).

\medskip\noindent
Soit $X$ un schéma projectif lisse non vide sur $k$,
équidimensionnel de dimension $d$. D'après
le lemme \ref{lemma-composition-correspondences}, on a
$$
\Mor(h(X)(d), \mathbf{1}) = \Mor((X, 1, d), (\Spec(k), 1, 0)) =
\text{Corr}^{-d}(X, \Spec(k)) = \CH_d(X)
$$
La classe du cycle $[X]$ dans $\CH_d(X)$ définit donc un morphisme
$h(X)(d) \to \mathbf{1}$. En appliquant $G$ et en prenant les composantes de degré $0$,
on obtient
$$
H^{2d}(X)(d) = G^0(h(X)(d)) \longrightarrow G^0(\mathbf{1}) = F
$$
Cette application $\int_X : H^{2d}(X)(d) \to F$ est la donnée (D3).

\medskip\noindent
Soit $X$ un schéma projectif lisse sur $k$,
non vide et équidimensionnel de dimension $d$. D'après le lemme \ref{lemma-dual},
$h(X)(d)$ est un dual à gauche de $h(X)$. Par conséquent,
$G(h(X)(d)) = H^*(X) \otimes_F F(d)[2d]$
est un dual à gauche de $H^*(X)$ dans la catégorie des espaces vectoriels gradués sur $F$.
Ici, $[n]$ désigne le foncteur de décalage des espaces vectoriels gradués.
D'après Homologie, lemme \ref{homology-lemma-left-dual-graded-vector-spaces},
on a $\sum_i \dim_F H^i(X) < \infty$, et
$\epsilon : h(X)(d) \otimes h(X) \to \mathbf{1}$ donne
des accouplements non dégénérés $H^{2d - i}(X)(d) \otimes_F H^i(X) \to F$.
Dans la démonstration du lemme \ref{lemma-dual}, nous avons vu que
$\epsilon$ est donné par $[\Delta]$ au moyen des identifications
$$
\Hom(h(X)(d) \otimes h(X), \mathbf{1}) =
\text{Corr}^{-d}(X \times X, \Spec(k)) =
\CH_d(X \times X)
$$
Ainsi, $\epsilon$ est la composée de $[X] : h(X)(d) \to \mathbf{1}$
et de $h(\Delta)(d) : h(X)(d) \otimes h(X) \to h(X)(d)$. Il s'ensuit
que les accouplements ci-dessus sont donnés par le cup-produit suivi de
$\int_X$. Cela démontre l'axiome (A).

\medskip\noindent
L'axiome (B) résulte de la compatibilité de $G$ avec les structures
tensorielles, supposée par hypothèse, et de notre construction du cup-produit ci-dessus.

\medskip\noindent
Axiome (C). Notre construction de $\gamma$ prend un cycle $\alpha$ sur $X$,
l'interprète comme une correspondance $a$ de $\Spec(k)$ dans $X$ d'un certain degré,
puis lui applique $G$. Si $f : Y \to X$ est un morphisme de schémas projectifs lisses
non vides et équidimensionnels sur $k$, alors
$f^!\alpha$ est l'image directe (!) de $\alpha$
par la correspondance $[\Gamma_f]$ de $X$ dans $Y$ ; voir
le lemme \ref{lemma-functor-and-cycles}. Ainsi,
$f^!\alpha$, considéré comme une correspondance de $\Spec(k)$ dans $Y$,
est égal à $a \circ [\Gamma_f]$ ; voir
le lemme \ref{lemma-composition-correspondences}.
Puisque $G$ est un foncteur, on en déduit que
$\gamma$ est compatible avec les images réciproques, ce qui est l'axiome (C)(a).

\medskip\noindent
Soit $f : Y \to X$ un morphisme de schémas projectifs lisses
non vides et équidimensionnels sur $k$,
et soit $\beta \in \CH^r(Y)$ un cycle sur $Y$. Il faut montrer que
$$
\int_Y \gamma(\beta) \cup f^*c = \int_X \gamma(f_*\beta) \cup c
$$
pour tout $c \in H^*(X)$. Soient $a, a^t, \eta_X, \eta_Y, [X], [Y]$
comme dans le lemme \ref{lemma-prep-dual}.
Soit $b$ le cycle $\beta$ considéré comme une correspondance de $\Spec(k)$ dans $Y$
de degré $r$. Alors $f_*\beta$, considéré comme une correspondance de
$\Spec(k)$ dans $X$, est égal à $a^t \circ b$ ; voir
les lemmes \ref{lemma-functor-and-cycles} et
\ref{lemma-composition-correspondences}.
L'égalité affichée ci-dessus sera vérifiée si l'on montre que
$$
h(X) = \mathbf{1} \otimes h(X)
\xrightarrow{b \otimes 1}
h(Y)(r) \otimes h(X)
\xrightarrow{1 \otimes a}
h(Y)(r) \otimes h(Y)
\xrightarrow{\eta_Y}
h(Y)(r)
\xrightarrow{[Y]}
\mathbf{1}(r - e)
$$
est égal à
$$
h(X) = \mathbf{1} \otimes h(X)
\xrightarrow{a^t \circ b \otimes 1}
h(X)(r + d - e) \otimes h(X)
\xrightarrow{\eta_X}
h(X)(r + d - e)
\xrightarrow{[X]}
\mathbf{1}(r - e)
$$
Cela résulte immédiatement du lemme \ref{lemma-prep-dual}.
L'axiome (C)(b) est donc vérifié.

\medskip\noindent
Pour démontrer l'axiome (C)(c), nous utilisons la discussion de
la remarque \ref{remark-replace-cup-product-classical}.
Il suffit donc de démontrer que $\gamma$ est compatible avec
les produits extérieurs. Soient $X$ et $Y$ des schémas projectifs lisses
non vides sur $k$, et soient $\alpha$ et $\beta$ des cycles sur ces schémas. Notons
$a$ et $b$ les correspondances associées de $\Spec(k)$ dans
$X$ et $Y$. Alors $\alpha \times \beta$ correspond à
la correspondance $a \otimes b$ de $\Spec(k)$ dans $X \otimes Y = X \times Y$.
La propriété voulue résulte donc de la compatibilité de $G$
avec les structures tensorielles de part et d'autre.

\medskip\noindent
L'axiome (C)(d) résulte de ce que le cycle $[\Spec(k)]$
correspond au morphisme identique de $h(\Spec(k))$.
Cela achève la démonstration du lemme.
\end{proof}

\begin{lemma}
\label{lemma-from-weil-to-functor}
Soit $k$ un corps et soit $F$ un corps de caractéristique $0$. À partir de données
(D0), (D1), (D2) et (D3) satisfaisant à (A), (B) et (C),
on peut construire un foncteur $\mathbf{Q}$-linéaire
$$
G : M_k \longrightarrow \text{espaces vectoriels gradués sur }F\text{}
$$
de catégories monoïdales symétriques tel que $H^*(X) = G(h(X))$.
\end{lemma}

\begin{proof}
La démonstration de ce lemme est la même que celle du
lemme \ref{lemma-from-weil-to-functor-classical} ;
nous conseillons vivement au lecteur de lire plutôt la démonstration de ce lemme.

\medskip\noindent
D'après le lemme \ref{lemma-characterize-motives}, il suffit de construire un foncteur
$G$ sur la catégorie des schémas projectifs lisses sur $k$,
dont les morphismes sont les correspondances de degré $0$, tel que
l'image de $G(c_2)$ dans $G(\mathbf{P}^1_k)$ soit un espace vectoriel gradué
inversible sur $F$.

\medskip\noindent
Soit $X$ un schéma projectif lisse sur $k$. Il existe une décomposition
canonique
$$
X = \coprod\nolimits_{0 \leq d \le \dim(X)} X_d
$$
en sous-schémas ouverts et fermés tels que $X_d$ soit équidimensionnel
de dimension $d$. Le lemme \ref{lemma-weil-additive} donne alors
$$
H^*(X) \longrightarrow \prod\nolimits_{0 \leq d \le \dim(X)} H^*(X_d)
$$
Si $Y$ est un second schéma projectif lisse sur $k$
et si l'on décompose de même $Y = \coprod Y_e$, alors
$$
\text{Corr}^0(X, Y) = \bigoplus \text{Corr}^0(X_d, Y_e)
$$
On a aussi $X \otimes Y = \coprod X_d \otimes Y_e$ dans
la catégorie des correspondances. Ces observations montrent
qu'il suffit de construire $G$ sur la catégorie dont les objets
sont les schémas projectifs lisses équidimensionnels sur $k$
et les morphismes les correspondances de degré $0$.
(Nous omettons certains détails.)

\medskip\noindent
Pour un schéma projectif lisse équidimensionnel
$X$ sur $k$, on pose $G(X) = H^*(X)$. Observons que $G(X) = 0$
si $X = \emptyset$ (lemme \ref{lemma-unit}). Les applications
de source ou de but $G(\emptyset)$ sont donc nulles, et nous pouvons
supposer désormais nos schémas non vides dans la discussion ci-dessous.

\medskip\noindent
Pour une correspondance $c \in \text{Corr}^0(X, Y)$ entre
schémas projectifs lisses non vides et équidimensionnels sur $k$,
considérons l'application $G(c) : G(X) = H^*(X) \to G(Y) = H^*(Y)$
donnée par la règle
$$
a \longmapsto
G(c)(a) = \text{pr}_{2, *}(\gamma(c) \cup \text{pr}_1^*a)
$$
Il est clair que $G(c)$ est additive en $c$, donc $\mathbf{Q}$-linéaire.
La compatibilité de $\gamma$ avec les images réciproques, les images directes
et les produits d'intersection, donnée par les axiomes (C)(a), (C)(b) et (C)(c),
montre que
$G(c' \circ c) = G(c') \circ G(c)$ si $c' \in \text{Corr}^0(Y, Z)$.
En effet, pour $a \in H^*(X)$, on a
\begin{align*}
(G(c') \circ G(c))(a)
& =
\text{pr}^{23}_{3, *}(\gamma(c') \cup
\text{pr}^{23, *}_2(\text{pr}^{12}_{2, *}(\gamma(c) \cup
\text{pr}^{12, *}_1a))) \\
& =
\text{pr}^{23}_{3, *}(\gamma(c') \cup
\text{pr}^{123}_{23, *}(\text{pr}^{123, *}_{12}(\gamma(c) \cup
\text{pr}^{12, *}_1 a))) \\
& =
\text{pr}^{23}_{3, *}
\text{pr}^{123}_{23, *}(
\text{pr}^{123, *}_{23}\gamma(c') \cup
\text{pr}^{123, *}_{12}\gamma(c) \cup
\text{pr}^{123, *}_1 a) \\
& =
\text{pr}^{23}_{3, *}
\text{pr}^{123}_{23, *}(
\gamma(\text{pr}^{123, *}_{23}c') \cup
\gamma(\text{pr}^{123, *}_{12}c) \cup
\text{pr}^{123, *}_1 a) \\
& =
\text{pr}^{13}_{3, *}
\text{pr}^{123}_{13, *}(
\gamma(\text{pr}^{123, *}_{23}c' \cdot \text{pr}^{123, *}_{12}c) \cup
\text{pr}^{123, *}_1 a) \\
& =
\text{pr}^{13}_{3, *}(
\gamma(\text{pr}^{123}_{13, *}(\text{pr}^{123, *}_{23}c' \cdot
\text{pr}^{123, *}_{12}c)) \cup
\text{pr}^{13, *}_1 a) \\
& =
G(c' \circ c)(a)
\end{align*},
avec des notations évidentes. La première égalité résulte des définitions.
La deuxième résulte de l'égalité
$\text{pr}^{23, *}_2 \circ \text{pr}^{12}_{2, *} =
\text{pr}^{123}_{23, *} \circ \text{pr}^{123, *}_{12}$,
qui découle immédiatement de la description de l'image directe
par les projections donnée dans le lemme \ref{lemma-pr2star}.
La troisième égalité résulte du lemme \ref{lemma-pushforward}
et du fait que $H^*$ est un foncteur.
La quatrième égalité résulte de l'axiome (C)(a) et du fait que
l'application de Gysin coïncide avec l'image réciproque plate pour les morphismes plats
(Homologie de Chow, lemme \ref{chow-lemma-lci-gysin-flat}).
La cinquième utilise l'axiome (C)(c), ainsi que
le lemme \ref{lemma-pushforward}, pour obtenir
$\text{pr}^{23}_{3, *} \circ \text{pr}^{123}_{23, *} =
\text{pr}^{13}_{3, *} \circ \text{pr}^{123}_{13, *}$.
La sixième utilise la formule de projection du
lemme \ref{lemma-pushforward}, ainsi que
l'axiome (C)(b), pour obtenir $
\text{pr}^{123}_{13, *}
\gamma(\text{pr}^{123, *}_{23}c' \cdot \text{pr}^{123, *}_{12}c) =
\gamma(\text{pr}^{123}_{13, *}(
\text{pr}^{123, *}_{23}c' \cdot \text{pr}^{123, *}_{12}c))$.
Enfin, la dernière égalité est la définition.

\medskip\noindent
Pour achever la démonstration que $G$ est un foncteur,
il faut montrer qu'il préserve les identités. Autrement dit, si
$1 = [\Delta] \in \text{Corr}^0(X, X)$ est l'identité dans la catégorie
des correspondances (lemme \ref{lemma-category-correspondences}),
il faut montrer que $G([\Delta]) = \text{id}$.
Cela résulte de la détermination
de $\gamma([\Delta])$ dans le lemme \ref{lemma-class-diagonal}
et du lemme \ref{lemma-pr2star}.
La construction de $G$ comme foncteur sur les
schémas projectifs lisses sur $k$ et les correspondances de degré $0$ est ainsi achevée.

\medskip\noindent
D'après le lemme \ref{lemma-base},
$G(\Spec(k)) = H^*(\Spec(k))$ est canoniquement isomorphe à $F$
comme $F$-algèbre. L'axiome de K\"unneth (B)(a)
montre que notre foncteur est compatible avec les produits tensoriels.
Il est donc un foncteur de catégories monoïdales symétriques.

\medskip\noindent
Il reste à vérifier que l'image de $G(c_2)$ dans
$G(\mathbf{P}^1_k) = H^*(\mathbf{P}^1_k)$
est un espace vectoriel gradué inversible sur $F$ (en particulier, nous ne savons pas encore
que $G$ se prolonge à $M_k$). D'après le lemme \ref{lemma-cohomology-P1},
la cohomologie n'est non nulle qu'en degrés $0$ et $2$,
où elle est de dimension $1$. L'égalité $1 = c_0 + c_2$ est une décomposition
de l'identité en somme d'idempotents orthogonaux dans
$\text{Corr}^0(\mathbf{P}^1_k, \mathbf{P}^1_k)$ ; voir
l'exemple \ref{example-decompose-P1}. De plus, on a $c_0 = a \circ b$, où
$a \in \text{Corr}^0(\Spec(k), \mathbf{P}^1_k)$ et
$b \in \text{Corr}^0(\mathbf{P}^1_k, \Spec(k))$, et où
$b \circ a = 1$ dans $\text{Corr}^0(\Spec(k), \Spec(k))$ ; voir la démonstration du
lemme \ref{lemma-inverse-h2}. Ainsi, $G(c_0)$ est le projecteur
sur la composante de degré $0$. Il s'ensuit que $G(c_2)$ est
le projecteur sur la composante de degré $2$, ce qui achève la démonstration.
\end{proof}

\begin{proposition}
\label{proposition-weil-cohomology-theory}
Soit $k$ un corps et soit $F$ un corps de caractéristique $0$. Il existe
une correspondance biunivoque ($1$ à $1$) entre :
\begin{enumerate}
\item les données (D0), (D1), (D2) et (D3) satisfaisant à (A), (B) et (C) ;
\item les foncteurs monoïdaux symétriques $\mathbf{Q}$-linéaires
$$
G : M_k \longrightarrow \text{espaces vectoriels gradués sur }F\text{}
$$
tels que $G(\mathbf{1}(1))$ ne soit non nul qu'en degré $-2$.
\end{enumerate}
\end{proposition}

\begin{proof}
Étant donné $G$ comme dans (2), on pose $H^*(X) = G(h(X))$ et l'on obtient des données
(D0), (D1), (D2) et (D3) satisfaisant à (A), (B) et (C) ;
voir le lemme \ref{lemma-from-functor-to-weil} et sa démonstration.

\medskip\noindent
Réciproquement, à partir de données (D0), (D1), (D2) et (D3)
satisfaisant à (A), (B) et (C), on obtient un foncteur $G$ comme dans (2)
par la construction de la démonstration du lemme \ref{lemma-from-weil-to-functor}.

\medskip\noindent
Nous omettons la démonstration détaillée que ces constructions
sont inverses l'une de l'autre.
\end{proof}











\section{Propriétés complémentaires}
\label{section-further}

\noindent
Nous démontrons dans cette section quelques conséquences supplémentaires
de données (D0), (D1), (D2) et (D3) satisfaisant à (A), (B) et (C),
comme à la section \ref{section-axioms}.

\begin{lemma}
\label{lemma-trace-disjoint-union}
Supposons que les données (D0), (D1), (D2) et (D3) satisfassent à (A), (B) et (C).
Soient $X, Y$ des schémas projectifs lisses non vides, tous deux équidimensionnels
de dimension $d$ sur $k$. Alors $\int_{X \amalg Y} = \int_X + \int_Y$.
\end{lemma}

\begin{proof}
Notons $i : X \to X \amalg Y$ et $j : Y \to X \amalg Y$ les coprojections.
D'après le lemme \ref{lemma-weil-additive}, l'application
$(i^*, j^*) : H^*(X \amalg Y) \to H^*(X) \times H^*(Y)$ est un isomorphisme.
L'énoncé signifie que, par l'isomorphisme
$(i^*, j^*) : H^{2d}(X \amalg Y)(d) \to H^{2d}(X)(d) \oplus H^{2d}(Y)(d)$,
l'application $\int_X + \int_Y$ se transforme en $\int_{X \amalg Y}$.
En effet, on a
$$
\int_{X \amalg Y} a =
\int_{X \amalg Y} i_*(i^*a) + j_*(j^*a) =
\int_X i^*a + \int_Y j^*a
$$,
où l'égalité $a = i_*(i^*a) + j_*(j^*a)$ a été établie dans
la démonstration du lemme \ref{lemma-weil-additive}.
\end{proof}

\begin{lemma}
\label{lemma-dim-0}
Supposons que les données (D0), (D1), (D2) et (D3) satisfassent à (A), (B) et (C).
Soit $X$ un schéma projectif lisse de dimension zéro sur $k$.
Alors :
\begin{enumerate}
\item $H^i(X) = 0$ pour $i \not = 0$ ;
\item $H^0(X)$ est une algèbre finie séparable sur $F$ ;
\item $\dim_F H^0(X) = \deg(X \to \Spec(F))$ ;
\item $\int_X : H^0(X) \to F$ est l'application trace ;
\item $\gamma([X]) = 1$ ;
\item $\int_X \gamma([X]) = \deg(X \to \Spec(k))$.
\end{enumerate}
\end{lemma}

\begin{proof}
On peut écrire $X = \Spec(k')$, où $k'$ est une algèbre finie séparable
sur $k$. Observons que $\deg(X \to \Spec(k)) = [k' : k]$.
Choisissons une extension galoisienne finie $k''/k$ contenant chacun des
facteurs de $k'$. (Rappelons qu'une $k$-algèbre finie séparable est
un produit d'extensions finies séparables de $k$.)
Posons $\Sigma = \Hom_k(k', k'')$. On obtient
$$
k' \otimes_k k'' = \prod\nolimits_{\sigma \in \Sigma} k''
$$
En posant $Y = \Spec(k'')$, l'axiome (B)(a) et le lemme \ref{lemma-weil-additive} donnent
$$
H^*(X) \otimes_F H^*(Y) =
\prod\nolimits_{\sigma \in \Sigma} H^*(Y)
$$
comme $F$-algèbres graduées commutatives. D'après le lemme \ref{lemma-unit},
la $F$-algèbre $H^*(Y)$ est non nulle. En comparant les dimensions des deux membres
de l'égalité affichée, on en déduit que $H^*(X)$ est concentré en degré $0$
et que $\dim_F H^0(X) = [k' : k]$. En appliquant cela à $Y$, on obtient
$H^*(Y) = H^0(Y)$. Puisque
$$
H^0(X) \otimes_F H^0(Y) = H^0(Y) \times \ldots \times H^0(Y)
$$
comme $F$-algèbres, il s'ensuit que $H^0(X)$ est une $F$-algèbre séparable,
car cela se vérifie après le changement de base fidèlement plat
$F \to H^0(Y)$.

\medskip\noindent
L'isomorphisme affiché ci-dessus est donné par l'application
$$
H^0(X) \otimes_F H^0(Y) \longrightarrow
\prod\nolimits_{\sigma \in \Sigma} H^0(Y),\quad
a \otimes b \longmapsto \prod\nolimits_\sigma \Spec(\sigma)^*a \cup b
$$
Par cet isomorphisme, on a $\int_{X \times Y} = \sum_\sigma \int_Y$, d'après
le lemme \ref{lemma-trace-disjoint-union}. Ainsi,
$$
\int_X a = \text{pr}_{1, *}(a \otimes 1) = \sum \Spec(\sigma)^*a
$$
dans $H^0(Y)$ ; la première égalité résulte du lemme \ref{lemma-pr2star}
et la seconde de l'observation précédente. Choisissons une
clôture algébrique $\overline{F}$ et
un homomorphisme de $F$-algèbres $\tau : H^0(Y) \to \overline{F}$.
Par changement de base, l'isomorphisme ci-dessus devient
$$
H^0(X) \otimes_F \overline{F} \longrightarrow
\prod\nolimits_{\sigma \in \Sigma} \overline{F},\quad
a \otimes b \longmapsto \prod\nolimits_\sigma \tau(\Spec(\sigma)^*a) b
$$
Il s'ensuit que les applications $a \mapsto \tau(\Spec(\sigma)^*a)$ forment l'ensemble complet
des plongements de $H^0(X)$ dans $\overline{F}$. En appliquant $\tau$
à la formule de $\int_X a$ obtenue ci-dessus, on en déduit
que $\int_X$ est l'application trace.
D'après le lemme \ref{lemma-unit}, on a $\gamma([X]) = 1$.
Enfin, on a $\int_X \gamma([X]) = \deg(X \to \Spec(k))$,
puisque $\gamma([X]) = 1$ et que la trace de $1$ est égale à $[k' : k]$.
\end{proof}

\begin{lemma}
\label{lemma-degrees-cycles}
Supposons que les données (D0), (D1), (D2) et (D3) satisfassent à (A), (B) et (C).
Soit $X$ un schéma projectif lisse non vide,
équidimensionnel de dimension $d$ sur $k$. Le diagramme
$$
\xymatrix{
\CH^d(X) \ar[r]_-\gamma \ar@{=}[d] &
H^{2d}(X)(d) \ar[d]^{\int_X} \\
\CH_0(X) \ar[r]^\deg & F
}
$$
est commutatif, où $\deg : \CH_0(X) \to \mathbf{Z}$ est le degré des
zéro-cycles étudié dans Homologie de Chow, section
\ref{chow-section-degree-zero-cycles}.
\end{lemma}

\begin{proof}
Soit $x$ un point fermé de $X$ dont le corps résiduel est séparable
sur $k$. Considérons $x$ comme un schéma et notons $i : x \to X$ le morphisme d'inclusion.
Pour éviter toute confusion, notons $\gamma' : \CH_0(x) \to H^0(x)$ l'application classe de cycle
de $x$. On a alors
$$
\int_X \gamma([x]) = \int_X \gamma(i_*[x]) =
\int_X i_*\gamma'([x]) = \int_x \gamma'([x]) = \deg(x \to \Spec(k))
$$
La deuxième égalité est l'axiome (C)(b), et la troisième
la définition de $i_*$ en cohomologie. La dernière égalité résulte du
lemme \ref{lemma-dim-0}. Cela démontre le lemme,
car $\CH_0(X)$ est engendré par les classes de points $x$ comme ci-dessus,
d'après le lemme \ref{lemma-generated-by-separable}.
\end{proof}

\begin{lemma}
\label{lemma-square-diagonal}
Supposons que les données (D0), (D1), (D2) et (D3) satisfassent à (A), (B) et (C).
Soit $X$ un schéma projectif lisse non vide sur $k$,
équidimensionnel de dimension $d$. On a
$$
\sum\nolimits_i (-1)^i\dim_F H^i(X) =
\deg(\Delta \cdot \Delta) = \deg(c_d(\mathcal{T}_{X/k}))
$$
\end{lemma}

\begin{proof}
Démontrons l'égalité de droite. On a
$[\Delta] \cdot [\Delta] = \Delta_*(\Delta^![\Delta])$
(Homologie de Chow, lemme \ref{chow-lemma-intersect-regularly-embedded}).
Puisque $\Delta_*$ préserve les degrés des $0$-cycles, il suffit de calculer
le degré de $\Delta^![\Delta]$. La classe $\Delta^![\Delta]$ est donnée
par le cap-produit de $[\Delta]$ avec
la classe de Chern de degré maximal du faisceau normal de $\Delta \subset X \times X$
(Homologie de Chow, lemme \ref{chow-lemma-gysin-fundamental}).
Puisque le faisceau conormal de $\Delta$ est $\Omega_{X/k}$
(Morphismes, lemme \ref{morphisms-lemma-differentials-diagonal}),
le faisceau normal est égal au faisceau tangent
$\mathcal{T}_{X/k} = \SheafHom_{\mathcal{O}_X}(\Omega_{X/k}, \mathcal{O}_X)$,
comme voulu.

\medskip\noindent
Démontrons l'égalité de gauche. D'après le lemme \ref{lemma-degrees-cycles}, on a
\begin{align*}
\deg([\Delta] \cdot [\Delta])
& =
\int_{X \times X} \gamma([\Delta]) \cup \gamma([\Delta]) \\
& =
\int_{X \times X} \Delta_*1 \cup \gamma([\Delta]) \\
& =
\int_{X \times X} \Delta_*(\Delta^*\gamma([\Delta])) \\
& =
\int_X \Delta^*\gamma([\Delta])
\end{align*}
Nous avons utilisé les lemmes \ref{lemma-push-unit} et
\ref{lemma-pushforward}.
Écrivons $\gamma([\Delta]) = \sum  e_{i, j} \otimes e'_{2d - i , j}$
comme dans le lemme \ref{lemma-class-diagonal}.
En rappelant que $\Delta^*$ est donné par le cup-produit
(remarque \ref{remark-replace-cup-product}), on obtient
$$
\int_X \sum\nolimits_{i, j} e_{i, j} \cup e'_{2d - i, j} =
\sum\nolimits_{i, j} \int_X e_{i, j} \cup e'_{2d - i, j} =
\sum\nolimits_{i, j} (-1)^i = \sum (-1)^i\beta_i
$$,
comme voulu.
\end{proof}

\begin{lemma}
\label{lemma-algebra-relations}
Soit $F$ un corps de caractéristique $0$.
Soient $F'$ et $F_i$, pour $i = 1, \ldots, r$,
des $F$-algèbres finies séparables. Soit $A$ une $F$-algèbre finie.
Soient $\sigma, \sigma' : A \to F'$ et $\sigma_i : A \to F_i$
des homomorphismes de $F$-algèbres. Supposons $\sigma$ et $\sigma'$ surjectifs.
S'il existe une relation
$$
\text{Tr}_{F'/F} \circ \sigma - \text{Tr}_{F'/F} \circ \sigma' =
n(\sum m_i \text{Tr}_{F_i/F} \circ \sigma_i)
$$,
où $n > 1$ et les $m_i$ sont des entiers, alors $\sigma = \sigma'$.
\end{lemma}

\begin{proof}
On peut écrire $A = \prod A_j$ comme un produit fini de
$F$-algèbres artiniennes locales $(A_j, \mathfrak m_j, \kappa_j)$ ; voir
Algèbre, lemme \ref{algebra-lemma-finite-dimensional-algebra} et
proposition \ref{algebra-proposition-dimension-zero-ring}.
Notons $A' = \prod \kappa_j$, où le produit porte sur les $j$
tels que $\kappa_j/k$ soit séparable. Chacune des applications
$\sigma, \sigma', \sigma_i$ se factorise alors par l'application
$A \to A'$. En remplaçant $A$ par $A'$, on peut supposer
que $A$ est une $F$-algèbre finie séparable.

\medskip\noindent
Choisissons une clôture algébrique $\overline{F}$. Posons
$\overline{A} = A \otimes_F \overline{F}$,
$\overline{F}' = F' \otimes_F \overline{F}$ et
$\overline{F}_i = F_i \otimes_F \overline{F}$.
Par changement de base de $\sigma$, $\sigma'$ et $\sigma_i$,
on obtient des homomorphismes de $\overline{F}$-algèbres $\overline{A} \to \overline{F}'$
et $\overline{A} \to \overline{F}_i$. De plus,
$\text{Tr}_{\overline{F}'/\overline{F}}$ est le changement de base
de $\text{Tr}_{F'/F}$, et de même pour
$\text{Tr}_{F_i/F}$. On peut donc remplacer
$F$ par $\overline{F}$ et se ramener au cas traité
dans le paragraphe suivant.

\medskip\noindent
Supposons $F$ algébriquement clos et $A$ une $F$-algèbre finie séparable.
Alors chacun de $A$, $F'$ et $F_i$ est un produit de copies de $F$.
Disons qu'un élément $e$ d'un produit
$F \times \ldots \times F$ de copies de $F$ est un idempotent minimal
s'il engendre l'un des facteurs, c'est-à-dire si
$e = (0, \ldots, 0, 1, 0, \ldots, 0)$. Soit $e \in A$ un idempotent minimal.
Puisque $\sigma$ et $\sigma'$
sont surjectifs, $\sigma(e)$ et $\sigma'(e)$ sont des idempotents
minimaux ou sont nuls. Si $\sigma \not = \sigma'$, on peut choisir
un idempotent minimal $e \in A$ tel que $\sigma(e) = 0$ et
$\sigma'(e) \not = 0$, ou inversement. Alors
$\text{Tr}_{F'/F}(\sigma(e)) = 0$ et
$\text{Tr}_{F'/F}(\sigma'(e)) = 1$, ou inversement.
D'autre part, $\sigma_i(e)$ est un idempotent,
donc $\text{Tr}_{F_i/F}(\sigma_i(e)) = r_i$ est un entier.
Il en résulte que
$$
-1 = \sum n m_i r_i = n (\sum m_i r_i)
\quad\text{ou}\quad
1 = \sum n m_i r_i = n (\sum m_i r_i)
$$,
ce qui est impossible.
\end{proof}

\begin{lemma}
\label{lemma-relations-classes-points}
Supposons que les données (D0), (D1), (D2) et (D3) satisfassent à (A), (B) et (C).
Soit $k'/k$ une extension finie séparable.
Soit $X$ un schéma projectif lisse sur $k'$.
Soient $x, x' \in X$ des points $k'$-rationnels.
Si $\gamma(x) \not = \gamma(x')$, alors
$[x] - [x']$ n'est divisible par aucun entier $n > 1$ dans $\CH_0(X)$.
\end{lemma}

\begin{proof}
Si $x$ et $x'$ appartiennent à des composantes irréductibles distinctes de $X$,
le résultat est évident. On peut donc supposer $X$ irréductible de dimension $d$.
Supposons $[x] - [x']$ divisible par $n > 1$ dans $\CH_0(X)$.
On peut écrire $[x] - [x'] = n(\sum m_i [x_i])$ dans $\CH_0(X)$,
pour certains points fermés $x_i \in X$
dont les corps résiduels sont séparables sur $k$, d'après
le lemme \ref{lemma-generated-by-separable}.
Alors
$$
\gamma([x]) - \gamma([x']) = n (\sum m_i \gamma([x_i]))
$$
dans $H^{2d}(X)(d)$. Notons $i^*, (i')^*, i_i^*$ les applications images réciproques
$H^0(X) \to H^0(x)$, $H^0(X) \to H^0(x')$ et $H^0(X) \to H^0(x_i)$.
Rappelons que $H^0(x)$ est une $F$-algèbre finie séparable
et que $\int_x : H^0(x) \to F$ est l'application trace
(lemme \ref{lemma-dim-0}), que nous noterons $\text{Tr}_x$.
Il en est de même pour $x'$ et $x_i$. Par la dualité de Poincar\'e sous la forme
de l'axiome (A)(b), l'égalité ci-dessus est duale de
$$
\text{Tr}_x \circ i^* - \text{Tr}_{x'} \circ (i')^* =
n(\sum m_i \text{Tr}_{x_i} \circ i_i^*)
$$,
égalité dans $\Hom_F(H^0(X), F)$. Enfin, observons que
$i^*$ et $(i')^*$ sont surjectives, puisque $x$ et $x'$ sont des points $k'$-rationnels
et que les composées $H^0(\Spec(k')) \to H^0(X) \to H^0(x)$
et $H^0(\Spec(k')) \to H^0(X) \to H^0(x')$ sont donc des isomorphismes.
Le lemme \ref{lemma-algebra-relations} donne $i^* = (i')^*$,
ce qui contredit l'hypothèse $\gamma([x]) \not = \gamma([x'])$.
\end{proof}

\begin{lemma}
\label{lemma-classes-points}
Supposons que les données (D0), (D1), (D2) et (D3) satisfassent à (A), (B) et (C).
Soit $k'/k$ une extension finie séparable et soit $X$ un schéma projectif lisse
géométriquement irréductible sur $k'$, de dimension $d$.
Alors $\gamma : \CH_0(X) \to H^{2d}(X)(d)$ se factorise par
$\deg : \CH_0(X) \to \mathbf{Z}$.
\end{lemma}

\begin{proof}
D'après le lemme \ref{lemma-generated-by-separable}, il suffit de montrer que,
pour des points fermés $x, x' \in X$ dont les corps résiduels sont séparables sur $k$,
on a $\deg(x') \gamma([x]) = \deg(x) \gamma([x'])$.

\medskip\noindent
Ramenons-nous d'abord au cas de points $k'$-rationnels. Soit $k''/k'$ une
extension galoisienne telle que $\kappa(x)$ et $\kappa(x')$ se plongent dans $k''$
sur $k$. Posons $Y = X \times_{\Spec(k')} \Spec(k'')$ et notons $p : Y \to X$
la projection. Par le choix de $k''/k'$, il existe un
point $k''$-rationnel $y$, resp.\ $y'$, de $Y$, d'image $x$, resp.\ $x'$.
Alors $p_*[y] = [k'' : \kappa(x)][x]$ et
$p_*[y'] = [k'' : \kappa(x')][x']$ dans $\CH_0(X)$.
Par la compatibilité avec les images directes de l'axiome (C)(b),
il suffit de démontrer $\gamma([y]) = \gamma([y'])$ dans $\CH^{2d}(Y)(d)$.
On est ainsi ramené à la discussion du paragraphe suivant.

\medskip\noindent
Supposons $x$ et $x'$ des points $k'$-rationnels. D'après
le lemme \ref{lemma-divide-difference-points},
il existe une extension finie séparable de corps $k''/k'$
telle que l'image réciproque $[y] - [y']$
de la différence $[x] - [x']$ devienne divisible
par un entier $n > 1$ sur $Y = X \times_{\Spec(k')} \Spec(k'')$.
(Notons que $y, y' \in Y$ sont des points $k''$-rationnels.)
D'après le lemme \ref{lemma-relations-classes-points}, on a
$\gamma([y]) = \gamma([y'])$ dans $H^{2d}(Y)(d)$.
Par la compatibilité avec l'image directe de l'axiome (C)(b),
on en déduit la même assertion pour $x$ et $x'$.
\end{proof}

\begin{lemma}
\label{lemma-injective}
Supposons que les données (D0), (D1), (D2) et (D3) satisfassent à (A), (B) et (C). Soit
$f : X \to Y$ un morphisme dominant de schémas projectifs lisses irréductibles
sur $k$. Alors $H^*(Y) \to H^*(X)$ est injective.
\end{lemma}

\begin{proof}
Il existe un sous-schéma fermé intègre $Z \subset X$ de même
dimension que $Y$ et qui se surjecte sur $Y$. Ainsi, $f_*[Z] = m[Y]$ pour un certain $m > 0$.
Alors $f_* \gamma([Z]) = m \gamma([Y]) = m$ dans $H^*(Y)$, d'après
le lemme \ref{lemma-unit}. Par la formule de projection
(lemme \ref{lemma-pushforward}),
on a donc $f_*(f^*a \cup \gamma([Z])) = m a$, et l'on conclut.
\end{proof}

\begin{lemma}
\label{lemma-otimes}
Supposons que les données (D0), (D1), (D2) et (D3) satisfassent à (A), (B) et (C). Soient
$k''/k'/k$ des algèbres finies séparables et soit $X$ un
schéma projectif lisse sur $k'$. Alors
$$
H^*(X) \otimes_{H^0(\Spec(k'))} H^0(\Spec(k'')) =
H^*(X \times_{\Spec(k')} \Spec(k''))
$$
\end{lemma}

\begin{proof}
Nous utiliserons les résultats du lemme \ref{lemma-dim-0} sans le mentionner de nouveau.
Écrivons
$$
k' \otimes_k k'' = k'' \times l
$$,
pour une certaine $k'$-algèbre finie séparable $l$. Posons
$F' = H^0(\Spec(k'))$, $F'' = H^0(\Spec(k''))$ et $G = H^0(\Spec(l))$.
Puisque $\Spec(k') \times \Spec(k'') = \Spec(k'') \amalg \Spec(l)$,
l'axiome (B)(a) et le lemme \ref{lemma-weil-additive}
donnent
$$
F' \otimes_F F'' = F'' \times G
$$
L'application de gauche à droite identifie $F''$ à $F' \otimes_{F'} F''$.
De même, on a
$$
H^*(X) \otimes_F F'' = H^*(X \times_{\Spec(k')} \Spec(k''))
\times H^*(X \times_{\Spec(k')} \Spec(l))
$$
comme modules sur $F' \otimes_F F'' = F'' \times G$. Le lemme est démontré.
\end{proof}






















\section{Théories cohomologiques de Weil, II}
\label{section-old}

\noindent
Pour nous, une théorie cohomologique de Weil sera l'analogue d'une
théorie cohomologique de Weil classique (section \ref{section-axioms-classical})
lorsque le corps de base $k$ n'est pas algébriquement clos.
À la section \ref{section-axioms}, nous avons énoncé des axiomes garantissant
que notre théorie cohomologique provient d'un foncteur monoïdal symétrique
sur la catégorie des motifs sur $k$. Il manque encore à ces axiomes
la condition $H^i(X) = 0$ pour $i < 0$, ainsi
qu'une condition sur $H^{2d}(X)(d)$ pour $X$ équidimensionnel de dimension $d$,
correspondant aux axiomes classiques (A)(c) et (A)(d).
Montrons d'abord au lecteur qu'il est nécessaire d'imposer
de telles conditions.

\begin{example}
\label{example-weird-weil}
Soient $k = \mathbf{C}$ et $F = \mathbf{C}$ tous deux égaux au corps
des nombres complexes. Pour $X$ projectif lisse sur $k$, notons
$H^{p, q}(X) = H^q(X, \Omega^p_{X/k})$. Soit $(H')^*$ le foncteur
qui envoie $X$ sur $(H')^*(X) = \bigoplus H^{p, q}(X)$, muni du
cup-produit usuel.
C'est une théorie cohomologique de Weil classique (insérer ici une référence ultérieure).
D'après la proposition \ref{proposition-weil-cohomology-theory-classical},
on obtient un foncteur monoïdal symétrique $\mathbf{Q}$-linéaire $G'$ de $M_k$
dans la catégorie des espaces vectoriels gradués sur $F$. Bien entendu, dans ce cas,
pour tout $M$ dans $M_k$, la valeur $G'(M)$ est naturellement bigraduée, c'est-à-dire
que l'on a
$$
(G')(M) = \bigoplus (G')^{p, q}(M),\quad
(G')^n = \bigoplus\nolimits_{n = p + q} (G')^{p, q}(M)
$$,
où $(G')^{p, q}$ est placé en degré total $p + q$, comme indiqué.
Construisons maintenant un foncteur monoïdal symétrique $\mathbf{Q}$-linéaire
$G$ à valeurs dans les espaces vectoriels gradués sur $F$ en posant
$$
G^n(M) = \bigoplus\nolimits_{n = 3p - q} (G')^{p, q}(M)
$$
Nous omettons la vérification que cela définit un foncteur monoïdal
symétrique (un point technique est que le choix des nombres impairs
$3$ et $-1$ ci-dessus assure la compatibilité du foncteur $G$ avec
les contraintes de commutativité).
Observons que $G(\mathbf{1}(1))$ est toujours placé en degré $-2$ !
Ainsi, par le lemme \ref{lemma-from-functor-to-weil-classical},
on obtient un foncteur $H^*$, des classes de cycles $\gamma$ et des applications traces
satisfaisant à tous les axiomes classiques (A), (B), (C), sauf peut-être
aux axiomes classiques (A)(a) et (A)(d).
Cependant, si $E$ est une courbe elliptique sur $k$, on trouve
$\dim H^{-1}(E) = 1$ : l'axiome (A)(a) est donc effectivement violé.
\end{example}

\begin{lemma}
\label{lemma-H-0-separable}
Supposons que les données (D0), (D1), (D2) et (D3) satisfassent à (A), (B) et (C).
Soit $X$ un schéma projectif lisse sur $k$.
Posons $k' = \Gamma(X, \mathcal{O}_X)$. Les conditions suivantes sont équivalentes :
\begin{enumerate}
\item il existe un nombre fini de points fermés $x_1, \ldots, x_r \in X$,
dont les corps résiduels sont séparables sur $k$, tels que
$H^0(X) \to H^0(x_1) \oplus \ldots \oplus H^0(x_r)$ soit injective ;
\item l'application $H^0(\Spec(k')) \to H^0(X)$ est un isomorphisme.
\end{enumerate}
Si ces conditions sont vérifiées, $H^0(X)$ est une algèbre finie séparable sur $F$.
Si $X$ est équidimensionnel de dimension $d$, les conditions (1) et (2)
équivalent aussi à :
\begin{enumerate}
\item[(3)] les classes des points fermés engendrent $H^{2d}(X)(d)$
comme module sur $H^0(X)$.
\end{enumerate}
\end{lemma}

\begin{proof}
Observons que l'énoncé a un sens, car $k'$ est une algèbre finie séparable
sur $k$ (Variétés, lemme
\ref{varieties-lemma-proper-geometrically-reduced-global-sections}) ;
ainsi, $\Spec(k')$ est lisse et projectif sur $k$.
La compatibilité de $H^*$ avec les sommes directes
(lemmes \ref{lemma-weil-additive} et \ref{lemma-trace-disjoint-union})
montre qu'il suffit de démontrer le lemme lorsque $X$ est connexe.
On peut donc supposer $X$ irréductible, et il faut établir
l'équivalence de (1), (2) et (3). Posons $d = \dim(X)$.
Alors $k'$ est un corps, extension finie séparable
de $k$, et $X$ est géométriquement irréductible sur $k'$ ; voir
Variétés, lemmes
\ref{varieties-lemma-proper-geometrically-reduced-global-sections} et
\ref{varieties-lemma-baby-stein}.

\medskip\noindent
D'après le lemme \ref{lemma-generated-by-separable}, on peut supposer que les points
fermés de (3) ont des corps résiduels séparables sur $k$.
Les axiomes (A)(a) et (A)(b) montrent que (1) et (3) sont équivalentes.

\medskip\noindent
Si (2) est vérifiée, choisissons un point fermé quelconque $x \in X$ dont le corps résiduel
est fini séparable sur $k'$. Alors
$H^0(\Spec(k')) = H^0(X) \to H^0(x)$ est injective, par exemple d'après
le lemme \ref{lemma-injective}.

\medskip\noindent
Supposons vérifiées les conditions équivalentes (1) et (3). Choisissons
$x_1, \ldots, x_r \in X$ comme dans (1), puis une extension finie séparable
$k''/k'$. D'après le lemme \ref{lemma-otimes}, on a
$$
H^0(X) \otimes_{H^0(\Spec(k'))} H^0(\Spec(k'')) =
H^0(X \times_{\Spec(k')} \Spec(k''))
$$
Pour montrer que
$H^0(\Spec(k')) \to H^0(X)$ est un isomorphisme,
on peut donc remplacer $k'$ par $k''$. On peut ainsi supposer que $x_1, \ldots, x_r$
sont des points $k'$-rationnels (chaque $x_i$ est remplacé par plusieurs
points, de sorte que $r$ augmente à cette étape). D'après le lemme \ref{lemma-classes-points},
$\gamma(x_1) = \gamma(x_2) = \ldots = \gamma(x_r)$.
Par l'axiome (A)(b), toutes les applications $H^0(X) \to H^0(x_i)$
sont égales. Cela signifie que (2) est vérifiée.

\medskip\noindent
Enfin, le lemme \ref{lemma-dim-0} entraîne
que $H^0(X)$ est une $F$-algèbre séparable si (1) est vérifiée.
\end{proof}

\begin{lemma}
\label{lemma-negative-cohomology}
Supposons que les données (D0), (D1), (D2) et (D3) satisfassent à (A), (B) et (C).
S'il existe un schéma projectif lisse $Y$ sur $k$ tel que
$H^i(Y)$ soit non nul pour un certain $i < 0$, il existe un schéma projectif lisse
équidimensionnel $X$ sur $k$ pour lequel les conditions équivalentes
du lemme \ref{lemma-H-0-separable} ne sont pas vérifiées pour $X$.
\end{lemma}

\begin{proof}
D'après le lemme \ref{lemma-weil-additive}, on peut supposer $Y$ irréductible,
donc a fortiori équidimensionnel. Si $i$ est impair, en remplaçant
$Y$ par $Y \times Y$, on obtient un exemple où $Y$ est équidimensionnel
et où $i = -2l$ pour un certain $l > 0$. Posons $X = Y \times (\mathbf{P}^1_k)^l$.
L'axiome (B)(a) donne
$$
H^0(X) \supset H^0(Y) \oplus
H^i(Y) \otimes_F H^2(\mathbf{P}^1_k)^{\otimes_F l}
$$,
les deux facteurs directs étant non nuls. Il est donc clair que $H^0(X)$ ne peut être
isomorphe à $H^0$ du spectre de
$\Gamma(X, \mathcal{O}_X) = \Gamma(Y, \mathcal{O}_Y)$,
puisque ce dernier s'envoie dans le premier facteur direct.
\end{proof}

\noindent
Nous sommes ainsi conduits à poser enfin la définition suivante.

\begin{definition}
\label{definition-weil-cohomology-theory}
Soit $k$ un corps et soit $F$ un corps de caractéristique $0$.
Une {\it théorie cohomologique de Weil} sur $k$ à coefficients dans $F$
est définie par des données (D0), (D1), (D2) et (D3) satisfaisant
à la dualité de Poincar\'e, à la formule de K\"unneth et à la compatibilité
avec les classes de cycles, plus précisément aux axiomes (A), (B) et (C)
de la section \ref{section-axioms},
et telles que les conditions équivalentes (1) et (2) du
lemme \ref{lemma-H-0-separable} soient vérifiées pour tout schéma projectif lisse $X$ sur $k$.
\end{definition}

\noindent
D'après le lemme \ref{lemma-negative-cohomology}, cela signifie aussi que tous
les groupes de cohomologie de degré négatif sont nuls. En particulier, si $k$ est
algébriquement clos, une théorie cohomologique de Weil ci-dessus, munie d'un isomorphisme
$F \to F(1)$, revient à se donner une théorie cohomologique de Weil classique.

\begin{remark}
\label{remark-betti-numbers-in-some-sense}
Soit $H^*$ une théorie cohomologique de Weil
(définition \ref{definition-weil-cohomology-theory}).
Soit $X$ un schéma projectif lisse géométriquement irréductible
de dimension $d$ sur $k'$, où $k'/k$ est une extension finie séparable de corps.
Supposons
$$
H^0(\Spec(k')) = F_1 \times \ldots \times F_r
$$,
pour certains corps $F_i$. On peut alors écrire
$$
H^*(X) = \prod\nolimits_{i = 1, \ldots, r}
H^*(X) \otimes_{H^0(\Spec(k'))} F_i
$$
La dernière hypothèse de la définition \ref{definition-weil-cohomology-theory}
signifie que $H^0(X)$ est libre de rang $1$ sur $\prod F_i$.
Autrement dit, chacun des facteurs
$H^0(X) \otimes_{H^0(\Spec(k'))} F_i$ est de dimension $1$ sur $F_i$.
La dualité de Poincar\'e donne la même assertion pour
la cohomologie en degré $2d$.
En revanche, on ignore s'il en est de même aux autres degrés.
Plus précisément, pour $0 < n < \dim(X)$, on ne sait pas si les entiers
$$
\dim_{F_i} H^n(X) \otimes_{H^0(\Spec(k'))} F_i
$$
sont indépendants de $i$ ! Cette question est étroitement liée à la question
ouverte suivante : étant donnés un corps de base algébriquement clos $\overline{k}$,
un corps de caractéristique zéro $F$, une théorie cohomologique de Weil classique
$H^*$ sur $\overline{k}$ de corps de coefficients $F$ et une variété projective lisse
$X$ sur $\overline{k}$, les nombres de Betti de $X$,
$$
\beta_i = \dim_F H^i(X)
$$,
sont-ils indépendants de $F$ et de la théorie cohomologique de Weil $H^*$ ?
\end{remark}

\begin{proposition}
\label{proposition-weil-cohomology-theory-again}
Soit $k$ un corps et soit $F$ un corps de caractéristique $0$.
Une théorie cohomologique de Weil revient à se donner un foncteur
monoïdal symétrique $\mathbf{Q}$-linéaire
$$
G : M_k \longrightarrow \text{espaces vectoriels gradués sur }F\text{}
$$
tel que :
\begin{enumerate}
\item $G(\mathbf{1}(1))$ ne soit non nul qu'en degré $-2$ ;
\item pour tout schéma projectif lisse $X$ sur $k$, en posant
$k' = \Gamma(X, \mathcal{O}_X)$, l'homomorphisme
$G(h(\Spec(k'))) \to G(h(X))$ d'espaces vectoriels gradués sur $F$
soit un isomorphisme en degré $0$.
\end{enumerate}
\end{proposition}

\begin{proof}
C'est une conséquence immédiate de la proposition \ref{proposition-weil-cohomology-theory}
et de la définition \ref{definition-weil-cohomology-theory}. On pourrait bien entendu
remplacer (2) par la condition que $G(h(X)) \to \bigoplus G(h(x_i))$
soit injective en degré $0$ pour un choix de points fermés
$x_1, \ldots, x_r \in X$ dont les corps résiduels sont séparables sur $k$.
\end{proof}







\section{Classes de Chern}
\label{section-chern}

\noindent
Nous expliquons dans cette section comment une première classe de Chern et
une formule du fibré projectif permettent d'obtenir toutes les classes de Chern.
On pourra consulter \cite{Grothendieck-chern}, bien que nos axiomes
soient légèrement différents.

\medskip\noindent
Soit $\mathcal{C}$ une catégorie de schémas possédant les propriétés suivantes :
\begin{enumerate}
\item Tout $X \in \Ob(\mathcal{C})$ est quasi-compact et quasi-séparé.
\item Si $X \in \Ob(\mathcal{C})$ et si $U \subset X$ est ouvert et fermé,
alors $U \to X$ est un morphisme de $\mathcal{C}$. Si $X' \to X$ est un morphisme
de $\mathcal{C}$ se factorisant par $U$, alors $X' \to U$ est un morphisme
de $\mathcal{C}$.
\item Si $X \in \Ob(\mathcal{C})$ et si $\mathcal{E}$ est un
$\mathcal{O}_X$-module localement libre de type fini, alors :
\begin{enumerate}
\item $p : \mathbf{P}(\mathcal{E}) \to X$ est un morphisme de $\mathcal{C}$ ;
\item pour un morphisme $f : X' \to X$ de $\mathcal{C}$,
le morphisme induit $\mathbf{P}(f^*\mathcal{E}) \to \mathbf{P}(\mathcal{E})$
est un morphisme de $\mathcal{C}$ ;
\item si $\mathcal{E} \to \mathcal{F}$ est une surjection sur un autre
$\mathcal{O}_X$-module localement libre de type fini, l'immersion fermée
$\mathbf{P}(\mathcal{F}) \to \mathbf{P}(\mathcal{E})$
est un morphisme de $\mathcal{C}$.
\end{enumerate}
\end{enumerate}
Supposons ensuite donné un foncteur contravariant $A$ de
la catégorie $\mathcal{C}$ dans la catégorie des algèbres graduées.
Ici, une algèbre graduée $A$ est une
$\mathbf{Z}$-algèbre unitaire associative, non nécessairement commutative, $A$, munie d'une graduation
$A = \bigoplus_{i \geq 0} A^i$. Pour un morphisme
$f : X' \to X$ de $\mathcal{C}$, nous notons $f^* : A(X) \to A(X')$
l'homomorphisme d'algèbres induit.
Nous noterons le produit de $a, b \in A(X)$ par $a \cup b$.
Enfin, nous supposons
donnée, pour tout objet $X$ de $\mathcal{C}$, une application additive
$$
c_1^A : \Pic(X) \longrightarrow A^1(X)
$$
Nous supposons vérifiés les axiomes suivants :
\begin{enumerate}
\item Pour $X \in \Ob(\mathcal{C})$ et $\mathcal{L} \in \Pic(X)$,
l'élément $c_1^A(\mathcal{L})$ appartient au centre de l'algèbre $A(X)$.
\item Si $X \in \Ob(\mathcal{C})$ et si $X = U \amalg V$, avec $U$ et $V$
ouverts et fermés, alors $A(X) = A(U) \times A(V)$ par les applications induites
$A(X) \to A(U)$ et $A(X) \to A(V)$.
\item Si $f : X' \to X$ est un morphisme de $\mathcal{C}$ et si $\mathcal{L}$
est un $\mathcal{O}_X$-module inversible, alors $f^*c_1^A(\mathcal{L}) =
c_1^A(f^*\mathcal{L})$.
\item Pour $X \in \Ob(\mathcal{C})$ et un $\mathcal{O}_X$-module localement libre
$\mathcal{E}$ de rang constant $r$, considérons le morphisme
$p : P = \mathbf{P}(\mathcal{E}) \to X$ de $\mathcal{C}$.
Alors l'application
$$
\bigoplus\nolimits_{i = 0, \ldots, r - 1} A(X)
\longrightarrow A(P),\quad
(a_0, \ldots, a_{r - 1}) \longmapsto
\sum c_1^A(\mathcal{O}_P(1))^i \cup p^*(a_i)
$$
est bijective.
\item Soit $X \in \Ob(\mathcal{C})$ et soit $\mathcal{E} \to \mathcal{F}$
une surjection de $\mathcal{O}_X$-modules localement libres de type fini,
de rangs $r + 1$ et $r$. Notons
$i : P' = \mathbf{P}(\mathcal{F}) \to \mathbf{P}(\mathcal{E}) = P$
le morphisme d'inclusion correspondant. C'est un morphisme de $\mathcal{C}$
qui présente $P'$ comme un diviseur de Cartier effectif sur $P$. Alors, pour
$a \in A(P)$ tel que $i^*a = 0$, on a
$a \cup c_1^A(\mathcal{O}_P(P')) = 0$.
\end{enumerate}
Pour énoncer notre résultat, rappelons que $\textit{Vect}(X)$ désigne
la catégorie (exacte) des $\mathcal{O}_X$-modules localement libres de type fini.
Dans Catégories dérivées de schémas, section \ref{perfect-section-K-groups},
nous avons défini le $K$-groupe de degré zéro
$K_0(\textit{Vect}(X))$ de cette catégorie.
Nous avons également vu que $K_0(\textit{Vect}(X))$ est un anneau ; voir
Catégories dérivées de schémas, remarque \ref{perfect-remark-K-ring}.

\begin{proposition}
\label{proposition-chern-class}
Dans la situation ci-dessus, il existe une unique règle associant
à tout $X \in \Ob(\mathcal{C})$ une « classe de Chern totale »
$$
c^A : K_0(\textit{Vect}(X)) \longrightarrow  \prod\nolimits_{i \geq 0} A^i(X)
$$
possédant les propriétés suivantes :
\begin{enumerate}
\item Pour $X \in \Ob(\mathcal{C})$, on a
$c^A(\alpha + \beta) = c^A(\alpha) c^A(\beta)$
et $c^A(0) = 1$.
\item Si $f : X' \to X$ est un morphisme de $\mathcal{C}$, alors
$f^* \circ c^A =  c^A \circ f^*$.
\item Pour $X \in \Ob(\mathcal{C})$ et $\mathcal{L} \in \Pic(X)$,
on a $c^A([\mathcal{L}]) = 1 + c_1^A(\mathcal{L})$.
\end{enumerate}
\end{proposition}

\begin{proof}
Soit $X \in \Ob(\mathcal{C})$ et soit $\mathcal{E}$ un
$\mathcal{O}_X$-module localement libre de type fini. Montrons d'abord comment définir
un élément $c^A(\mathcal{E}) \in A(X)$.

\medskip\noindent
Soit d'abord $X = \bigcup X_r$ la décomposition en
sous-schémas ouverts et fermés tels que $\mathcal{E}|_{X_r}$ soit
de rang constant $r$. Puisque $X$ est quasi-compact, cette décomposition
est finie. Ainsi, $A(X) = \prod A(X_r)$. Il suffit donc de définir
$c^A(\mathcal{E})$ lorsque $\mathcal{E}$ est de rang constant $r$. Dans
ce cas, soit $p : P \to X$ le fibré projectif de $\mathcal{E}$.
On définit de manière unique des éléments $c_i^A(\mathcal{E}) \in A^i(X)$
pour $i \geq 0$, tels que $c_0^A(\mathcal{E}) = 1$ et que l'équation
\begin{equation}
\label{equation-chern-classes}
\sum\nolimits_{i = 0}^r
(-1)^i c_1(\mathcal{O}_P(1))^i \cup p^*c^A_{r - i}(\mathcal{E})
= 0
\end{equation}
soit vérifiée. Comme d'habitude, on pose
$c^A(\mathcal{E}) = c_0^A(\mathcal{E}) + c_1^A(\mathcal{E}) + \ldots
+ c_r^A(\mathcal{E})$ dans $A(X)$.

\medskip\noindent
Si $\mathcal{E}$ est inversible, alors
$c^A(\mathcal{E}) = 1 + c_1^A(\mathcal{L})$.
Cela résulte immédiatement de la construction ci-dessus.

\medskip\noindent
Les éléments $c_i^A(\mathcal{E})$ appartiennent au centre de $A(X)$.
Pour le montrer, on peut supposer $\mathcal{E}$ de rang constant $r$.
Soit $p : P \to X$ le fibré projectif correspondant.
Si $a \in A(X)$, alors $p^*a \cup (-1)^r c_1(\mathcal{O}_P(1))^r =
(-1)^r c_1(\mathcal{O}_P(1))^r \cup p^*a$ ; la même propriété vaut donc
pour tous les autres termes de l'expression définissant $c_i^A(\mathcal{E})$,
ce qui permet de conclure.

\medskip\noindent
Si $f : X' \to X$ est un morphisme de $\mathcal{C}$, alors
$f^*c_i^A(\mathcal{E}) = c_i^A(f^*\mathcal{E})$.
Pour le montrer, on peut supposer $\mathcal{E}$ de rang constant $r$.
Soient $p : P \to X$ et $p' : P' \to X'$ les fibrés
projectifs associés à $\mathcal{E}$ et à $f^*\mathcal{E}$.
Le morphisme induit $g : P' \to P$ est un morphisme de $\mathcal{C}$.
L'image réciproque par $g$ de l'égalité définissant $c_i^A(\mathcal{E})$
est l'équation correspondante pour $f^*\mathcal{E}$, d'où la conclusion.

\medskip\noindent
Soit $X \in \Ob(\mathcal{C})$. Considérons une suite exacte courte
$$
0 \to \mathcal{L} \to \mathcal{E} \to \mathcal{F} \to 0
$$
de $\mathcal{O}_X$-modules localement libres de type fini, avec $\mathcal{L}$ inversible.
Alors
$$
c^A(\mathcal{E}) = c^A(\mathcal{L}) c^A(\mathcal{F})
$$
En effet, par la construction de $c^A_i$, on peut supposer $\mathcal{E}$
de rang constant $r + 1$ et $\mathcal{F}$ de rang constant $r$.
L'inclusion
$$
i : P' = \mathbf{P}(\mathcal{F}) \longrightarrow \mathbf{P}(\mathcal{E}) = P
$$
est un morphisme de $\mathcal{C}$ et est le schéma des zéros d'une section
régulière du module inversible
$\mathcal{L}^{\otimes -1} \otimes \mathcal{O}_P(1)$.
L'élément
$$
\sum\nolimits_{i = 0}^r (-1)^i c_1^A(\mathcal{O}_P(1))^i \cup
p^*c^A_i(\mathcal{F})
$$
a une image réciproque nulle sur $P'$ par définition. On obtient donc
$$
\left(c_1^A(\mathcal{O}_P(1)) - c_1^A(\mathcal{L})\right) \cup
\left(\sum\nolimits_{i = 0}^r (-1)^i c_1^A(\mathcal{O}_P(1))^i \cup
p^*c^A_i(\mathcal{F})\right) = 0
$$
dans $A^*(P)$, par l'hypothèse (5) sur notre cohomologie $A$.
Par définition de $c_1^A(\mathcal{E})$,
cela donne l'égalité voulue.

\medskip\noindent
Soit $X \in \Ob(\mathcal{C})$. Considérons une suite exacte courte
$$
0 \to \mathcal{E} \to \mathcal{F} \to \mathcal{G} \to 0
$$
de $\mathcal{O}_X$-modules localement libres de type fini. Alors
$$
c^A(\mathcal{F}) = c^A(\mathcal{E}) c^A(\mathcal{G})
$$
En effet, par la construction de $c^A_i$, on peut supposer
$\mathcal{E}$, $\mathcal{F}$ et $\mathcal{G}$
de rangs constants $r$, $s$ et $t$. On raisonne par récurrence sur $r$.
Le cas $r = 1$ a été traité ci-dessus. Si $r > 1$, il suffit de vérifier
l'égalité après image réciproque par le morphisme $\mathbf{P}(\mathcal{E}^\vee) \to X$.
On peut donc supposer qu'il existe un sous-module inversible
$\mathcal{L} \subset \mathcal{E}$ tel que
$\mathcal{E}' = \mathcal{E}/\mathcal{L}$ et
$\mathcal{F}' = \mathcal{E}/\mathcal{L}$ soient tous deux localement libres de type fini
(de rangs $s - 1$ et $t - 1$). Alors
$$
c^A(\mathcal{E}) = c^A(\mathcal{L}) c^A(\mathcal{E}')
\quad\text{et}\quad
c^A(\mathcal{F}) = c^A(\mathcal{L}) c^A(\mathcal{F}')
$$
Puisque l'on a la suite exacte courte
$$
0 \to \mathcal{E}' \to \mathcal{F}' \to \mathcal{G} \to 0
$$,
l'hypothèse de récurrence donne
$$
c^A(\mathcal{F}') = c^A(\mathcal{E}') c^A(\mathcal{G})
$$
Le résultat découle donc d'un calcul formel.

\medskip\noindent
Nous pouvons maintenant, pour $X \in \Ob(\mathcal{C})$,
définir $c^A : K_0(\textit{Vect}(X)) \to A(X)$.
On envoie le générateur $[\mathcal{E}]$ sur $c^A(\mathcal{E})$
et l'on prolonge multiplicativement. Ainsi, par exemple,
$c^A(-[\mathcal{E}]) = c^A(\mathcal{E})^{-1}$ est l'inverse
formel de $a^A([\mathcal{E}])$.
La multiplicativité sur les suites exactes courtes établie ci-dessus
assure que cette définition convient.

\medskip\noindent
Unicité. Soit $X \in \Ob(\mathcal{C})$ et soit $\mathcal{E}$
un $\mathcal{O}_X$-module localement libre de type fini. Nous voulons montrer
que les conditions (1), (2) et (3) du lemme déterminent de manière unique
$c^A([\mathcal{E}])$. On peut supposer $\mathcal{E}$
de rang constant $r$, ce qui utilise déjà (2). On raisonne alors par récurrence sur $r$.
Si $r = 1$, l'unicité résulte de (3).
Si $r > 1$, on passe, à l'aide de (2), à l'image réciproque sur le fibré projectif $p : P \to X$ ;
on peut ainsi supposer que l'on dispose d'une suite exacte courte
$0 \to \mathcal{E}' \to \mathcal{E} \to \mathcal{E}'' \to 0$,
où $\mathcal{E}'$ et $\mathcal{E}''$ sont de rangs inférieurs.
Par l'hypothèse de récurrence, $c^A(\mathcal{E}')$ et $c^A(\mathcal{E}'')$
sont déterminés de manière unique. L'unicité pour $\mathcal{E}$ résulte alors
de l'axiome (1).
\end{proof}

\begin{lemma}
\label{lemma-splitting-principle}
Dans la situation ci-dessus, soit $X \in \Ob(\mathcal{C})$. Soit $\mathcal{E}_i$
une famille finie de $\mathcal{O}_X$-modules localement libres de rangs $r_i$.
Il existe un morphisme $p : P \to X$ dans $\mathcal{C}$ tel que :
\begin{enumerate}
\item $p^* : A(X) \to A(P)$ soit injective ;
\item chaque $p^*\mathcal{E}_i$ admette une filtration dont les quotients successifs
$\mathcal{L}_{i, 1}, \ldots, \mathcal{L}_{i, r_i}$
soient des $\mathcal{O}_P$-modules inversibles.
\end{enumerate}
\end{lemma}

\begin{proof}
On peut supposer $r_i \geq 1$ pour tout $i$. On démontre le lemme par récurrence
sur $\sum (r_i - 1)$. Si cet entier vaut $0$, alors $\mathcal{E}_i$
est inversible pour tout $i$, et l'on conclut en prenant $\pi = \text{id}_X$.
Sinon, choisissons un $i$ tel que $r_i > 1$ et considérons
le fibré projectif $p : P \to X$ associé à $\mathcal{E}_i$.
On a une suite exacte courte
$$
0 \to \mathcal{F} \to p^*\mathcal{E}_i \to \mathcal{O}_P(1) \to 0
$$
de $\mathcal{O}_P$-modules localement libres de type fini, de rangs $r_i - 1$,
$r_i$ et $1$. Observons que $p^* : A(X) \to A(P)$ est injective
par hypothèse. En appliquant l'hypothèse de récurrence aux
$\mathcal{O}_P$-modules localement libres de type fini $\mathcal{F}$ et $p^*\mathcal{E}_{i'}$
pour $i' \not = i$, on obtient un morphisme $p' : P' \to P$ possédant
les propriétés de l'énoncé. La composée
$p \circ p' : P' \to X$ convient alors.
\end{proof}

\begin{lemma}
\label{lemma-chern-classes-E-tensor-L}
Soit $X \in \Ob(\mathcal{C})$. Soit $\mathcal{E}$ un
$\mathcal{O}_X$-module localement libre de type fini et soit $\mathcal{L}$ un
$\mathcal{O}_X$-module inversible. Alors
$$
c^A_i({\mathcal E} \otimes {\mathcal L})
=
\sum\nolimits_{j = 0}^i
\binom{r - i + j}{j} c^A_{i - j}({\mathcal E}) \cup c^A_1({\mathcal L})^j
$$
\end{lemma}

\begin{proof}
Par la construction de $c^A_i$, on peut supposer $\mathcal{E}$
de rang constant $r$. Soient $p : P \to X$ et $p' : P' \to X$
les fibrés projectifs associés à $\mathcal{E}$ et à
$\mathcal{E} \otimes \mathcal{L}$.
Il existe alors un isomorphisme $g : P \to P'$ tel que
$g^*\mathcal{O}_{P'}(1) = \mathcal{O}_P(1) \otimes p^*\mathcal{L}$.
Voir Constructions, lemme \ref{constructions-lemma-twisting-and-proj}.
Ainsi,
$$
g^*c_1^A(\mathcal{O}_{P'}(1)) =
c_1^A(\mathcal{O}_P(1)) + p^*c_1^A(\mathcal{L})
$$
L'égalité voulue en découle formellement, en utilisant la définition
des classes de Chern par l'équation (\ref{equation-chern-classes}).
\end{proof}

\begin{proposition}
\label{proposition-chern-character}
Dans la situation ci-dessus, supposons que $A(X)$ soit une $\mathbf{Q}$-algèbre pour tout
$X \in \Ob(\mathcal{C})$. Il existe alors une unique règle associant
à tout $X \in \Ob(\mathcal{C})$ un « caractère de Chern »
$$
ch^A : K_0(\textit{Vect}(X)) \longrightarrow
\prod\nolimits_{i \geq 0} A^i(X)
$$
possédant les propriétés suivantes :
\begin{enumerate}
\item $ch^A$ est un homomorphisme d'anneaux pour tout $X \in \Ob(\mathcal{C})$.
\item Si $f : X' \to X$ est un morphisme de $\mathcal{C}$, alors
$f^* \circ ch^A =  ch^A \circ f^*$.
\item Pour $X \in \Ob(\mathcal{C})$ et $\mathcal{L} \in \Pic(X)$,
on a $ch^A([\mathcal{L}]) = \exp(c_1^A(\mathcal{L}))$.
\end{enumerate}
\end{proposition}

\begin{proof}
Soit $X \in \Ob(\mathcal{C})$ et soit $\mathcal{E}$ un
$\mathcal{O}_X$-module localement libre de type fini. Montrons d'abord comment définir
le rang $r^A(\mathcal{E}) \in A^0(X)$. Soit $X = \bigcup X_r$
la décomposition en sous-schémas ouverts et fermés tels que
$\mathcal{E}|_{X_r}$ soit de rang constant $r$. Puisque $X$ est quasi-compact, cette
décomposition est finie ; écrivons-la $X = X_0 \amalg X_1 \amalg \ldots \amalg X_n$.
Alors $A(X) = A(X_0) \times A(X_1) \times \ldots \times A(X_n)$. On peut donc
définir $r^A(\mathcal{E}) = (0, 1, \ldots, n) \in A^0(X)$.

\medskip\noindent
Soient $P_p(c_1, \ldots, c_p)$ les polynômes construits dans
Homologie de Chow, exemple \ref{chow-example-power-sum}.
On peut alors définir
$$
ch^A(\mathcal{E}) = r^A(\mathcal{E}) +
\sum\nolimits_{i \geq 1} (1/i!)
P_i(c^A_1(\mathcal{E}), \ldots, c^A_i(\mathcal{E}))
\in \prod\nolimits_{i \geq 0} A^i(X)
$$,
où les $ci^A$ sont les classes de Chern de la
proposition \ref{proposition-chern-class}.
Les propriétés (2) et (3) du lemme sont immédiatement vérifiées.

\medskip\noindent
Il reste à démontrer les trois assertions suivantes :
\begin{enumerate}
\item Si $0 \to \mathcal{E}_1 \to \mathcal{E} \to \mathcal{E}_2 \to 0$
est une suite exacte courte de $\mathcal{O}_X$-modules localement libres de type fini
sur $X \in \Ob(\mathcal{C})$, alors
$ch^A(\mathcal{E}) = ch^A(\mathcal{E}_1) + ch^A(\mathcal{E}_2)$.
\item Si $\mathcal{E}_1$ et $\mathcal{E}_2 \to 0$ sont des
$\mathcal{O}_X$-modules localement libres de type fini sur $X \in \Ob(\mathcal{C})$, alors
$ch^A(\mathcal{E}_1 \otimes \mathcal{E}_2) =
ch^A(\mathcal{E}_1) ch^A(\mathcal{E}_2)$.
\end{enumerate}
La première montrera que $ch^A$ se factorise par
$K_0(\textit{Vect}(X))$, et la première et la seconde ensemble
montreront que $ch^A$ est un homomorphisme d'anneaux.

\medskip\noindent
Pour démontrer ces assertions, on peut se ramener au cas où $\mathcal{E}_1$
et $\mathcal{E}_2$ sont de rangs constants $r_1$ et $r_2$. Les
égalités dans $A^0(X)$ sont alors immédiates. Pour les égalités en degrés
supérieurs, le lemme \ref{lemma-splitting-principle} permet de
supposer que $\mathcal{E}_1$ et $\mathcal{E}_2$ admettent des filtrations
dont les composantes graduées sont des modules inversibles
$\mathcal{L}_{1, j}$, pour $j = 1, \ldots, r_1$, et
$\mathcal{L}_{2, j}$, pour $j = 1, \ldots, r_2$.
La multiplicativité des classes de Chern donne
$$
c_i^A(\mathcal{E}_1) =
s_i(c_1^A(\mathcal{L}_{1, 1}), \ldots, c_1^A(\mathcal{L}_{1, r_1}))
$$,
où $s_i$ est la $i$-ième fonction symétrique élémentaire, comme dans
Homologie de Chow, exemple \ref{chow-example-power-sum}.
Il en est de même pour $c_i^A(\mathcal{E}_2)$. Dans le cas (1), on obtient
$$
c_i^A(\mathcal{E}) =
s_i(c_1^A(\mathcal{L}_{1, 1}), \ldots, c_1^A(\mathcal{L}_{1, r_1}),
c_1^A(\mathcal{L}_{2, 1}), \ldots, c_1^A(\mathcal{L}_{2, r_2}))
$$,
et, dans le cas (2),
$$
c_i^A(\mathcal{E}_1 \otimes \mathcal{E}_2) =
s_i(c_1^A(\mathcal{L}_{1, 1}) + c_1^A(\mathcal{L}_{2, 1}),
\ldots, c_1^A(\mathcal{L}_{1, r_1}) + c_1^A(\mathcal{L}_{2, r_2}))
$$
Par définition des polynômes $P_i$, cela signifie que
$$
P_i(c^A_1(\mathcal{E}_1), \ldots, c^A_i(\mathcal{E}_1)) =
\sum\nolimits_{j = 1, \ldots, r_1} c_1^A(\mathcal{L}_{1, j})^i
$$,
et de même pour $\mathcal{E}_2$. Dans le cas (1), on a aussi
$$
P_i(c^A_1(\mathcal{E}), \ldots, c^A_i(\mathcal{E})) =
\sum\nolimits_{j = 1, \ldots, r_1} c_1^A(\mathcal{L}_{1, j})^i +
\sum\nolimits_{j = 1, \ldots, r_2} c_1^A(\mathcal{L}_{2, j})^i
$$
Dans le cas (2), on obtient de même
$$
P_i(c^A_1(\mathcal{E}_1 \otimes \mathcal{E}_2), \ldots,
c^A_i(\mathcal{E}_1 \otimes \mathcal{E}_2)) =
\sum\nolimits_{j = 1, \ldots, r_1}
\sum\nolimits_{j' = 1, \ldots, r_2}
(c_1^A(\mathcal{L}_{1, j}) + c_1^A(\mathcal{L}_{2, j'}))^i
$$
Les égalités voulues découlent donc d'identités élémentaires
entre polynômes symétriques.

\medskip\noindent
Nous omettons la démonstration de l'unicité.
\end{proof}

\begin{lemma}
\label{lemma-adams-and-chern}
Dans la situation ci-dessus, soit $X \in \Ob(\mathcal{C})$.
Si $\psi^2$ est défini comme dans
Homologie de Chow, lemme \ref{chow-lemma-second-adams-operator},
et si $c^A$ et $ch^A$ sont définis comme dans les
propositions \ref{proposition-chern-class} et
\ref{proposition-chern-character},
alors on a $c^A_i(\psi^2(\alpha)) = 2^i c^A_i(\alpha)$ et
$ch^A_i(\psi^2(\alpha)) = 2^i ch^A_i(\alpha)$
pour tout $\alpha \in K_0(\textit{Vect}(X))$.
\end{lemma}

\begin{proof}
Observons que l'application $\prod_{i \geq 0} A^i(X) \to \prod_{i \geq 0} A^i(X)$,
qui multiplie par $2^i$ sur $A^i(X)$, est un homomorphisme d'anneaux. Puisque $\psi^2$
est également un homomorphisme d'anneaux, il suffit donc de démontrer les formules
sur des générateurs additifs de $K_0(\textit{Vect}(X))$. On peut ainsi supposer $\alpha = [\mathcal{E}]$
pour un $\mathcal{O}_X$-module localement libre de type fini $\mathcal{E}$.
La construction des classes de Chern de $\mathcal{E}$ ramène immédiatement
au cas où $\mathcal{E}$ est de rang constant $r$.
On peut alors choisir un morphisme projectif lisse $p : P \to X$
tel que la restriction $A^*(X) \to A^*(P)$ soit injective
et que $p^*\mathcal{E}$ admette une filtration finie dont
les composantes graduées soient des $\mathcal{O}_P$-modules inversibles $\mathcal{L}_j$ ; voir
le lemme \ref{lemma-splitting-principle}. Alors
$[p^*\mathcal{E}] = \sum [\mathcal{L}_j]$, d'où
$\psi^2([p^*\mathcal{E}]) = \sum [\mathcal{L}_j^{\otimes 2}]$
par définition de $\psi^2$. En posant $x_j  = c^A_1(\mathcal{L}_j)$,
on a
$$
c^A(\alpha) = \prod (1 + x_j)
\quad\text{et}\quad
c^A(\psi^2(\alpha)) = \prod (1 + 2 x_j)
$$
dans $\prod A^i(P)$, et l'on a
$$
ch^A(\alpha) = \sum \exp(x_j)
\quad\text{et}\quad
ch^A(\psi^2(\alpha)) = \sum \exp(2 x_j)
$$
dans $\prod A^i(P)$. Le résultat voulu découle de ces formules.
\end{proof}







\section{Puissances extérieures et K-groupes}
\label{section-lambda-operations}

\noindent
Nous nous limitons aux constructions nécessaires pour définir les opérateurs lambda.
Soit $X$ un schéma. Rappelons que $\textit{Vect}(X)$ désigne
la catégorie des $\mathcal{O}_X$-modules localement libres de type fini.
Rappelons également que nous avons construit un $K$-groupe de degré zéro
$K_0(\textit{Vect}(X))$ associé à cette catégorie dans
Catégories dérivées de schémas, section \ref{perfect-section-K-groups}.
Enfin, $K_0(\textit{Vect}(X))$ est un anneau ; voir
Catégories dérivées de schémas, remarque \ref{perfect-remark-K-ring}.

\begin{lemma}
\label{lemma-lambda-operations}
Soit $X$ un schéma. Il existe des applications
$$
\lambda^r : K_0(\textit{Vect}(X)) \longrightarrow K_0(\textit{Vect}(X))
$$
qui envoient $[\mathcal{E}]$ sur $[\wedge^r(\mathcal{E})]$
lorsque $\mathcal{E}$ est un $\mathcal{O}_X$-module localement libre de type fini,
et qui sont compatibles avec les images réciproques.
\end{lemma}

\begin{proof}
Considérons l'anneau $R = K_0(\textit{Vect}(X))[[t]]$, où $t$ est une
indéterminée. Pour un $\mathcal{O}_X$-module localement libre de type fini
$\mathcal{E}$, posons
$$
c(\mathcal{E}) = \sum\nolimits_{i = 0}^\infty [\wedge^i(\mathcal{E})] t^i
$$
dans $R$. Montrons que, pour une suite exacte courte
$$
0 \to \mathcal{E}' \to \mathcal{E} \to \mathcal{E}'' \to 0
$$
de $\mathcal{O}_X$-modules localement libres de type fini,
on a $c(\mathcal{E}) = c(\mathcal{E}') c(\mathcal{E}'')$.
Cette assertion entraîne que $c$ se prolonge en une application
$$
c :  K_0(\textit{Vect}(X)) \longrightarrow R
$$
qui transforme l'addition dans $K_0(\textit{Vect}(X))$ en la multiplication dans $R$.
En écrivant $c(\alpha) = \sum \lambda^i(\alpha) t^i$, on obtient les opérateurs
$\lambda^i$ cherchés.

\medskip\noindent
Pour démontrer l'assertion, considérons la suite exacte courte comme une
filtration de $\mathcal{E}$ à $2$ crans. On obtient une filtration
induite de $\wedge^r(\mathcal{E})$ à $r + 1$ crans, dont les
sous-quotients sont
$$
\wedge^r(\mathcal{E}'),
\wedge^{r - 1}(\mathcal{E}') \otimes \mathcal{E}'',
\wedge^{r - 2}(\mathcal{E}') \otimes \wedge^2(\mathcal{E}''), \ldots
\wedge^r(\mathcal{E}'')
$$
Ainsi, $[\wedge^r(\mathcal{E})]$ est égal à
$$
\sum\nolimits_{i = 0}^r
[\wedge^{r - i}(\mathcal{E}')] [\wedge^i(\mathcal{E}'')]
$$,
et le résultat s'en déduit aisément par un calcul algébrique élémentaire.
\end{proof}










\section{Théories cohomologiques de Weil, III}
\label{section-c1}

\noindent
Soit $k$ un corps et soit $F$ un corps de caractéristique zéro.
Supposons que l'on dispose des données suivantes :
\begin{enumerate}
\item[(D0)] Un espace vectoriel de dimension $1$ sur $F$, noté $F(1)$.
\item[(D1)] Un foncteur contravariant $H^*(-)$ de la catégorie des schémas projectifs
lisses sur $k$ dans la catégorie des
$F$-algèbres graduées commutatives.
\item[(D2')] Pour tout schéma projectif lisse $X$ sur $k$, un homomorphisme
$c_1^H : \Pic(X) \to H^2(X)(1)$ de groupes abéliens.
\end{enumerate}
Nous employons la terminologie, les notations et les conventions relatives
à (D0) et (D1) exposées à la section \ref{section-axioms}.
Pour un schéma projectif lisse $X$ sur $k$ et un
$\mathcal{O}_X$-module inversible $\mathcal{L}$, la classe de cohomologie
$c_1^H(\mathcal{L}) \in H^2(X)(1)$ de (D2')
est parfois appelée la {\it première classe de Chern de $\mathcal{L}$
en cohomologie}.

\medskip\noindent
Voici la liste des axiomes.
\begin{enumerate}
\item[(A1)] $H^*$ est compatible avec les sommes finies.
\item[(A2)] $c_1^H$ est compatible avec les images réciproques.
\item[(A3)] Soit $X$ un schéma projectif lisse sur $k$.
Soit $\mathcal{E}$ un $\mathcal{O}_X$-module localement libre de rang $r \geq 1$.
Considérons le morphisme $p : P = \mathbf{P}(\mathcal{E}) \to X$.
Alors l'application
$$
\bigoplus\nolimits_{i = 0, \ldots, r - 1} H^*(X)(-i)
\longrightarrow H^*(P),\quad
(a_0, \ldots, a_{r - 1}) \longmapsto
\sum c_1^H(\mathcal{O}_P(1))^i \cup p^*(a_i)
$$
est un isomorphisme d'espaces vectoriels sur $F$.
\item[(A4)] Soit $i : Y \to X$ l'inclusion d'un diviseur de Cartier
effectif sur $k$, avec $X$ et $Y$ lisses et projectifs
sur $k$. Pour $a \in H^*(X)$ tel que
$i^*a = 0$, on a $a \cup c_1^H(\mathcal{O}_X(Y)) = 0$.
\item[(A5)] $H^*$ est compatible avec les produits finis.
\item[(A6)] Soit $X$ un schéma projectif lisse non vide sur $k$,
équidimensionnel de dimension $d$. Il existe une application
$F$-linéaire $\lambda : H^{2d}(X)(d) \to F$ telle que
$(\text{id} \otimes \lambda) \gamma([\Delta]) =  1$ dans $H^*(X)$.
\item[(A7)] Si $b : X' \to X$ est l'éclatement d'un centre
lisse dans un schéma projectif lisse $X$ sur $k$\footnote{Alors $X'$
est également lisse et projectif sur $k$ ; voir
Compléments sur les morphismes, lemme \ref{more-morphisms-lemma-blowup}.}, alors
$b^* : H^*(X) \to H^*(X')$ est injective.
\item[(A8)] Si $X$ est un schéma projectif lisse sur $k$ et si
$k' = \Gamma(X, \mathcal{O}_X)$, alors l'application $H^0(\Spec(k')) \to H^0(X)$
est un isomorphisme.
\item[(A9)] Soit $X$ un schéma projectif lisse non vide sur $k$,
équidimensionnel de dimension $d$. Soit $i : Y \to X$ un diviseur de Cartier
effectif non vide, lisse sur $k$. Pour $a \in H^{2d - 2}(X)(d - 1)$,
on a $\lambda_Y(i^*(a)) = \lambda_X(a \cup c_1^H(\mathcal{O}_X(Y))$, où
$\lambda_Y$ et $\lambda_X$ sont celles de l'axiome (A6) pour $X$ et $Y$.
\end{enumerate}
Précisons le sens de chacun de ces axiomes.
Les axiomes (A3), (A4) et (A7) sont clairs tels qu'ils sont énoncés.

\medskip\noindent
À propos de (A1). Cela signifie que $H^*(\emptyset) = 0$ et que
$(i^*, j^*) : H^*(X \amalg Y) \to H^*(X) \times H^*(Y)$
est un isomorphisme, où $i$ et $j$ sont les coprojections.

\medskip\noindent
À propos de (A2). Cela signifie que, pour un morphisme $f : X \to Y$ de schémas projectifs
lisses sur $k$ et un $\mathcal{O}_Y$-module inversible $\mathcal{N}$,
on a $f^*c_1^H(\mathcal{L}) = c_1^H(f^*\mathcal{L})$.

\medskip\noindent
À propos de (A5). Cela signifie que $H^*(\Spec(k)) = F$ et que, pour $X$ et $Y$ projectifs
lisses sur $k$, l'application $H^*(X) \otimes_F H^*(Y) \to H^*(X \times Y)$,
$a \otimes b \mapsto p^*(a) \cup q^*(b)$, est un isomorphisme,
où $p$ et $q$ sont les projections.

\medskip\noindent
À propos de (A6). Soit $X$ un schéma projectif lisse non vide sur $k$,
équidimensionnel de dimension $d$. D'après le lemme \ref{lemma-cycle-classes},
sous les axiomes (A1) -- (A4), on peut considérer la classe de la diagonale
$$
\gamma([\Delta]) \in
H^{2d}(X \times X)(d) = \bigoplus\nolimits_i H^i(X) \otimes_F H^{2d - i}(X)(d)
$$,
où la décomposition tensorielle provient de l'axiome (A5).
Étant donnée une application $F$-linéaire $\lambda : H^{2d}(X)(d) \to F$, on peut aussi considérer
$\lambda$ comme une application $F$-linéaire $\lambda : H^*(X)(d) \to F$, en la précomposant
avec la projection sur $H^{2d}(X)(d)$. Avec cette convention, l'axiome (A6)
a bien un sens.

\medskip\noindent
À propos de (A8). Soit $X$ un schéma projectif lisse sur $k$.
Alors $k' = \Gamma(X, \mathcal{O}_X)$ est une
$k$-algèbre finie séparable (Variétés, lemme
\ref{varieties-lemma-proper-geometrically-reduced-global-sections}) ;
ainsi, $\Spec(k')$ est lisse et projectif sur $k$.
On peut donc appliquer $H^*$ à $\Spec(k')$, et l'axiome (A8) a bien un sens.

\medskip\noindent
À propos de (A9). Nous verrons dans la remarque \ref{remark-trace} que, sous
les axiomes (A1) -- (A7), l'application $\lambda$ de l'axiome (A6) est unique.

\begin{lemma}
\label{lemma-chern-classes}
Supposons que les données (D0), (D1) et (D2') satisfassent aux axiomes (A1), (A2), (A3) et (A4).
Il existe une unique règle associant à tout schéma projectif lisse $X$ sur $k$
un homomorphisme d'anneaux
$$
ch^H :
K_0(\textit{Vect}(X))
\longrightarrow
\prod\nolimits_{i \geq 0} H^{2i}(X)(i)
$$,
compatible avec les images réciproques et tel que
$ch^H(\mathcal{L}) = \exp(c_1^H(\mathcal{L}))$
pour tout $\mathcal{O}_X$-module inversible $\mathcal{L}$.
\end{lemma}

\begin{proof}
Cela résulte immédiatement de la proposition \ref{proposition-chern-character},
appliquée à la catégorie des schémas projectifs lisses sur $k$,
au foncteur $A : X \mapsto \bigoplus_{i \geq 0} H^{2i}(X)(i)$
et à l'application $c_1^H$.
\end{proof}

\begin{lemma}
\label{lemma-cycle-classes}
Supposons que les données (D0), (D1) et (D2') satisfassent aux axiomes (A1), (A2), (A3) et (A4).
Il existe une unique règle associant à tout schéma projectif lisse $X$ sur $k$
un homomorphisme d'anneaux gradués
$$
\gamma : \CH^*(X) \longrightarrow \bigoplus\nolimits_{i \geq 0} H^{2i}(X)(i)
$$,
compatible avec les images réciproques et tel que $ch^H(\alpha) = \gamma(ch(\alpha))$
pour $\alpha$ dans $K_0(\textit{Vect}(X))$.
\end{lemma}

\begin{proof}
Rappelons que l'on dispose d'un isomorphisme
$$
K_0(\textit{Vect}(X)) \otimes \mathbf{Q}
\longrightarrow \CH^*(X) \otimes \mathbf{Q},\quad
\alpha \longmapsto ch(\alpha) \cap [X]
$$ ;
voir Homologie de Chow, lemme \ref{chow-lemma-K-tensor-Q}. C'est un isomorphisme
d'anneaux d'après Homologie de Chow, remarque \ref{chow-remark-chern-character-K}.
On définit $\gamma$ par la formule $\gamma(\alpha) = ch^H(\alpha')$,
où $ch^H$ est celui du lemme \ref{lemma-chern-classes} et où
$\alpha' \in K_0(\textit{Vect}(X))$ est tel que
$ch(\alpha') \cap [X] = \alpha$ dans $\CH^*(X) \otimes \mathbf{Q}$.

\medskip\noindent
La construction $\alpha \mapsto \gamma(\alpha)$ est compatible
avec les images réciproques, car $ch^H$ et la formation des classes de Chern
le sont tous deux ; voir
le lemme \ref{lemma-chern-classes} et
Homologie de Chow, remarque \ref{chow-remark-gysin-chern-classes}.

\medskip\noindent
Il reste à voir que $\gamma$ respecte la graduation.
Soit $\psi^2 : K_0(\textit{Vect}(X)) \to K_0(\textit{Vect}(X))$
le deuxième opérateur d'Adams ; voir Homologie de Chow,
lemme \ref{chow-lemma-second-adams-operator}.
Si $\alpha \in \CH^i(X)$ et si
$\alpha' \in K_0(\textit{Vect}(X)) \otimes \mathbf{Q}$
est l'unique élément tel que $ch(\alpha') \cap [X] = \alpha$,
nous avons vu dans
Homologie de Chow, section \ref{chow-section-intersection-regular},
que $\psi^2(\alpha') = 2^i \alpha'$.
On en déduit $ch^H(\alpha') \in H^{2i}(X)(i)$
par le lemme \ref{lemma-adams-and-chern}, comme voulu.
\end{proof}

\begin{lemma}
\label{lemma-divide-pullback-good-blowing-up}
Soit $b : X' \to X$ l'éclatement d'un schéma projectif lisse
sur $k$ le long d'un sous-schéma fermé lisse $Z \subset X$.
Considérons le diagramme
$$
\xymatrix{
E \ar[r]_j \ar[d]_\pi & X' \ar[d]^b \\
Z \ar[r]^i & X
}
$$
Supposons qu'il existe un élément de $K_0(X)$ dont la restriction à
$Z$ soit égale à la classe de $\mathcal{C}_{Z/X}$ dans $K_0(Z)$.
Supposons que toute composante irréductible de $Z$ soit de codimension $r$ dans $X$.
Il existe alors un cycle $\theta \in \CH^{r - 1}(X')$
tel que $b^![Z] = [E] \cdot \theta$ dans $\CH^r(X')$ et
$\pi_*j^!(\theta) = [Z]$ dans $\CH^r(Z)$.
\end{lemma}

\begin{proof}
Le schéma $X$ est lisse et projectif sur $k$, de sorte que
$K_0(X) = K_0(\textit{Vect}(X))$. Voir
Catégories dérivées de schémas, lemmes
\ref{perfect-lemma-resolution-property-ample} et
\ref{perfect-lemma-K-is-old-K}.
Soit $\alpha \in K_0(\text{Vect}(X))$ un élément
dont la restriction à $Z$ est $\mathcal{C}_{Z/X}$.
D'après Homologie de Chow, lemme \ref{chow-lemma-minus-adams-operator},
il existe un élément $\alpha^\vee$ dont la restriction est
$\mathcal{C}_{Z/X}^\vee$. Par la formule de l'éclatement
(Homologie de Chow, lemme \ref{chow-lemma-blow-up-formula}),
on a
$$
b^![Z] = b^!i_*[Z] = j_* res(b^!)([Z]) =
j_*(c_{r - 1}(\mathcal{F}^\vee) \cap \pi^*[Z]) =
j_*(c_{r - 1}(\mathcal{F}^\vee) \cap [E])
$$,
où $\mathcal{F}$ est le noyau de la surjection
$\pi^*\mathcal{C}_{Z/X} \to \mathcal{C}_{E/X'}$.
Observons que $b^*\alpha^\vee - [\mathcal{O}_{X'}(E)]$
est un élément de $K_0(\text{Vect}(X'))$
dont la restriction est $[\pi^*\mathcal{C}_{Z/X}^\vee] - [\mathcal{C}_{E/X'}^\vee] =
[\mathcal{F}^\vee]$ sur $E$. Puisque le cap-produit par les classes de Chern
commute à $j_*$, l'expression ci-dessus est égale à
$$
c_{r - 1}(b^*\alpha^\vee - [\mathcal{O}_{X'}(E)]) \cap [E]
$$
dans le groupe de Chow de $X'$. Ainsi, en posant
$$
\theta = c_{r - 1}(b^*\alpha^\vee - [\mathcal{O}_{X'}(E)]) \cap [X']
$$,
on obtient la première relation $\theta \cdot [E] = b^![Z]$,
par exemple d'après Homologie de Chow, lemme \ref{chow-lemma-identify-chow-for-smooth}.
Pour la seconde relation, observons que
$$
j^!\theta = j^!(c_{r - 1}(b^*\alpha^\vee - [\mathcal{O}_{X'}(E)]) \cap [X'])
= c_{r - 1}(\mathcal{F}^\vee) \cap j^![X'] =
c_{r - 1}(\mathcal{F}^\vee) \cap [E]
$$
dans les groupes de Chow de $E$. Pour montrer que son image par $\pi_*$ est égale à $[Z]$,
il suffit de montrer que le degré du cycle de codimension $r - 1$
$(-1)^{r - 1}c_{r - 1}(\mathcal{F}) \cap [E]$ sur les fibres de $\pi$ vaut $1$.
Nous omettons ce calcul.
\end{proof}

\begin{lemma}
\label{lemma-A5-A6-imply}
Supposons que les données (D0), (D1) et (D2') satisfassent aux axiomes (A1) -- (A4)
et (A7). Soit $X$ un schéma projectif lisse sur $k$ et soit $Z \subset X$
un sous-schéma fermé lisse, tel que toute composante irréductible de $Z$
soit de codimension $r$ dans $X$. Supposons que la classe de
$\mathcal{C}_{Z/X}$ dans $K_0(Z)$ soit la restriction d'un élément de $K_0(X)$.
Si $a \in H^*(X)$ et si $a|_Z = 0$ dans $H^*(Z)$, alors
$\gamma([Z]) \cup a = 0$.
\end{lemma}

\begin{proof}
Soit $b : X' \to X$ l'éclatement. D'après (A7), il suffit de montrer
que
$$
b^*(\gamma([Z]) \cup a) = b^*\gamma([Z]) \cup b^*a = 0
$$
Par le lemme \ref{lemma-divide-pullback-good-blowing-up}, on a
$$
b^*\gamma([Z]) = \gamma(b^*[Z]) =
\gamma([E] \cdot \theta) =
\gamma([E]) \cup \gamma(\theta)
$$
Puisque $b^*a$ se restreint à zéro sur $E$ et que
$\gamma([E]) = c^H_1(\mathcal{O}_{X'}(E))$, l'axiome (A4) donne le résultat voulu.
\end{proof}

\begin{lemma}
\label{lemma-poincare-duality}
Supposons que les données (D0), (D1) et (D2') satisfassent aux axiomes (A1) -- (A7).
Alors l'axiome (A) de la section \ref{section-axioms} est vérifié avec
$\int_X = \lambda$ comme dans l'axiome (A6).
\end{lemma}

\begin{proof}
Soit $X$ un schéma projectif lisse non vide sur $k$,
équidimensionnel de dimension $d$. Nous allons montrer que l'espace vectoriel gradué sur $F$
$H^*(X)(d)[2d]$ est un dual à gauche de $H^*(X)$. Cela donnera le résultat voulu par
Homologie, lemme \ref{homology-lemma-left-dual-graded-vector-spaces}.
Nous utiliserons l'axiome (A5), qui donne en particulier
$$
H^*(X \times X)(d) =
\bigoplus H^i(X) \otimes H^j(X)(d) =
\bigoplus H^i(X)(d) \otimes H^j(X)
$$
Définissons une application
$$
\eta : F \longrightarrow H^*(X \times X)(d)
$$
par multiplication par $\gamma([\Delta]) \in H^{2d}(X \times X)(d)$.
D'autre part, définissons une application
$$
\epsilon :
H^*(X \times X)(d) \longrightarrow H^*(X)(d) \xrightarrow{\lambda} F
$$
en prenant d'abord l'image réciproque $\Delta^*$ par le morphisme diagonal
$\Delta : X \to X \times X$, puis l'application $F$-linéaire
$\lambda : H^{2d}(X)(d) \to F$ de l'axiome (A6), précomposée
avec la projection $H^*(X)(d) \to H^{2d}(X)(d)$.
Pour montrer que $H^*(X)(d)$ est un dual à gauche de $H^*(X)$,
il faut montrer que la composée des applications
$$
\eta \otimes 1 :
H^*(X) \longrightarrow H^*(X \times X \times X)(d)
$$
et
$$
1 \otimes \epsilon : H^*(X \times X \times X)(d) \longrightarrow H^*(X)
$$
est l'identité. Pour $a \in H^*(X)$,
la composée envoie $a$ sur
$$
(1 \otimes \lambda)(\Delta_{23}^*(q_{12}^*\gamma([\Delta]) \cup q_3^*a)) =
(1 \otimes \lambda)(\gamma([\Delta]) \cup p_2^*a)
$$,
où $q_i : X \times X \times X \to X$ et
$q_{ij} : X \times X \times X \to X \times X$ sont les projections,
$\Delta_{23} : X \times X \to X \times X \times X$ est la diagonale et
$p_i : X \times X \to X$ sont les projections.
L'égalité résulte de $\Delta_{23}^*(q_{12}^*\gamma([\Delta]) =
\Delta_{23}^*\gamma([\Delta \times X]) = \gamma([\Delta])$
et de $\Delta_{23}^* q_3^*a = p_2^*a$.
Puisque $\gamma([\Delta]) \cup p_1^*a = \gamma([\Delta]) \cup p_2^*a$
(voir ci-dessous), l'expression ci-dessus se simplifie en
$$
(1 \otimes \lambda)(\gamma([\Delta]) \cup p_1^*a) = a
$$
par notre choix de $\lambda$, comme voulu. La seconde condition,
$(\epsilon \otimes 1) \circ (1 \otimes \eta) = \text{id}$,
de Catégories, définition \ref{categories-definition-dual},
se démontre exactement de la
même manière.

\medskip\noindent
Notons que $p_1^*a$ et $\text{pr}_2^*a$ se restreignent à la même
classe de cohomologie sur $\Delta \subset X \times X$. De plus, on
a $\mathcal{C}_{\Delta/X \times X} = \Omega^1_\Delta$, qui
est la restriction de $p_1^*\Omega^1_X$. Le
lemme \ref{lemma-A5-A6-imply} donne donc
$\gamma([\Delta]) \cup p_1^*a = \gamma([\Delta]) \cup p_2^*a$,
ce qui achève la démonstration.
\end{proof}

\begin{remark}[Unicité des applications traces]
\label{remark-trace}
Supposons que les données (D0), (D1) et (D2') satisfassent aux axiomes (A1) -- (A7).
Soit $X$ un schéma projectif lisse sur $k$, non vide
et équidimensionnel de dimension $d$. En combinant les arguments
des démonstrations du lemme \ref{lemma-poincare-duality} et de
Homologie, lemme \ref{homology-lemma-left-dual-graded-vector-spaces},
on voit que
$$
\gamma([\Delta]) \in \bigoplus\nolimits_i H^i(X) \otimes H^{2d - i}(X)(d)
$$
définit une dualité parfaite entre $H^i(X)$ et $H^{2d - i}(X)(d)$
pour tout $i$.
En particulier, l'application linéaire $\int_X = \lambda : H^{2d}(X)(d) \to F$ de
l'axiome (A6) est unique ! Nous appellerons désormais l'application linéaire $\int_X$
l'application trace de $X$.
\end{remark}

\begin{lemma}
\label{lemma-trace-product}
Supposons que les données (D0), (D1) et (D2') satisfassent aux axiomes (A1) -- (A7).
Alors l'axiome (B) de la section \ref{section-axioms} est vérifié.
\end{lemma}

\begin{proof}
L'axiome (B)(a) résulte immédiatement de l'axiome (A5).
Soient $X$ et $Y$ des schémas projectifs lisses non vides sur $k$,
équidimensionnels de dimensions $d$ et $e$. Pour établir l'axiome (B)(b),
observons que la diagonale $\Delta_{X \times Y}$ de $X \times Y$
est le produit d'intersection des images réciproques des diagonales
$\Delta_X$ de $X$ et $\Delta_Y$ de $Y$ par les projections
$p : X \times Y \times X \times Y \to X \times X$ et
$q : X \times Y \times X \times Y \to Y \times Y$.
La compatibilité de $\gamma$ avec les produits d'intersection entraîne alors
que
$$
\gamma([\Delta_{X \times Y}]) \in
H^{2d + 2e}(X \times Y \times X \times Y)(d + e)
$$
est le cup-produit des images réciproques de $\gamma([\Delta_X])$
et de $\gamma([\Delta_Y])$ par $p$ et $q$. Écrivons
$$
\gamma([\Delta_{X \times Y}]) = \sum \eta_{X \times Y, i}
\text{ avec }
\eta_{X \times Y, i} \in
H^i(X \times Y) \otimes H^{2d + 2e - i}(X \times Y)(d + e)
$$,
et de même $\gamma([\Delta_X]) = \sum \eta_{X, i}$ et
$\gamma([\Delta_Y]) = \sum \eta_{Y, i}$. L'observation précédente
donne
$$
\eta_{X \times Y, 0} =
\sum\nolimits_{i \in \mathbf{Z}} p^*\eta_{X, i} \cup q^*\eta_{Y, -i}
$$
(Si notre théorie cohomologique s'annule en degrés négatifs, ce qui sera
le cas dans presque toutes les situations, seul le terme $i = 0$ contribue,
et $\eta_{X \times Y, 0}$ appartient à
$H^0(X) \otimes H^0(Y) \otimes H^{2d}(X)(d) \otimes H^{2e}(Y)(e)$, comme attendu ;
mais nous n'en avons pas besoin.) Puisque $\lambda_X : H^{2d}(X)(d) \to F$ et
$\lambda_Y : H^{2e}(Y)(e) \to F$ envoient $\eta_{X, 0}$ et $\eta_{Y, 0}$
sur $1$ dans $H^*(X)$ et $H^*(Y)$, on voit que $\lambda_X \otimes \lambda_Y$
envoie $\eta_{X \times Y, 0}$ sur $1$ dans
$H^*(X) \otimes H^*(Y) = H^*(X \times Y)$, ce qui achève la démonstration.
\end{proof}

\begin{lemma}
\label{lemma-trace-base}
Supposons que les données (D0), (D1) et (D2') satisfassent aux axiomes (A1) -- (A7).
Alors l'axiome (C)(d) de la section \ref{section-axioms} est vérifié.
\end{lemma}

\begin{proof}
On a $\gamma([\Spec(k)]) = 1 \in H^*(\Spec(k))$ par construction.
Puisque
$$
H^0(\Spec(k)) = F,\quad
H^0(\Spec(k) \times \Spec(k)) = H^0(\Spec(k)) \otimes_F H^0(\Spec(k))
$$,
l'application $\int_{\Spec(k)} = \lambda$ de l'axiome (A6) doit envoyer $1$ sur $1$,
car nous avons vu que
$\int_{\Spec(k) \times \Spec(k)} = \int_{\Spec(k)} \int_{\Spec(k)}$
au lemme \ref{lemma-trace-product}.
\end{proof}

\noindent
Supposons que les données (D0), (D1) et (D2') satisfassent aux axiomes (A1) -- (A7).
On obtient alors les données (D0), (D1), (D2) et (D3) de la
section \ref{section-axioms},
satisfaisant aux axiomes (A), (B), (C)(a), (C)(c) et (C)(d)
de la section \ref{section-axioms} ; voir
les lemmes \ref{lemma-poincare-duality}, \ref{lemma-trace-product} et
\ref{lemma-trace-base}.
De plus, on dispose des images directes $f_* : H^*(X) \to H^*(Y)$
construites à la section \ref{section-axioms}. Le seul axiome de la
section \ref{section-axioms}
qui ne soit pas encore établi est l'axiome (C)(b).

\begin{lemma}
\label{lemma-ok-for-projective-bundle}
Supposons que les données (D0), (D1) et (D2') satisfassent aux axiomes (A1) -- (A7).
Soit $p : P \to X$ comme dans l'axiome (A3), avec $X$ non vide et équidimensionnel.
Alors $\gamma$ commute à l'image directe par $p$.
\end{lemma}

\begin{proof}
Il suffit de le démontrer sur des générateurs de $\CH_*(P)$.
Il suffit donc de le démontrer pour une classe de cycle de la
forme $\xi^i \cdot p^*\alpha$, où $0 \leq i \leq r - 1$
et $\alpha \in \CH_*(X)$. Notons que $p_*(\xi^i \cdot p^*\alpha) = 0$
si $i < r - 1$, et que $p_*(\xi^{r - 1} \cdot p^*\alpha) = \alpha$.
D'autre part, on a
$\gamma(\xi^i \cdot p^*\alpha) = c^i \cup p^*\gamma(\alpha)$,
et la formule de projection (lemme \ref{lemma-pushforward})
donne
$$
p_*\gamma(\xi^i \cdot p^*\alpha) = p_*(c^i) \cup \gamma(\alpha)
$$
Il suffit ainsi de montrer que $p_*c^i = 0$ pour $i < r - 1$ et
que $p_*c^{r - 1} = 1$. De manière équivalente, il suffit de démontrer que
$\lambda_P : H^{2d + 2r - 2}(P)(d + r - 1) \to F$, définie par
les règles
$$
\lambda_P(c^i \cup p^*(a)) =
\left\{
\begin{matrix}
0 & \text{si} & i < r - 1 \\
\int_X(a) & \text{si} & i = r - 1
\end{matrix}
\right.
$$,
satisfait à la condition de l'axiome (A5). Cela résulte du
calcul de la classe de la diagonale de $P$ au
lemme \ref{lemma-diagonal-projective-bundle}.
\end{proof}

\begin{lemma}
\label{lemma-integrate-1}
Supposons que les données (D0), (D1) et (D2') satisfassent aux axiomes (A1) -- (A7).
Si $k'/k$ est une extension galoisienne, on a
$\int_{\Spec(k')} 1 = [k' : k]$.
\end{lemma}

\begin{proof}
Observons que
$$
\Spec(k') \times \Spec(k') =
\coprod\nolimits_{\sigma \in \text{Gal}(k'/k)}
(\Spec(\sigma) \times \text{id})^{-1} \Delta
$$
comme cycles sur $\Spec(k') \times \Spec(k')$.
Notre construction de $\gamma$ envoie toujours $[X]$ sur $1$ dans $H^0(X)$. Ainsi,
$1 \otimes 1 = 1 = \sum (\Spec(\sigma) \times \text{id})^*\gamma([\Delta])$.
Notons $\lambda : H^0(\Spec(k')) \to F$ l'application de
l'axiome (A6), c'est-à-dire que $(\text{id} \otimes \lambda)(\gamma(\Delta)) = 1$
dans $H^0(\Spec(k'))$. On obtient
\begin{align*}
\lambda(1) 1
& =
(\text{id} \otimes \lambda)(1 \otimes 1) \\
& =
(\text{id} \otimes \lambda)(
\sum (\Spec(\sigma) \times \text{id})^*\gamma([\Delta])) \\
& =
\sum (\Spec(\sigma) \times \text{id})^*(
(\text{id} \otimes \lambda)(\gamma([\Delta])) \\
& =
\sum (\Spec(\sigma) \times \text{id})^*(1) \\
& =
[k' : k]
\end{align*}
Puisque $\lambda$ est une autre notation pour $\int_{\Spec(k')}$
(remarque \ref{remark-trace}), la démonstration est achevée.
\end{proof}

\begin{lemma}
\label{lemma-enough}
Supposons que les données (D0), (D1) et (D2') satisfassent aux axiomes (A1) -- (A7).
Pour montrer que $\gamma$ commute à l'image directe, il suffit
de montrer que $i_*(1) = \gamma([Z])$ lorsque $i : Z \to X$ est une immersion
fermée de schémas projectifs lisses non vides et équidimensionnels sur $k$.
\end{lemma}

\begin{proof}
Nous utiliserons sans le mentionner de nouveau le fait que l'application
classe de cycle $\gamma$ construite au lemme \ref{lemma-cycle-classes}
est compatible avec les produits d'intersection et les images réciproques,
et que nous avons déjà établi les axiomes
(A), (B), (C)(a), (C)(c) et (C)(d) de la section \ref{section-axioms} ; voir
le lemme \ref{lemma-poincare-duality},
la remarque \ref{remark-trace} et
les lemmes \ref{lemma-trace-product} et \ref{lemma-trace-base}.
En particulier, on peut utiliser, par exemple, le lemme \ref{lemma-pushforward}
pour voir que l'image directe sur $H^*$ est compatible avec la composition
et satisfait à la formule de projection.

\medskip\noindent
Soit $f : X \to Y$ un morphisme de schémas projectifs lisses
non vides et équidimensionnels sur $k$.
Nous cherchons à montrer $f_*\gamma(\alpha) = \gamma(f_*\alpha)$
pour toute classe de cycle $\alpha$ sur $X$.
On peut écrire $\alpha$ comme une combinaison $\mathbf{Q}$-linéaire de produits de
classes de Chern de $\mathcal{O}_X$-modules localement libres
(Homologie de Chow, lemme \ref{chow-lemma-K-tensor-Q}).
On peut donc supposer que $\alpha$ est un produit de classes de Chern de
$\mathcal{O}_X$-modules localement libres de type fini
$\mathcal{E}_1, \ldots, \mathcal{E}_r$.
Choisissons $p : P \to X$ comme dans le principe de scindage
(Homologie de Chow, lemme \ref{chow-lemma-splitting-principle}).
D'après Homologie de Chow, remarque \ref{chow-remark-the-proof-shows-more},
$p$ est une composée de fibrés projectifs et
$\alpha = p_*(\xi_1 \cap \ldots \cap \xi_d \cap \cdot p^*\alpha)$,
où les $\xi_i$ sont des premières classes de Chern de modules inversibles.
D'après le lemme \ref{lemma-ok-for-projective-bundle},
$p_*$ commute aux classes de cycles.
Il suffit donc de démontrer la propriété pour la composée
$f \circ p$. Puisque $p^*\mathcal{E}_1, \ldots, p^*\mathcal{E}_r$
admettent des filtrations dont les quotients successifs sont des modules
inversibles, on est ramené au cas où $\alpha$ est
de la forme $\xi_1 \cap \ldots \cap \xi_t \cap [X]$,
pour certaines premières classes de Chern $\xi_i$ de modules inversibles $\mathcal{L}_i$.

\medskip\noindent
Supposons
$\alpha = c_1(\mathcal{L}_1) \cap \ldots \cap c_1(\mathcal{L}_t) \cap [X]$,
pour certains modules inversibles $\mathcal{L}_i$ sur $X$.
Soit $\mathcal{L}$ un $\mathcal{O}_X$-module inversible ample.
Pour $n \gg 0$, les $\mathcal{O}_X$-modules inversibles
$\mathcal{L}^{\otimes n}$ et
$\mathcal{L}_1 \otimes \mathcal{L}^{\otimes n}$ sont tous deux
très amples sur $X$ relativement à $k$ ; voir
Morphismes, lemme \ref{morphisms-lemma-invertible-add-enough-ample-very-ample}.
Puisque $c_1(\mathcal{L}_1) = c_1(\mathcal{L}_1 \otimes \mathcal{L}^{\otimes n})
- c_1(\mathcal{L}^{\otimes n})$, on est ramené au cas où
$\mathcal{L}_1$ est très ample. En répétant ce raisonnement avec $\mathcal{L}_i$
pour $i = 2, \ldots, t$, on se ramène au cas où $\mathcal{L}_i$
est très ample sur $X$ relativement à $k$ pour tout $i = 1, \ldots, t$.

\medskip\noindent
Supposons $k$ infini et $\alpha = 
c_1(\mathcal{L}_1) \cap \ldots \cap c_1(\mathcal{L}_t) \cap [X]$,
pour certains modules inversibles très amples $\mathcal{L}_i$ sur $X$ relativement à $k$.
Par Bertini, sous la forme de Variétés, lemme \ref{varieties-lemma-bertini},
on peut trouver successivement des sections régulières $s_i$ de $\mathcal{L}_i$
telles que les schémas $Z(s_1) \cap \ldots \cap Z(s_i)$
soient lisses sur $k$ et de codimension $i$ dans $X$.
Par la construction du cap-produit avec la première classe
d'un module inversible (qui remonte à
Homologie de Chow, définition \ref{chow-definition-divisor-invertible-sheaf}),
on est ainsi ramené au cas où $\alpha = [Z]$,
pour un sous-schéma fermé lisse non vide $Z \subset X$
équidimensionnel.

\medskip\noindent
Supposons $\alpha = [Z]$, où $Z \subset X$ est un sous-schéma fermé lisse.
Choisissons une immersion fermée $X \to \mathbf{P}^n$. On peut factoriser $f$ sous la forme
$$
X \to Y \times \mathbf{P}^n \to Y
$$
Puisque le résultat est connu pour le second morphisme par
le lemme \ref{lemma-ok-for-projective-bundle},
il suffit de le démontrer lorsque
$\alpha = [Z]$, où $i : Z \to X$ est une immersion fermée
et $f$ une immersion fermée.
Alors $j = f \circ i$ est également une immersion fermée.
L'hypothèse appliquée à $i$ et à $j$ donne la conclusion.

\medskip\noindent
Il reste à démontrer le lemme lorsque $k$ est fini. Nous conseillons vivement au
lecteur de passer la fin de la démonstration. Tout ce qui précède reste valable,
sauf que nous ne savons pas choisir les sections
$s_i$ qui découpent $Z$ sur $k$, puisque $k$ est fini. Nous savons toutefois
que l'on peut trouver $s_i$ sur la clôture algébrique $\overline{k}$
de $k$, par le même lemme. Il existe donc
une extension finie $k'/k$ telle que nos sections $s_i$
soient définies sur $k'$. Notons $\pi : X_{k'} \to X$ la projection.
Les arguments ci-dessus donnent la conclusion voulue,
sous l'hypothèse du lemme,
pour le cycle $\pi^*\alpha$ et le morphisme
$f \circ \pi : X_{k'} \to Y$.
On a $\pi_*\pi^*\alpha = [k' : k] \alpha$ ; voir
Homologie de Chow, lemme \ref{chow-lemma-finite-flat}.
D'autre part, on a
$$
\pi_*\gamma(\pi^*\alpha) = \pi_*\pi^*\gamma(\alpha) =
\gamma(\alpha) \pi_*1
$$
par la formule de projection de notre théorie cohomologique. Observons
que $\pi$ est une projection (!) ; on a donc
$\pi_*(1) = \int_{\Spec(k')}(1) 1$ d'après
le lemme \ref{lemma-pr2star}. Pour achever la démonstration dans
le cas d'un corps fini, il suffit donc de démontrer
$\int_{\Spec(k')}(1) = [k' : k]$, ce que nous faisons au
lemme \ref{lemma-integrate-1}.
\end{proof}

\noindent
Dans les lemmes ci-dessous, nous utilisons les grassmanniennes définies et construites
dans Constructions, section \ref{constructions-section-grassmannian}.

\begin{lemma}
\label{lemma-grassmanian}
Supposons que les données (D0), (D1) et (D2') vérifient les axiomes (A1) -- (A7).
Étant donnés des entiers $0 < l < n$ et un schéma projectif lisse
équidimensionnel non vide $X$ sur $k$, considérons la projection
$p : X \times \mathbf{G}(l, n) \to X$.
Alors $\gamma$ commute à l'image directe par $p$.
\end{lemma}

\begin{proof}
Si $l = 1$ ou $l = n - 1$, alors $p$ est un fibré projectif et
le résultat découle du lemme \ref{lemma-ok-for-projective-bundle}.
En général, il existe un morphisme
$$
h : Y \to X \times \mathbf{G}(l, n)
$$
tel que $h$ et $p \circ h$ soient tous deux des composés de fibrés
en espaces projectifs. En effet, notons $\mathbf{G}(1, 2, \ldots, l; n)$
la variété de drapeaux partiels. Alors le morphisme
$$
\mathbf{G}(1, 2, \ldots, l; n) \to \mathbf{G}(l, n)
$$
est un composé de fibrés en espaces projectifs, et le morphisme
structural $\mathbf{G}(1, 2, \ldots, l; n) \to \Spec(k)$
est lui aussi de cette forme. On peut donc poser $Y = X \times \mathbf{G}(1, 2, \ldots, l; n)$.
Puisque tout cycle sur $X \times \mathbf{G}(l, n)$ est l'image directe
d'un cycle sur $Y$, le résultat pour $Y \to X$ et celui pour
$Y \to X \times \mathbf{G}(l, n)$ entraînent le résultat pour $p$.
\end{proof}

\begin{lemma}
\label{lemma-enough-better}
Supposons que les données (D0), (D1) et (D2') vérifient les axiomes (A1) -- (A7).
Pour montrer que $\gamma$ commute à l'image directe, il suffit
de montrer que $i_*(1) = \gamma([Z])$ lorsque $i : Z \to X$ est une immersion
fermée de schémas projectifs lisses équidimensionnels non vides sur $k$
telle que la classe de $\mathcal{C}_{Z/X}$ dans $K_0(Z)$ soit
l'image réciproque d'une classe de $K_0(X)$.
\end{lemma}

\begin{proof}
D'après le lemme \ref{lemma-enough}, il suffit de montrer que $i_*(1) = \gamma([Z])$
lorsque $i : Z \to X$ est une immersion fermée de schémas
projectifs lisses équidimensionnels non vides
sur $k$. Supposons $Z$ de codimension $r$ dans $X$.
Soit $\mathcal{L}$ un module inversible suffisamment ample sur $X$.
Choisissons $n > 0$ et une surjection
$$
\mathcal{O}_Z^{\oplus n} \to \mathcal{C}_{Z/X} \otimes \mathcal{L}|_Z
$$
Ceci donne un morphisme $g : Z \to \mathbf{G}(n - r, n)$
vers la grassmannienne sur $k$, voir
Constructions, section \ref{constructions-section-grassmannian}.
Considérons la composée
$$
Z \to X \times \mathbf{G}(n - r, n) \to X
$$
L'image directe par le second morphisme est compatible aux classes
de cycles d'après le lemme \ref{lemma-grassmanian}. Le faisceau conormal $\mathcal{C}$
de l'immersion fermée $Z \to X \times \mathbf{G}(n - r, n)$ s'insère dans
une suite exacte courte
$$
0 \to \mathcal{C}_{Z/X} \to \mathcal{C} \to
g^*\Omega_{\mathbf{G}(n - r, n)} \to 0
$$
Voir Compléments sur les morphismes, lemme
\ref{more-morphisms-lemma-two-unramified-morphisms-formally-smooth}.
Puisque $\mathcal{C}_{Z/X} \otimes \mathcal{L}|_Z$ est l'image
réciproque d'un faisceau localement libre de rang fini sur $\mathbf{G}(n - r, n)$,
on en déduit que la classe de $\mathcal{C}$ dans $K_0(Z)$
est l'image réciproque d'une classe de $K_0(X \times \mathbf{G}(n - r, n))$.
On a donc la propriété pour $Z \to X \times \mathbf{G}(n - r, n)$,
ce qui permet de conclure.
\end{proof}

\begin{lemma}
\label{lemma-injective-H0}
Supposons que les données (D0), (D1) et (D2') vérifient les axiomes (A1) -- (A7).
Si $k''/k'/k$ sont des extensions finies séparables de corps, alors
$H^0(\Spec(k')) \to H^0(\Spec(k''))$ est injective.
\end{lemma}

\begin{proof}
On peut remplacer $k''$ par sa clôture normale sur $k$,
qui est galoisienne sur $k$, voir
Corps, lemme \ref{fields-lemma-normal-closure-galois}.
Alors $k''$ est aussi galoisien sur $k'$, voir
Corps, lemme \ref{fields-lemma-galois-goes-up}.
On en déduit un isomorphisme
$$
k' \otimes_k k'' \longrightarrow
\prod\nolimits_{\sigma \in \text{Gal}(k''/k')} k'',\quad
\eta \otimes \zeta \longmapsto (\eta \sigma(\zeta))_\sigma
$$
Il donne un isomorphisme
$\coprod_\sigma \Spec(k'') \to \Spec(k') \times \Spec(k'')$
qui induit en cohomologie l'isomorphisme
$$
H^*(\Spec(k')) \otimes_F H^*(\Spec(k''))
\to
\prod\nolimits_\sigma H^*(\Spec(k'')),\quad
a' \otimes a'' \longmapsto (\pi^*a' \cup \Spec(\sigma)^*a'')_\sigma
$$
où $\pi : \Spec(k'') \to \Spec(k')$ est le morphisme
correspondant à l'inclusion de $k'$ dans $k''$.
On conclut en prenant $a'' = 1$.
\end{proof}

\begin{lemma}
\label{lemma-pushforward-blowup}
Supposons que les données (D0), (D1) et (D2') vérifient les axiomes (A1) -- (A8).
Soit $b : X' \to X$ l'éclatement d'un schéma projectif lisse $X$
sur $k$, non vide, équidimensionnel de dimension $d$,
le long d'un centre lisse partout non dense $Z$. Alors $b_*(1) = 1$.
\end{lemma}

\begin{proof}
On peut remplacer $X$ par une composante connexe de $X$ (certains détails
sont omis). On peut donc supposer $X$ connexe, donc irréductible.
Posons $k' = \Gamma(X, \mathcal{O}_X) = \Gamma(X', \mathcal{O}_{X'})$ ;
nous omettons la démonstration de cette égalité. Choisissons un point fermé $x' \in X'$
qui n'appartient pas au diviseur exceptionnel et dont le corps résiduel
$k''$ est séparable sur $k$ ; ceci est possible d'après
Variétés, lemme \ref{varieties-lemma-smooth-separable-closed-points-dense}.
Notons $x \in X$ son image (dont le corps résiduel est bien entendu
égal lui aussi à $k''$). Considérons le diagramme
$$
\xymatrix{
x' \times X' \ar[r] \ar[d] & X' \times X' \ar[d] \\
x \times X \ar[r] & X \times X
}
$$
L'image réciproque de la classe de la diagonale $\Delta = \Delta_X$ est la classe du
« point diagonal » $\delta_x : x \to x \times X$, et de même pour la classe de
la diagonale $\Delta'$. D'autre part, l'image réciproque du point diagonal $\delta_x$
par la flèche verticale de gauche est le point diagonal $\delta_{x'}$.
Écrivons $\gamma([\Delta]) = \sum \eta_i$ avec
$\eta_i \in H^i(X) \otimes H^{2d - i}(X)(d)$ et
$\gamma([\Delta']) = \sum \eta'_i$ avec
$\eta'_i \in H^i(X') \otimes H^{2d - i}(X')(d)$.
Les arguments précédents montrent que $\eta_0$ et $\eta'_0$ ont pour image la même
classe dans
$$
H^0(x') \otimes_F H^{2d}(X')(d)
$$
On a $H^0(\Spec(k')) = H^0(X) = H^0(X')$ d'après l'axiome (A8).
D'après le lemme \ref{lemma-injective-H0}, cet espace commun s'injecte
dans $H^0(x')$. On en déduit que $\eta_0$ a pour image $\eta'_0$ par l'application
$$
H^0(X) \otimes_F H^{2d}(X)(d)
\longrightarrow
H^0(X') \otimes_F H^{2d}(X')(d)
$$
Cela signifie que $\int_X$ est égal à $\int_{X'}$ composé avec
l'application d'image réciproque. Le lemme est démontré.
\end{proof}

\begin{lemma}
\label{lemma-done}
Supposons que les données (D0), (D1) et (D2') vérifient les axiomes (A1) -- (A8).
Alors l'application classe de cycle $\gamma$ commute à l'image directe.
\end{lemma}

\begin{proof}
Soit $i : Z \to X$ comme dans le lemme \ref{lemma-enough-better}. Considérons
le diagramme
$$
\xymatrix{
E \ar[r]_j \ar[d]_\pi & X' \ar[d]^b \\
Z \ar[r]^i & X
}
$$
Soit $\theta \in \CH^{r - 1}(X')$ comme dans le
lemme \ref{lemma-divide-pullback-good-blowing-up}.
Alors $\pi_*j^!\theta = [Z]$ dans $\CH_*(Z)$ entraîne que
$\pi_*\gamma(j^!\theta) = 1$ d'après le lemme \ref{lemma-ok-for-projective-bundle},
car $\pi$ est un fibré en espaces projectifs.
On obtient donc
$$
i_*(1) = i_*(\pi_*(\gamma(j^!\theta))) =
b_*j_*(j^*\gamma(\theta)) =
b_*(j_*(1) \cup \gamma(\theta))
$$
On a $j_*(1) = \gamma([E])$ d'après (A9). Ceci est donc égal à
$$
b_*(\gamma([E]) \cup \gamma(\theta)) =
b_*(\gamma([E] \cdot \theta)) =
b_*(\gamma(b^*[Z])) =
b_*b^*\gamma([Z]) = b_*(1) \cup \gamma([Z])
$$
Comme $b_*(1) = 1$ d'après le lemme \ref{lemma-pushforward-blowup}, la
démonstration est achevée.
\end{proof}

\begin{proposition}
\label{proposition-get-weil}
Supposons que les données (D0), (D1) et (D2') vérifient les axiomes (A1) -- (A8).
On obtient alors une théorie cohomologique de Weil.
\end{proposition}

\begin{proof}
Les axiomes (A), (B), (C)(a), (C)(c) et (C)(d) de la
section \ref{section-axioms} résultent des
lemmes \ref{lemma-poincare-duality}, \ref{lemma-trace-product} et
\ref{lemma-trace-base}.
L'axiome (C)(b) résulte du
lemme \ref{lemma-done}.
Enfin, la condition supplémentaire de la
définition \ref{definition-weil-cohomology-theory}
est satisfaite, car elle coïncide avec notre axiome (A8).
\end{proof}

\noindent
Le lemme suivant est parfois utile pour montrer que l'on obtient une
théorie cohomologique de Weil sur un corps non algébriquement clos, en se
ramenant à un corps algébriquement clos.

\begin{lemma}
\label{lemma-check-over-extension}
Soit $k'/k$ une extension de corps. Soit $F'/F$ une extension
de corps de caractéristique $0$. Supposons donnés
\begin{enumerate}
\item des données (D0), (D1), (D2') pour $k$ et $F$, notées
$F(1), H^*, c_1^H$,
\item des données (D0), (D1), (D2') pour $k'$ et $F'$, notées
$F'(1), (H')^*, c_1^{H'}$, et
\item un isomorphisme $F(1) \otimes_F F' \to F'(1)$, des isomorphismes fonctoriels
$H^*(X) \otimes_F F' \to (H')^*(X_{k'})$ sur la catégorie des schémas projectifs lisses
$X$ sur $k$, tels que les diagrammes
$$
\xymatrix{
\Pic(X) \ar[r]_{c_1^H} \ar[d] & H^2(X)(1) \ar[d] \\
\Pic(X_{k'}) \ar[r]^{c_1^{H'}} & (H')^2(X_{k'})(1)
}
$$
commutent.
\end{enumerate}
Dans ce cas, si $F'(1), (H')^*, c_1^{H'}$ vérifient les axiomes (A1) -- (A9),
il en est de même de $F(1), H^*, c_1^H$.
\end{lemma}

\begin{proof}
Nous vérifions les axiomes un à un.

\medskip\noindent
Axiome (A1). Il faut montrer que $H^*(\emptyset) = 0$ et que
$(i^*, j^*) : H^*(X \amalg Y) \to H^*(X) \times H^*(Y)$
est un isomorphisme, où $i$ et $j$ sont les coprojections.
Par fonctorialité des isomorphismes
$H^*(X) \otimes_F F' \to (H')^*(X_{k'})$, cela
résulte de l'algèbre linéaire : si $\varphi : V \to W$ est une application $F$-linéaire
d'espaces vectoriels sur $F$, alors $\varphi$ est un isomorphisme si et seulement si
$\varphi_{F'} : V \otimes_F F' \to W \otimes_F F'$ est un isomorphisme.

\medskip\noindent
Axiome (A2). Cela signifie qu'étant donnés un morphisme $f : X \to Y$ de schémas projectifs lisses
sur $k$ et un $\mathcal{O}_Y$-module inversible $\mathcal{N}$,
on a $f^*c_1^H(\mathcal{L}) = c_1^H(f^*\mathcal{L})$. Ceci résulte immédiatement
de la propriété correspondante pour $c_1^{H'}$, des diagrammes commutatifs
du lemme et du fait que l'application canonique
$V \to V \otimes_F F'$ est injective pour tout espace vectoriel sur $F$, noté $V$.

\medskip\noindent
Axiome (A3). Il découle du principe énoncé dans la démonstration de
l'axiome (A1) et de la compatibilité de $c_1^H$ et $c_1^{H'}$.

\medskip\noindent
Axiome (A4). Soit $i : Y \to X$ l'inclusion d'un diviseur de Cartier
effectif sur $k$, où $X$ et $Y$ sont tous deux projectifs et lisses
sur $k$. Pour $a \in H^*(X)$ tel que
$i^*a = 0$, il faut montrer que $a \cup c_1^H(\mathcal{O}_X(Y)) = 0$.
Notons $a' \in (H')^*(X_{k'})$ l'image de $a$.
L'hypothèse entraîne que $(i')^*a' = 0$, où $i' : Y_{k'} \to X_{k'}$
est le changement de base de $i$. On obtient donc
$a' \cup c_1^{H'}(\mathcal{O}_{X_{k'}}(Y_{k'})) = 0$ par l'axiome
pour $(H')^*$. Puisque $a' \cup c_1^{H'}(\mathcal{O}_{X_{k'}}(Y_{k'}))$
est l'image de $a \cup c_1^H(\mathcal{O}_X(Y))$, on conclut par
le principe énoncé dans la démonstration de l'axiome (A2).

\medskip\noindent
Axiome (A5). Cela signifie que $H^*(\Spec(k)) = F$ et que, pour $X$ et $Y$ projectifs
lisses sur $k$, l'application $H^*(X) \otimes_F H^*(Y) \to H^*(X \times Y)$,
$a \otimes b \mapsto p^*(a) \cup q^*(b)$ est un isomorphisme,
où $p$ et $q$ sont les projections. Ceci découle du principe
énoncé dans la démonstration de l'axiome (A1).

\medskip\noindent
Interrompons le cours de l'argumentation pour montrer que, pour tout
schéma projectif lisse $X$ sur $k$, le diagramme
$$
\xymatrix{
\CH^*(X) \ar[r]_-\gamma \ar[d]_{g^*} & \bigoplus H^{2i}(X)(i) \ar[d] \\
\CH^*(X_{k'}) \ar[r]^-{\gamma'} & \bigoplus (H')^{2i}(X_{k'})(i)
}
$$
commute. Remarquons que l'on dispose de $\gamma$, puisque l'on a déjà établi les axiomes
(A1) -- (A4) ; voir le lemme \ref{lemma-cycle-classes}.
La flèche verticale de gauche est d'ailleurs celle qui est étudiée dans
Homologie de Chow, section \ref{chow-section-change-base},
pour le morphisme de schémas de base $g : \Spec(k') \to \Spec(k)$.
Plus précisément, c'est l'application donnée dans
Homologie de Chow, lemme \ref{chow-lemma-pullback-base-change}.
Prenons $\alpha \in \CH^*(X)$. Écrivons $\alpha = ch(\beta) \cap [X]$
dans $\CH^*(X) \otimes \mathbf{Q}$,
pour un certain $\beta \in K_0(\textit{Vect}(X)) \otimes \mathbf{Q}$,
de sorte que $\gamma(\alpha) = ch^{H}(\beta)$ ; telle est notre construction de $\gamma$.
Puisque l'application de Homologie de Chow, lemme \ref{chow-lemma-pullback-base-change},
est compatible au cap-produit par les classes de Chern d'après
Homologie de Chow, lemme \ref{chow-lemma-pullback-base-change-chern-classes},
on voit que $g^*\alpha = ch((X_{k'} \to X)^*\beta) \cap [X_{k'}]$.
Ainsi $\gamma'(g^*\alpha) = ch^{H'}((X_{k'} \to X)^*\beta)$.
Le diagramme sera donc commutatif si, pour tout
$\mathcal{O}_X$-module localement libre $\mathcal{E}$ de rang $r$ et tout $0 \leq i \leq r$,
l'élément $c_i^H(\mathcal{E})$ de $H^{2i}(X)(i)$
a pour image l'élément $c_i^{H'}(\mathcal{E}_{k'})$ de $(H')^{2i}(X_{k'})(i)$.
Comme on dispose de la formule du fibré en espaces projectifs aussi bien pour
$X$ que pour $X'$, on peut, pour le montrer, remplacer $X$ par un fibré en espaces projectifs
sur $X$ un nombre fini de fois. On peut donc supposer
que $\mathcal{E}$ possède une filtration dont les quotients gradués sont des
$\mathcal{O}_X$-modules inversibles
$\mathcal{L}_1, \ldots, \mathcal{L}_r$.
Voir Homologie de Chow, lemme \ref{chow-lemma-splitting-principle} et
remarque \ref{chow-remark-the-proof-shows-more}.
Alors $c^H_i(\mathcal{E}$ est le $i$-ième polynôme symétrique élémentaire
en $c^H_1(\mathcal{L}_1), \ldots, c^H_1(\mathcal{L}_r)$,
et l'on conclut par l'hypothèse de compatibilité des
premières classes de Chern.

\medskip\noindent
Axiome (A6). Supposons donnés des espaces vectoriels sur $F$,
$V$, $W$, un élément $v \in V$ et un tenseur $\xi \in V \otimes_F W$.
Notons $V' = V \otimes_F F'$, $W' = W \otimes_F F'$, et $v'$, $\xi'$
les images de $v$, $\xi$ dans $V'$, $V' \otimes_{F'} W'$. Le principe d'algèbre linéaire
utilisé pour démontrer l'axiome (A6) est le suivant :
il existe une application $F$-linéaire $\lambda : W \to F$ telle que
$(1 \otimes \lambda)\xi = v$ si et seulement s'il existe une application $F'$-linéaire
$\lambda' : W \otimes_F F' \to F'$ telle que
$(1 \otimes \lambda')\xi' = v'$.

\medskip\noindent
Soit $X$ un schéma projectif lisse équidimensionnel non vide
sur $k$, de dimension $d$. Notons $\gamma = \gamma([\Delta])$
dans $H^{2d}(X \times X)(d)$ (les produits fibrés sans indication de base seront pris sur $k$).
Remarquons, ou rappelons, que ceci a un sens, puisque l'on a déjà établi les axiomes (A1) -- (A4) ;
voir le lemme \ref{lemma-cycle-classes}. On peut écrire
$$
\gamma = \sum \gamma_i, \quad
\gamma_i \in H^i(X) \otimes_F H^{2d - i}(X)(d)
$$
dans la décomposition de K\"unneth. De même, notons
$\gamma' = \gamma([\Delta']) = \sum \gamma'_i$
dans $(H')^{2d}(X_{k'} \times_{k'} X_{k'})(d)$.
D'après le principe d'algèbre linéaire rappelé ci-dessus, il suffit
de montrer que $\gamma_0$ a pour image $\gamma'_0$ dans
$(H')^0(X) \otimes_{F'} (H')^{2d}(X')(d)$.
La compatibilité des décompositions de K\"unneth
montre qu'il suffit de vérifier que $\gamma$ a pour image
$\gamma'$ dans
$$
(H')^{2d}(X_{k'} \times_{k'} X_{k'})(d) = (H')^{2d}((X \times X)_{k'})(d)
$$
Puisque $\Delta_{k'} = \Delta'$, ceci résulte de la discussion précédente.

\medskip\noindent
Axiome (A7). Il résulte du fait d'algèbre linéaire suivant : une
application linéaire $V \to W$ d'espaces vectoriels sur $F$ est injective
si et seulement si $V \otimes_F F' \to W \otimes_F F'$ est injective.

\medskip\noindent
Axiome (A8). Il résulte du fait d'algèbre linéaire utilisé dans
la démonstration de l'axiome (A1).

\medskip\noindent
Axiome (A9). Soit $X$ un schéma projectif lisse non vide sur $k$,
équidimensionnel de dimension $d$. Soit $i : Y \to X$ un diviseur de Cartier
effectif non vide, lisse sur $k$.
Soient $\lambda_Y$ et $\lambda_X$ comme dans l'axiome (A6) pour $X$ et $Y$.
Il faut montrer que, pour $a \in H^{2d - 2}(X)(d - 1)$,
on a $\lambda_Y(i^*(a)) = \lambda_X(a \cup c_1^H(\mathcal{O}_X(Y))$.
D'après la remarque \ref{remark-trace},
on sait que $\lambda_X : H^{2d}(X)(d) \to F$ et
$\lambda_Y : H^{2d - 2}(Y)(d - 1)$ sont uniquement
déterminés par la condition de l'axiome (A6).
Cela étant, il résulte de la démonstration de l'axiome (A6) pour $H^*$ ci-dessus
que $\lambda_X \otimes \text{id}_{F'}$ correspond à $\lambda_{X_{k'}}$
par l'identification donnée $H^{2d}(X)(d) \otimes_F F' = H^{2d}(X_{k'})(d)$.
Ainsi, le fait que l'axiome (A9) soit vérifié pour $F'(1), (H')^*, c_1^{H'}$
entraîne qu'il l'est pour $F(1), H^*, c_1^H$, par une chasse au diagramme.
Ceci achève la démonstration du lemme.
\end{proof}















\begin{multicols}{2}[\section{Autres chapitres}]
\noindent
Préliminaires
\begin{enumerate}
\item \hyperref[introduction-section-phantom]{Introduction}
\item \hyperref[conventions-section-phantom]{Conventions}
\item \hyperref[sets-section-phantom]{Théorie des ensembles}
\item \hyperref[categories-section-phantom]{Catégories}
\item \hyperref[topology-section-phantom]{Topologie}
\item \hyperref[sheaves-section-phantom]{Faisceaux sur les espaces}
\item \hyperref[sites-section-phantom]{Sites et faisceaux}
\item \hyperref[stacks-section-phantom]{Champs}
\item \hyperref[fields-section-phantom]{Corps}
\item \hyperref[algebra-section-phantom]{Algèbre commutative}
\item \hyperref[brauer-section-phantom]{Groupes de Brauer}
\item \hyperref[homology-section-phantom]{Algèbre homologique}
\item \hyperref[derived-section-phantom]{Catégories dérivées}
\item \hyperref[simplicial-section-phantom]{Méthodes simpliciales}
\item \hyperref[more-algebra-section-phantom]{Compléments d'algèbre}
\item \hyperref[smoothing-section-phantom]{Lissage des morphismes d'anneaux}
\item \hyperref[modules-section-phantom]{Faisceaux de modules}
\item \hyperref[sites-modules-section-phantom]{Modules sur les sites}
\item \hyperref[injectives-section-phantom]{Objets injectifs}
\item \hyperref[cohomology-section-phantom]{Cohomologie des faisceaux}
\item \hyperref[sites-cohomology-section-phantom]{Cohomologie sur les sites}
\item \hyperref[dga-section-phantom]{Algèbre différentielle graduée}
\item \hyperref[dpa-section-phantom]{Algèbre à puissances divisées}
\item \hyperref[sdga-section-phantom]{Faisceaux différentiels gradués}
\item \hyperref[hypercovering-section-phantom]{Hyperrecouvrements}
\end{enumerate}
Schémas
\begin{enumerate}
\setcounter{enumi}{25}
\item \hyperref[schemes-section-phantom]{Schémas}
\item \hyperref[constructions-section-phantom]{Constructions de schémas}
\item \hyperref[properties-section-phantom]{Propriétés des schémas}
\item \hyperref[morphisms-section-phantom]{Morphismes de schémas}
\item \hyperref[coherent-section-phantom]{Cohomologie des schémas}
\item \hyperref[divisors-section-phantom]{Diviseurs}
\item \hyperref[limits-section-phantom]{Limites de schémas}
\item \hyperref[varieties-section-phantom]{Variétés}
\item \hyperref[topologies-section-phantom]{Topologies sur les schémas}
\item \hyperref[descent-section-phantom]{Descente}
\item \hyperref[perfect-section-phantom]{Catégories dérivées de schémas}
\item \hyperref[more-morphisms-section-phantom]{Compléments sur les morphismes}
\item \hyperref[flat-section-phantom]{Compléments sur la platitude}
\item \hyperref[groupoids-section-phantom]{Schémas en groupoïdes}
\item \hyperref[more-groupoids-section-phantom]{Compléments sur les schémas en groupoïdes}
\item \hyperref[etale-section-phantom]{Morphismes étales de schémas}
\end{enumerate}
Thèmes de théorie des schémas
\begin{enumerate}
\setcounter{enumi}{41}
\item \hyperref[chow-section-phantom]{Homologie de Chow}
\item \hyperref[intersection-section-phantom]{Théorie de l'intersection}
\item \hyperref[pic-section-phantom]{Schémas de Picard des courbes}
\item \hyperref[weil-section-phantom]{Théories cohomologiques de Weil}
\item \hyperref[adequate-section-phantom]{Modules adéquats}
\item \hyperref[dualizing-section-phantom]{Complexes dualisants}
\item \hyperref[duality-section-phantom]{Dualité pour les schémas}
\item \hyperref[discriminant-section-phantom]{Discriminants et différentes}
\item \hyperref[derham-section-phantom]{Cohomologie de de Rham}
\item \hyperref[local-cohomology-section-phantom]{Cohomologie locale}
\item \hyperref[algebraization-section-phantom]{Géométrie algébrique et formelle}
\item \hyperref[curves-section-phantom]{Courbes algébriques}
\item \hyperref[resolve-section-phantom]{Résolution des surfaces}
\item \hyperref[models-section-phantom]{Réduction semi-stable}
\item \hyperref[functors-section-phantom]{Foncteurs et morphismes}
\item \hyperref[equiv-section-phantom]{Catégories dérivées de variétés}
\item \hyperref[pione-section-phantom]{Groupes fondamentaux de schémas}
\item \hyperref[etale-cohomology-section-phantom]{Cohomologie étale}
\item \hyperref[crystalline-section-phantom]{Cohomologie cristalline}
\item \hyperref[proetale-section-phantom]{Cohomologie pro-étale}
\item \hyperref[relative-cycles-section-phantom]{Cycles relatifs}
\item \hyperref[more-etale-section-phantom]{Compléments de cohomologie étale}
\item \hyperref[trace-section-phantom]{La formule des traces}
\end{enumerate}
Espaces algébriques
\begin{enumerate}
\setcounter{enumi}{64}
\item \hyperref[spaces-section-phantom]{Espaces algébriques}
\item \hyperref[spaces-properties-section-phantom]{Propriétés des espaces algébriques}
\item \hyperref[spaces-morphisms-section-phantom]{Morphismes d'espaces algébriques}
\item \hyperref[decent-spaces-section-phantom]{Espaces algébriques décents}
\item \hyperref[spaces-cohomology-section-phantom]{Cohomologie des espaces algébriques}
\item \hyperref[spaces-limits-section-phantom]{Limites d'espaces algébriques}
\item \hyperref[spaces-divisors-section-phantom]{Diviseurs sur les espaces algébriques}
\item \hyperref[spaces-over-fields-section-phantom]{Espaces algébriques sur des corps}
\item \hyperref[spaces-topologies-section-phantom]{Topologies sur les espaces algébriques}
\item \hyperref[spaces-descent-section-phantom]{Descente et espaces algébriques}
\item \hyperref[spaces-perfect-section-phantom]{Catégories dérivées d'espaces}
\item \hyperref[spaces-more-morphisms-section-phantom]{Compléments sur les morphismes d'espaces}
\item \hyperref[spaces-flat-section-phantom]{Platitude sur les espaces algébriques}
\item \hyperref[spaces-groupoids-section-phantom]{Groupoïdes dans les espaces algébriques}
\item \hyperref[spaces-more-groupoids-section-phantom]{Compléments sur les groupoïdes dans les espaces}
\item \hyperref[bootstrap-section-phantom]{Amorçage}
\item \hyperref[spaces-pushouts-section-phantom]{Sommes amalgamées d'espaces algébriques}
\end{enumerate}
Thèmes de géométrie
\begin{enumerate}
\setcounter{enumi}{81}
\item \hyperref[spaces-chow-section-phantom]{Groupes de Chow des espaces}
\item \hyperref[groupoids-quotients-section-phantom]{Quotients de groupoïdes}
\item \hyperref[spaces-more-cohomology-section-phantom]{Compléments sur la cohomologie des espaces}
\item \hyperref[spaces-simplicial-section-phantom]{Espaces simpliciaux}
\item \hyperref[spaces-duality-section-phantom]{Dualité pour les espaces}
\item \hyperref[formal-spaces-section-phantom]{Espaces algébriques formels}
\item \hyperref[restricted-section-phantom]{Algébrisation des espaces formels}
\item \hyperref[spaces-resolve-section-phantom]{Résolution des surfaces revisitée}
\end{enumerate}
Théorie des déformations
\begin{enumerate}
\setcounter{enumi}{89}
\item \hyperref[formal-defos-section-phantom]{Théorie formelle des déformations}
\item \hyperref[defos-section-phantom]{Théorie des déformations}
\item \hyperref[cotangent-section-phantom]{Le complexe cotangent}
\item \hyperref[examples-defos-section-phantom]{Problèmes de déformation}
\end{enumerate}
Champs algébriques
\begin{enumerate}
\setcounter{enumi}{93}
\item \hyperref[algebraic-section-phantom]{Champs algébriques}
\item \hyperref[examples-stacks-section-phantom]{Exemples de champs}
\item \hyperref[stacks-sheaves-section-phantom]{Faisceaux sur les champs algébriques}
\item \hyperref[criteria-section-phantom]{Critères de représentabilité}
\item \hyperref[artin-section-phantom]{Axiomes d'Artin}
\item \hyperref[quot-section-phantom]{Espaces de Quot et de Hilbert}
\item \hyperref[stacks-properties-section-phantom]{Propriétés des champs algébriques}
\item \hyperref[stacks-morphisms-section-phantom]{Morphismes de champs algébriques}
\item \hyperref[stacks-limits-section-phantom]{Limites de champs algébriques}
\item \hyperref[stacks-cohomology-section-phantom]{Cohomologie des champs algébriques}
\item \hyperref[stacks-perfect-section-phantom]{Catégories dérivées de champs}
\item \hyperref[stacks-introduction-section-phantom]{Introduction aux champs algébriques}
\item \hyperref[stacks-more-morphisms-section-phantom]{Compléments sur les morphismes de champs}
\item \hyperref[stacks-geometry-section-phantom]{La géométrie des champs}
\end{enumerate}
Thèmes de théorie des espaces de modules
\begin{enumerate}
\setcounter{enumi}{107}
\item \hyperref[moduli-section-phantom]{Champs de modules}
\item \hyperref[moduli-curves-section-phantom]{Espaces de modules de courbes}
\end{enumerate}
Divers
\begin{enumerate}
\setcounter{enumi}{109}
\item \hyperref[examples-section-phantom]{Exemples}
\item \hyperref[exercises-section-phantom]{Exercices}
\item \hyperref[guide-section-phantom]{Guide bibliographique}
\item \hyperref[desirables-section-phantom]{Desiderata}
\item \hyperref[coding-section-phantom]{Style de programmation}
\item \hyperref[obsolete-section-phantom]{Obsolète}
\item \hyperref[fdl-section-phantom]{Licence de documentation libre GNU}
\item \hyperref[index-section-phantom]{Index généré automatiquement}
\end{enumerate}
\end{multicols}


\bibliography{my}
\bibliographystyle{amsalpha}

\end{document}
